% Options for packages loaded elsewhere
% Options for packages loaded elsewhere
\PassOptionsToPackage{unicode}{hyperref}
\PassOptionsToPackage{hyphens}{url}
\PassOptionsToPackage{dvipsnames,svgnames,x11names}{xcolor}
%
\documentclass[
  portuguese,
  letterpaper,
  openany,
  oneside,
  12pt,
  a4paper]{book}
\usepackage{xcolor}
\usepackage[left=2cm,right=2cm,top=2.5cm,bottom=2.5cm]{geometry}
\usepackage{amsmath,amssymb}
\setcounter{secnumdepth}{5}
\usepackage{iftex}
\ifPDFTeX
  \usepackage[T1]{fontenc}
  \usepackage[utf8]{inputenc}
  \usepackage{textcomp} % provide euro and other symbols
\else % if luatex or xetex
  \usepackage{unicode-math} % this also loads fontspec
  \defaultfontfeatures{Scale=MatchLowercase}
  \defaultfontfeatures[\rmfamily]{Ligatures=TeX,Scale=1}
\fi
\usepackage{lmodern}
\ifPDFTeX\else
  % xetex/luatex font selection
\fi
% Use upquote if available, for straight quotes in verbatim environments
\IfFileExists{upquote.sty}{\usepackage{upquote}}{}
\IfFileExists{microtype.sty}{% use microtype if available
  \usepackage[]{microtype}
  \UseMicrotypeSet[protrusion]{basicmath} % disable protrusion for tt fonts
}{}
\makeatletter
\@ifundefined{KOMAClassName}{% if non-KOMA class
  \IfFileExists{parskip.sty}{%
    \usepackage{parskip}
  }{% else
    \setlength{\parindent}{0pt}
    \setlength{\parskip}{6pt plus 2pt minus 1pt}}
}{% if KOMA class
  \KOMAoptions{parskip=half}}
\makeatother
% Make \paragraph and \subparagraph free-standing
\makeatletter
\ifx\paragraph\undefined\else
  \let\oldparagraph\paragraph
  \renewcommand{\paragraph}{
    \@ifstar
      \xxxParagraphStar
      \xxxParagraphNoStar
  }
  \newcommand{\xxxParagraphStar}[1]{\oldparagraph*{#1}\mbox{}}
  \newcommand{\xxxParagraphNoStar}[1]{\oldparagraph{#1}\mbox{}}
\fi
\ifx\subparagraph\undefined\else
  \let\oldsubparagraph\subparagraph
  \renewcommand{\subparagraph}{
    \@ifstar
      \xxxSubParagraphStar
      \xxxSubParagraphNoStar
  }
  \newcommand{\xxxSubParagraphStar}[1]{\oldsubparagraph*{#1}\mbox{}}
  \newcommand{\xxxSubParagraphNoStar}[1]{\oldsubparagraph{#1}\mbox{}}
\fi
\makeatother

\usepackage{color}
\usepackage{fancyvrb}
\newcommand{\VerbBar}{|}
\newcommand{\VERB}{\Verb[commandchars=\\\{\}]}
\DefineVerbatimEnvironment{Highlighting}{Verbatim}{commandchars=\\\{\}}
% Add ',fontsize=\small' for more characters per line
\usepackage{framed}
\definecolor{shadecolor}{RGB}{241,243,245}
\newenvironment{Shaded}{\begin{snugshade}}{\end{snugshade}}
\newcommand{\AlertTok}[1]{\textcolor[rgb]{0.68,0.00,0.00}{#1}}
\newcommand{\AnnotationTok}[1]{\textcolor[rgb]{0.37,0.37,0.37}{#1}}
\newcommand{\AttributeTok}[1]{\textcolor[rgb]{0.40,0.45,0.13}{#1}}
\newcommand{\BaseNTok}[1]{\textcolor[rgb]{0.68,0.00,0.00}{#1}}
\newcommand{\BuiltInTok}[1]{\textcolor[rgb]{0.00,0.23,0.31}{#1}}
\newcommand{\CharTok}[1]{\textcolor[rgb]{0.13,0.47,0.30}{#1}}
\newcommand{\CommentTok}[1]{\textcolor[rgb]{0.37,0.37,0.37}{#1}}
\newcommand{\CommentVarTok}[1]{\textcolor[rgb]{0.37,0.37,0.37}{\textit{#1}}}
\newcommand{\ConstantTok}[1]{\textcolor[rgb]{0.56,0.35,0.01}{#1}}
\newcommand{\ControlFlowTok}[1]{\textcolor[rgb]{0.00,0.23,0.31}{\textbf{#1}}}
\newcommand{\DataTypeTok}[1]{\textcolor[rgb]{0.68,0.00,0.00}{#1}}
\newcommand{\DecValTok}[1]{\textcolor[rgb]{0.68,0.00,0.00}{#1}}
\newcommand{\DocumentationTok}[1]{\textcolor[rgb]{0.37,0.37,0.37}{\textit{#1}}}
\newcommand{\ErrorTok}[1]{\textcolor[rgb]{0.68,0.00,0.00}{#1}}
\newcommand{\ExtensionTok}[1]{\textcolor[rgb]{0.00,0.23,0.31}{#1}}
\newcommand{\FloatTok}[1]{\textcolor[rgb]{0.68,0.00,0.00}{#1}}
\newcommand{\FunctionTok}[1]{\textcolor[rgb]{0.28,0.35,0.67}{#1}}
\newcommand{\ImportTok}[1]{\textcolor[rgb]{0.00,0.46,0.62}{#1}}
\newcommand{\InformationTok}[1]{\textcolor[rgb]{0.37,0.37,0.37}{#1}}
\newcommand{\KeywordTok}[1]{\textcolor[rgb]{0.00,0.23,0.31}{\textbf{#1}}}
\newcommand{\NormalTok}[1]{\textcolor[rgb]{0.00,0.23,0.31}{#1}}
\newcommand{\OperatorTok}[1]{\textcolor[rgb]{0.37,0.37,0.37}{#1}}
\newcommand{\OtherTok}[1]{\textcolor[rgb]{0.00,0.23,0.31}{#1}}
\newcommand{\PreprocessorTok}[1]{\textcolor[rgb]{0.68,0.00,0.00}{#1}}
\newcommand{\RegionMarkerTok}[1]{\textcolor[rgb]{0.00,0.23,0.31}{#1}}
\newcommand{\SpecialCharTok}[1]{\textcolor[rgb]{0.37,0.37,0.37}{#1}}
\newcommand{\SpecialStringTok}[1]{\textcolor[rgb]{0.13,0.47,0.30}{#1}}
\newcommand{\StringTok}[1]{\textcolor[rgb]{0.13,0.47,0.30}{#1}}
\newcommand{\VariableTok}[1]{\textcolor[rgb]{0.07,0.07,0.07}{#1}}
\newcommand{\VerbatimStringTok}[1]{\textcolor[rgb]{0.13,0.47,0.30}{#1}}
\newcommand{\WarningTok}[1]{\textcolor[rgb]{0.37,0.37,0.37}{\textit{#1}}}

\usepackage{longtable,booktabs,array}
\usepackage{calc} % for calculating minipage widths
% Correct order of tables after \paragraph or \subparagraph
\usepackage{etoolbox}
\makeatletter
\patchcmd\longtable{\par}{\if@noskipsec\mbox{}\fi\par}{}{}
\makeatother
% Allow footnotes in longtable head/foot
\IfFileExists{footnotehyper.sty}{\usepackage{footnotehyper}}{\usepackage{footnote}}
\makesavenoteenv{longtable}
\usepackage{graphicx}
\makeatletter
\newsavebox\pandoc@box
\newcommand*\pandocbounded[1]{% scales image to fit in text height/width
  \sbox\pandoc@box{#1}%
  \Gscale@div\@tempa{\textheight}{\dimexpr\ht\pandoc@box+\dp\pandoc@box\relax}%
  \Gscale@div\@tempb{\linewidth}{\wd\pandoc@box}%
  \ifdim\@tempb\p@<\@tempa\p@\let\@tempa\@tempb\fi% select the smaller of both
  \ifdim\@tempa\p@<\p@\scalebox{\@tempa}{\usebox\pandoc@box}%
  \else\usebox{\pandoc@box}%
  \fi%
}
% Set default figure placement to htbp
\def\fps@figure{htbp}
\makeatother


% definitions for citeproc citations
\NewDocumentCommand\citeproctext{}{}
\NewDocumentCommand\citeproc{mm}{%
  \begingroup\def\citeproctext{#2}\cite{#1}\endgroup}
\makeatletter
 % allow citations to break across lines
 \let\@cite@ofmt\@firstofone
 % avoid brackets around text for \cite:
 \def\@biblabel#1{}
 \def\@cite#1#2{{#1\if@tempswa , #2\fi}}
\makeatother
\newlength{\cslhangindent}
\setlength{\cslhangindent}{1.5em}
\newlength{\csllabelwidth}
\setlength{\csllabelwidth}{3em}
\newenvironment{CSLReferences}[2] % #1 hanging-indent, #2 entry-spacing
 {\begin{list}{}{%
  \setlength{\itemindent}{0pt}
  \setlength{\leftmargin}{0pt}
  \setlength{\parsep}{0pt}
  % turn on hanging indent if param 1 is 1
  \ifodd #1
   \setlength{\leftmargin}{\cslhangindent}
   \setlength{\itemindent}{-1\cslhangindent}
  \fi
  % set entry spacing
  \setlength{\itemsep}{#2\baselineskip}}}
 {\end{list}}
\usepackage{calc}
\newcommand{\CSLBlock}[1]{\hfill\break\parbox[t]{\linewidth}{\strut\ignorespaces#1\strut}}
\newcommand{\CSLLeftMargin}[1]{\parbox[t]{\csllabelwidth}{\strut#1\strut}}
\newcommand{\CSLRightInline}[1]{\parbox[t]{\linewidth - \csllabelwidth}{\strut#1\strut}}
\newcommand{\CSLIndent}[1]{\hspace{\cslhangindent}#1}

\ifLuaTeX
\usepackage[bidi=basic]{babel}
\else
\usepackage[bidi=default]{babel}
\fi
% get rid of language-specific shorthands (see #6817):
\let\LanguageShortHands\languageshorthands
\def\languageshorthands#1{}


\setlength{\emergencystretch}{3em} % prevent overfull lines

\providecommand{\tightlist}{%
  \setlength{\itemsep}{0pt}\setlength{\parskip}{0pt}}



 


\usepackage{url}
\def\UrlBreaks{\do\/\do-}
\usepackage{adjustbox}
\def\pandocbounded#1{\adjustbox{max width=\linewidth, keepaspectratio}{#1}}
\usepackage{microtype}
\usepackage{amsmath,amssymb}
\usepackage{booktabs}
\usepackage{longtable}
\usepackage{array}
\usepackage{float}
\usepackage{subcaption}
\usepackage{xcolor}
\usepackage{fancyvrb}
\usepackage{csquotes}
\usepackage{emoji}
\setemojifont{TwemojiMozilla}
\definecolor{pdi-blue}{RGB}{21,101,192}
\definecolor{pdi-green}{RGB}{46,125,50}
\makeatletter
\@ifpackageloaded{tcolorbox}{}{\usepackage[skins,breakable]{tcolorbox}}
\@ifpackageloaded{fontawesome5}{}{\usepackage{fontawesome5}}
\definecolor{quarto-callout-color}{HTML}{909090}
\definecolor{quarto-callout-note-color}{HTML}{0758E5}
\definecolor{quarto-callout-important-color}{HTML}{CC1914}
\definecolor{quarto-callout-warning-color}{HTML}{EB9113}
\definecolor{quarto-callout-tip-color}{HTML}{00A047}
\definecolor{quarto-callout-caution-color}{HTML}{FC5300}
\definecolor{quarto-callout-color-frame}{HTML}{acacac}
\definecolor{quarto-callout-note-color-frame}{HTML}{4582ec}
\definecolor{quarto-callout-important-color-frame}{HTML}{d9534f}
\definecolor{quarto-callout-warning-color-frame}{HTML}{f0ad4e}
\definecolor{quarto-callout-tip-color-frame}{HTML}{02b875}
\definecolor{quarto-callout-caution-color-frame}{HTML}{fd7e14}
\makeatother
\makeatletter
\@ifpackageloaded{bookmark}{}{\usepackage{bookmark}}
\makeatother
\makeatletter
\@ifpackageloaded{caption}{}{\usepackage{caption}}
\AtBeginDocument{%
\ifdefined\contentsname
  \renewcommand*\contentsname{Índice}
\else
  \newcommand\contentsname{Índice}
\fi
\ifdefined\listfigurename
  \renewcommand*\listfigurename{Lista de Figuras}
\else
  \newcommand\listfigurename{Lista de Figuras}
\fi
\ifdefined\listtablename
  \renewcommand*\listtablename{Lista de Tabelas}
\else
  \newcommand\listtablename{Lista de Tabelas}
\fi
\ifdefined\figurename
  \renewcommand*\figurename{Figura}
\else
  \newcommand\figurename{Figura}
\fi
\ifdefined\tablename
  \renewcommand*\tablename{Tabela}
\else
  \newcommand\tablename{Tabela}
\fi
}
\@ifpackageloaded{float}{}{\usepackage{float}}
\floatstyle{ruled}
\@ifundefined{c@chapter}{\newfloat{codelisting}{h}{lop}}{\newfloat{codelisting}{h}{lop}[chapter]}
\floatname{codelisting}{Listagem}
\newcommand*\listoflistings{\listof{codelisting}{Lista de Listagens}}
\makeatother
\makeatletter
\makeatother
\makeatletter
\@ifpackageloaded{caption}{}{\usepackage{caption}}
\@ifpackageloaded{subcaption}{}{\usepackage{subcaption}}
\makeatother
\usepackage{bookmark}
\IfFileExists{xurl.sty}{\usepackage{xurl}}{} % add URL line breaks if available
\urlstyle{same}
\hypersetup{
  pdftitle={Processamento Digital de Imagens e Visão Computacional},
  pdfauthor={Francisco de Assis Zampirolli},
  pdflang={pt},
  colorlinks=true,
  linkcolor={blue},
  filecolor={Maroon},
  citecolor={Blue},
  urlcolor={blue},
  pdfcreator={LaTeX via pandoc}}


\title{Processamento Digital de Imagens e Visão Computacional}
\usepackage{etoolbox}
\makeatletter
\providecommand{\subtitle}[1]{% add subtitle to \maketitle
  \apptocmd{\@title}{\par {\large #1 \par}}{}{}
}
\makeatother
\subtitle{Abordagem Prática com Python}
\author{Francisco de Assis Zampirolli}
\date{2026-05-10}
\begin{document}
\frontmatter
\maketitle

\renewcommand*\contentsname{Índice}
{
\hypersetup{linkcolor=}
\setcounter{tocdepth}{2}
\tableofcontents
}

\mainmatter
\bookmarksetup{startatroot}

\chapter*{PDI e VC --- Python}\label{pdi-e-vc-python}
\addcontentsline{toc}{chapter}{PDI e VC --- Python}

\markboth{PDI e VC --- Python}{PDI e VC --- Python}

Bem-vindo ao livro de PDI e Visão Computacional --- versão Python /
Português.

\subsection*{Organização}\label{organizauxe7uxe3o}
\addcontentsline{toc}{subsection}{Organização}

\begin{itemize}
\tightlist
\item
  \textbf{Parte I --- Fundamentos de PDI} --- Representação,
  histogramas, filtragem, morfologia
\item
  \textbf{Parte II --- Visão Computacional} --- Segmentação,
  descritores, detecção, deep learning
\end{itemize}

\begin{center}\rule{0.5\linewidth}{0.5pt}\end{center}

\bookmarksetup{startatroot}

\chapter*{Prefácio}\label{prefuxe1cio}
\addcontentsline{toc}{chapter}{Prefácio}

\markboth{Prefácio}{Prefácio}

Este livro é um \textbf{projeto em construção} sobre
\textbf{Processamento Digital de Imagens (PDI) e Visão Computacional
(VC)}, concebido como material didático interativo para cursos de
graduação e pós-graduação em Computação, Engenharias e áreas afins.

\begin{quote}
\textbf{Destaque:} Este livro segue a metodologia de Zampirolli et al.
(2025) --- uma versão estendida de Zampirolli et al. (2024), que recebeu
o \textbf{prêmio de melhor artigo} na trilha de ``Recursos e Ambientes
Educacionais'' na \textbf{EduComp 2024} --- com a biblioteca
\texttt{morph.py}, desenvolvida pelo autor.
\end{quote}

\section*{Um livro vivo}\label{um-livro-vivo}
\addcontentsline{toc}{section}{Um livro vivo}

\markright{Um livro vivo}

Este não é um livro estático. O conteúdo está em constante evolução:
exemplos são refinados, novas seções são adicionadas e abordagens
pedagógicas são aprimoradas com base no feedback de alunos e
professores. Por isso, tratamos este material como um \emph{projeto
inicial} --- uma base sólida que continuará crescendo.

\section*{Uso de ferramentas de Inteligência
Artificial}\label{uso-de-ferramentas-de-inteliguxeancia-artificial}
\addcontentsline{toc}{section}{Uso de ferramentas de Inteligência
Artificial}

\markright{Uso de ferramentas de Inteligência Artificial}

Este livro foi produzido com o apoio de \textbf{ferramentas de
Inteligência Artificial} (IA) para aprimorar a qualidade do texto e do
código. A versão inicial dos capítulos passou por um processo de revisão
e correção utilizando múltiplos assistentes de IA, incluindo
\textbf{DeepSeek}, \textbf{Gemini} e \textbf{Claude}. Por meio de
\emph{prompts} iterativos, essas ferramentas ajudaram a:

\begin{itemize}
\tightlist
\item
  Refinar a clareza e a precisão das explicações;
\item
  Detectar e corrigir erros de sintaxe, lógica e digitação nos exemplos
  de código;
\item
  Sugerir melhorias na organização e na fluência do texto.
\end{itemize}

Além disso, \textbf{algumas imagens e ilustrações presentes no livro
foram geradas pelo Gemini}, sempre com o objetivo de facilitar a
compreensão de conceitos abstratos da área.

\begin{quote}
\textbf{Transparência:} O uso de IA não substitui a revisão final do
autor, nem a validação pelos colaboradores humanos. O material é
constantemente testado e ajustado para garantir correção técnica e
adequação didática. Acreditamos que a parceria entre autor e ferramentas
de IA enriquece o processo de criação, desde que usada de forma ética e
criteriosa.
\end{quote}

\section*{Como este livro é
produzido}\label{como-este-livro-uxe9-produzido}
\addcontentsline{toc}{section}{Como este livro é produzido}

\markright{Como este livro é produzido}

Todo o conteúdo é escrito em \textbf{Quarto}, um sistema de publicação
científica e técnica que permite renderizar o mesmo código-fonte para
múltiplos formatos:

\begin{longtable}[]{@{}
  >{\raggedright\arraybackslash}p{(\linewidth - 2\tabcolsep) * \real{0.4500}}
  >{\raggedright\arraybackslash}p{(\linewidth - 2\tabcolsep) * \real{0.5500}}@{}}
\toprule\noalign{}
\begin{minipage}[b]{\linewidth}\raggedright
Formato
\end{minipage} & \begin{minipage}[b]{\linewidth}\raggedright
Descrição
\end{minipage} \\
\midrule\noalign{}
\endhead
\bottomrule\noalign{}
\endlastfoot
\textbf{HTML} & Versão web completa, com navegação interativa,
disponível em
\href{https://fzampirolli.github.io/pdi-vc/}{fzampirolli.github.io/pdi-vc} \\
\textbf{PDF} & Versão para impressão ou leitura offline --- ex.
\href{https://github.com/fzampirolli/pdi-vc/blob/master/gen/book/py.pt/livro.pt.py.pdf}{gen/book/py.pt/livro.pt.py.pdf} \\
\textbf{Notebooks para alunos} & Versão \texttt{.ipynb} processada por
um filtro personalizado que resolve citações, numera figuras, tabelas e
equações, e gera referências formatadas no estilo ABNT --- ideal para
uso no Jupyter e Google Colab \\
\end{longtable}

O filtro de pós-processamento (\texttt{gerar\_notebooks\_alunos.py})
prepara os notebooks para distribuição aos alunos, garantindo que
funcionem perfeitamente em ambientes interativos mesmo sem depender do
Quarto instalado.

\section*{Sobre as linguagens de programação e
idiomas}\label{sobre-as-linguagens-de-programauxe7uxe3o-e-idiomas}
\addcontentsline{toc}{section}{Sobre as linguagens de programação e
idiomas}

\markright{Sobre as linguagens de programação e idiomas}

O \textbf{núcleo original} deste livro é escrito em \textbf{Português}
com exemplos em \textbf{Python}. Esta combinação (\texttt{pt\ +\ py}) é
a versão de referência, mantida e revisada diretamente pelos autor.

Para ampliar o alcance do material, versões em outras combinações de
idioma e linguagem de programação são geradas com o auxílio de
\textbf{APIs de Inteligência Artificial} (modelos de tradução e
conversão de código). Essas versões passam por um pipeline automatizado
que:

\begin{enumerate}
\def\labelenumi{\arabic{enumi}.}
\tightlist
\item
  \textbf{Traduz} o texto explicativo para o idioma de destino
\item
  \textbf{Converte} os exemplos de código Python para a linguagem alvo
  (ex: C++, R, Julia)
\item
  \textbf{Aplica} o mesmo pós-processamento de citações, figuras,
  tabelas e equações
\end{enumerate}

\subsection*{Exemplos de uso}\label{exemplos-de-uso}
\addcontentsline{toc}{subsection}{Exemplos de uso}

O projeto utiliza um \texttt{Makefile} com atalhos para desenvolvimento.
Veja alguns comandos comuns:

\begin{Shaded}
\begin{Highlighting}[]
\CommentTok{\# Versão de referência: Português + Python, HTML + PDF}
\FunctionTok{make}\NormalTok{ build{-}all}

\CommentTok{\# Apenas HTML (atualiza automaticamente ao salvar arquivos)}
\FunctionTok{make}\NormalTok{ html}

\CommentTok{\# Português + C++, apenas PDF}
\FunctionTok{make}\NormalTok{ build{-}pdf LANGS=cpp LOCALES=pt}

\CommentTok{\# Inglês + Python, apenas HTML}
\FunctionTok{make}\NormalTok{ build LANGS=py LOCALES=en}

\CommentTok{\# Inglês + C++, HTML + PDF}
\FunctionTok{make}\NormalTok{ build{-}all LANGS=cpp LOCALES=en}

\CommentTok{\# Todas as combinações disponíveis (py, cpp, java, c × pt, en, fr, es, it)}
\FunctionTok{make}\NormalTok{ full}

\CommentTok{\# Gerar notebooks para alunos (distribuição Jupyter/Colab)}
\FunctionTok{make}\NormalTok{ alunos}

\CommentTok{\# Publicar no GitHub Pages}
\FunctionTok{make}\NormalTok{ publish}
\end{Highlighting}
\end{Shaded}

\begin{quote}
\textbf{Dica:} Você pode combinar qualquer idioma com qualquer
linguagem. Os idiomas disponíveis são: \texttt{pt} (Português),
\texttt{en} (Inglês), \texttt{fr} (Francês), \texttt{es} (Espanhol),
\texttt{it} (Italiano). As linguagens disponíveis são: \texttt{py}
(Python), \texttt{cpp} (C++), \texttt{java} (Java), \texttt{c} (C).
\end{quote}

\begin{quote}
\textbf{Nota importante:} As versões geradas por IA são um \textbf{ponto
de partida} e podem conter imprecisões. Elas são disponibilizadas como
material complementar e devem ser revisadas antes do uso em sala de
aula. Contribuições da comunidade para refinar essas versões são muito
bem-vindas!
\end{quote}

\section*{Código aberto}\label{cuxf3digo-aberto}
\addcontentsline{toc}{section}{Código aberto}

\markright{Código aberto}

Este livro é um projeto de código aberto. Todo o conteúdo --- texto,
código-fonte, imagens e scripts de processamento --- está disponível
publicamente em:

\textbf{\href{https://github.com/fzampirolli/pdi-vc/}{github.com/fzampirolli/pdi-vc}}

Você pode:

\begin{itemize}
\tightlist
\item
  \textbf{Ler e estudar} gratuitamente
\item
  \textbf{Reportar erros} ou sugerir melhorias abrindo uma \emph{issue}
\item
  \textbf{Contribuir} com correções, novos exemplos ou traduções via
  \emph{pull requests}
\item
  \textbf{Adaptar} o material para suas próprias turmas
\end{itemize}

\section*{Como usar os notebooks}\label{como-usar-os-notebooks}
\addcontentsline{toc}{section}{Como usar os notebooks}

\markright{Como usar os notebooks}

Cada capítulo do livro está disponível como um \emph{notebook}
interativo. Se você nunca usou um notebook Jupyter ou Google Colab, veja
a seção \textbf{``Antes de começar: Notebooks em Python''} a seguir para
um guia rápido.

\section*{Agradecimentos}\label{agradecimentos}
\addcontentsline{toc}{section}{Agradecimentos}

\markright{Agradecimentos}

Agradecemos aos alunos, colegas e colaboradores que testam, questionam e
contribuem com este projeto. Este livro é feito para vocês e por vocês.

\begin{center}\rule{0.5\linewidth}{0.5pt}\end{center}

\section*{Antes de começar: Notebooks em
Python}\label{antes-de-comeuxe7ar-notebooks-em-python}
\addcontentsline{toc}{section}{Antes de começar: Notebooks em Python}

\markright{Antes de começar: Notebooks em Python}

Você já ouviu falar em \textbf{\emph{Literate Programming}} (Programação
Literária)? O conceito foi criado por Donald Knuth em 1984 (Knuth, 1984)
--- o mesmo que desenvolveu o sistema \textbf{TeX} para tipografia
digital. Knuth iniciou esse projeto no final da década de 70 após ficar
insatisfeito com a qualidade gráfica das novas edições de sua série
clássica, \emph{The Art of Computer Programming} (Knuth, 1997), buscando
uma forma de unir a precisão matemática à beleza estética dos livros
antigos. O conceito de Programação Literária não propõe apenas
``combinar código e documentação'', mas sim inverter a lógica: o
programa deve ser escrito para ser lido por humanos (como um ensaio), e
as partes de código são extraídas dele para a máquina. É um detalhe
sutil, mas que reforça o termo ``literária''.

O conteúdo deste livro é formado por \emph{Notebook}: um documento que
intercala textos explicativos --- as \textbf{células de texto} no
formato
\href{https://colab.research.google.com/notebooks/markdown_guide.ipynb}{Markdown}
--- com códigos de programas em \href{https://www.python.org/}{Python}
--- as \textbf{células de código}. Mas aceita também outras linguagens
de programação, mas o processo de execução é mais complicado.

É fácil distinguir as células de código: elas são precedidas por
\texttt{{[}\ {]}}. Para executar uma célula de código, basta
selecioná-la e pressionar Shift + Enter. Antes de começar, algumas
observações importantes sobre o ambiente:

\begin{itemize}
\tightlist
\item
  \textbf{No computador:} no Colab ou Jupyter, você também pode executar
  uma célula clicando no botão de $\triangleright$ (\emph{play}) que
  aparece ao passar o mouse sobre os colchetes \texttt{{[}\ {]}}.
\item
  \textbf{No celular:} no Colab, o atalho Shift + Enter pode não
  funcionar. Nesse caso, toque no botão de $\triangleright$
  (\emph{play}).
\item
  \textbf{Execução local:} você pode rodar este notebook no seu
  computador instalando o \href{https://jupyter.org}{Jupyter Notebook}.
\item
  \textbf{Google Colab:} para executar no seu navegador, vá em
  \textbf{Arquivo → Salvar uma cópia no Drive} e abra o arquivo
  \textbf{Cópia de aula01.ipynb}, que ficará na pasta \textbf{Colab
  Notebooks} do seu Drive.
\end{itemize}

O resultado será exibido logo abaixo da célula executada.

\begin{tcolorbox}[enhanced jigsaw, titlerule=0mm, left=2mm, opacityback=0, toprule=.15mm, arc=.35mm, opacitybacktitle=0.6, colbacktitle=quarto-callout-note-color!10!white, colframe=quarto-callout-note-color-frame, rightrule=.15mm, breakable, coltitle=black, bottomtitle=1mm, toptitle=1mm, title=\textcolor{quarto-callout-note-color}{\faInfo}\hspace{0.5em}{Nota sobre o formato}, colback=white, bottomrule=.15mm, leftrule=.75mm]

Você pode estar acessando este conteúdo de diferentes maneiras. Se
estiver executando este notebook no \href{https://jupyter.org/}{Jupyter}
ou no \href{https://colab.research.google.com/}{Google Colab}, as
células de código são interativas: você pode executá-las pressionando
Shift + Enter ou clicando no botão de \emph{play}. Se estiver lendo a
versão renderizada em HTML ou PDF, o código aparecerá como blocos
estáticos, mas o conteúdo e as explicações permanecem os mesmos --- a
única diferença é que você não poderá executar o código diretamente no
documento.

\end{tcolorbox}

\part{Parte I --- Fundamentos de PDI}

\chapter{Fundamentos e Primeiros
Passos}\label{fundamentos-e-primeiros-passos}

\href{https://colab.research.google.com/github/fzampirolli/pdi-vc/blob/master/notebooks_alunos/cap01/cap01_aluno.ipynb}{
\includegraphics[width=0.16\linewidth,height=\textheight,keepaspectratio]{cap01/imagens/colab-badge.png}}
\href{https://github.com/fzampirolli/pdi-vc}{
\includegraphics[width=0.2\linewidth,height=\textheight,keepaspectratio]{cap01/imagens/github-badge.png}}

Este capítulo inaugura a \textbf{Parte 1} do livro, focada nos
fundamentos do Processamento Digital de Imagens (PDI). Aqui, você
aprenderá como as imagens são representadas matematicamente e como
manipulá-las usando Python e a biblioteca \texttt{morph.py} (Zampirolli
et al., 2025).

\section{Objetivos}\label{objetivos}

Ao final deste capítulo, você será capaz de:

\begin{itemize}
\tightlist
\item
  Compreender a natureza física e matemática da imagem digital
  \(f(x,y)\).
\item
  Identificar as faixas do espectro eletromagnético relevantes para PDI.
\item
  Configurar o ambiente de desenvolvimento em Python.
\item
  Realizar operações básicas: leitura, exibição e salvamento de imagens.
\item
  Manipular estruturas de matrizes (NumPy) sem cair em armadilhas de
  memória.
\item
  Acessar e modificar intensidades de pixels individualmente.
\item
  Aplicar limiarização manual.
\end{itemize}

\section{\texorpdfstring{Antes de começar: \emph{Notebooks} em
Python}{Antes de começar: Notebooks em Python}}\label{antes-de-comeuxe7ar-notebooks-em-python-1}

Este material foi construído sob o conceito de \textbf{\emph{Literate
Programming}} (Programação Literária), idealizado por Donald Knuth na
década de 84 (Knuth, 1984). Knuth --- também criador do sistema
\textbf{TeX} para tipografia digital --- propôs que os programas fossem
escritos como uma narrativa lógica para seres humanos, intercalando
código e documentação.

Para executar uma célula, pressione Shift + Enter ou clique no botão
$\triangleright$.

\begin{tcolorbox}[enhanced jigsaw, titlerule=0mm, left=2mm, opacityback=0, toprule=.15mm, arc=.35mm, opacitybacktitle=0.6, colbacktitle=quarto-callout-note-color!10!white, colframe=quarto-callout-note-color-frame, rightrule=.15mm, breakable, coltitle=black, bottomtitle=1mm, toptitle=1mm, title=\textcolor{quarto-callout-note-color}{\faInfo}\hspace{0.5em}{Nota sobre o formato}, colback=white, bottomrule=.15mm, leftrule=.75mm]

Se você estiver lendo a versão renderizada (PDF ou HTML), o código
aparecerá como blocos estáticos. Para interagir, acesse o link do
\textbf{Google Colab} no topo da página.

\end{tcolorbox}

\section{Fundamentos}\label{fundamentos}

O estudo de sistemas baseados em imagens compreende um ecossistema de
disciplinas integradas que transformam dados visuais brutos em
conhecimento estruturado. Enquanto algumas áreas focam na geração de
representações, outras dedicam-se ao tratamento e à análise desses dados
para dar suporte a aplicações tecnológicas complexas.

O diagrama apresentado no Figura~\ref{fig-1-representacao} estabelece a
distinção e complementaridade entre o Processamento de Imagens (PDI) e a
Visão Computacional (VC). O PDI, destacado em \textbf{verde}, tem como
foco a transformação de imagem para imagem, visando melhoria de
qualidade ou pré-processamento, como a remoção de ruídos e realce de
contraste.

Em contrapartida, a VC, assinalada em \textbf{azul}, foca na
interpretação do conteúdo visual para extrair modelos ou informações,
como o reconhecimento de objetos e gestos. A região de intersecção
ilustra a sinergia entre as áreas, onde o PDI prepara os dados visuais
para a interpretação pela VC. O mapa também demonstra as interconexões
de ambas as disciplinas com áreas como Robótica, Computação Gráfica,
Inteligência Artificial e Neurociência.

\begin{figure}

\centering{

\includegraphics[width=1\linewidth,height=\textheight,keepaspectratio]{cap01/imagens/fig-1-pdf-vc.png}

}

\caption{\label{fig-1-representacao}Diagrama relacional detalhando as
distinções fundamentais, sinergias e interconexões entre Processamento
de Imagens e Visão Computacional no contexto de sistemas baseados em
imagens.}

\end{figure}%

\subsection{\texorpdfstring{\emoji{eye} Visão
Computacional}{ Visão Computacional}}\label{visuxe3o-computacional}

\begin{itemize}
\tightlist
\item
  \textbf{Foco:} \emph{Imagem → Modelo} (caminho inverso da Computação
  Gráfica).\\
\item
  \textbf{Objetivo:} Extrair informação de alto nível a partir de
  imagens ou vídeos.\\
\item
  \textbf{Aplicações típicas:}

  \begin{itemize}
  \tightlist
  \item
    Robótica -- detecção de obstáculos, localização e navegação
    autônoma.
  \item
    Vigilância e inspeção -- reconhecimento de eventos, leitura de
    placas, controle de qualidade.
  \item
    Sensoriamento remoto -- análise de imagens de satélite, mapeamento
    ambiental.
  \item
    Imagens médicas -- detecção de tumores, segmentação de órgãos,
    auxílio ao diagnóstico.
  \item
    Interação humano‑computador -- reconhecimento de gestos, expressões
    faciais, rastreamento ocular.
  \end{itemize}
\item
  \textbf{Relação com outras áreas:} utiliza técnicas de Aprendizado de
  Máquina e IA para classificar e interpretar cenas; serve como
  ``olhos'' da Robótica.
\end{itemize}

\subsection{\texorpdfstring{\emoji{framed-picture} Processamento Digital
de Imagens
(PDI)}{ Processamento Digital de Imagens (PDI)}}\label{processamento-digital-de-imagens-pdi}

\begin{itemize}
\tightlist
\item
  \textbf{Foco:} \emph{Imagem → Imagem} (geralmente -- transformação de
  uma imagem em outra).\\
\item
  \textbf{Objetivo:} Melhorar a qualidade visual ou extrair
  características de baixo nível.\\
\item
  \textbf{Aplicações comuns:}

  \begin{itemize}
  \tightlist
  \item
    Eliminação de ruídos (filtros de média, mediana, gaussiano).
  \item
    Melhoria de contraste (equalização de histograma, ajuste gamma).
  \item
    Detecção de bordas (Sobel, Canny, Laplaciano).
  \item
    Segmentação (limiarização, crescimento de regiões,
    \emph{watershed}).
  \item
    Transformações geométricas (redimensionamento, rotação, correção de
    perspectiva).
  \end{itemize}
\item
  \textbf{Relação com outras áreas:}

  \begin{itemize}
  \tightlist
  \item
    É a \textbf{base} para a maioria dos sistemas de Visão Computacional
    (pré‑processamento).\\
  \item
    A Computação Gráfica frequentemente aplica PDI para
    pós‑processamento (ex.: suavização, realce).\\
  \item
    Técnicas de IA podem otimizar parâmetros de processamento (ex.:
    aprendizado de filtros).
  \end{itemize}
\end{itemize}

\subsection{Como as ``Outras Áreas'' se
conectam}\label{como-as-outras-uxe1reas-se-conectam}

\begin{longtable}[]{@{}
  >{\raggedright\arraybackslash}p{(\linewidth - 2\tabcolsep) * \real{0.2389}}
  >{\raggedright\arraybackslash}p{(\linewidth - 2\tabcolsep) * \real{0.7611}}@{}}
\toprule\noalign{}
\begin{minipage}[b]{\linewidth}\raggedright
Área
\end{minipage} & \begin{minipage}[b]{\linewidth}\raggedright
Relação com PDI e VC
\end{minipage} \\
\midrule\noalign{}
\endhead
\bottomrule\noalign{}
\endlastfoot
\textbf{Inteligência Artificial} & Fornece modelos (redes neurais, SVM)
que interpretam saídas da VC. \\
\textbf{Robótica} & Consome dados de VC para tomar decisões (navegação,
manipulação). \\
\textbf{Aprendizado de Máquina} & Usa descritores extraídos pelo PDI/VC
para treinar classificadores. \\
\textbf{Computação Gráfica} & Caminho inverso: \emph{modelo → imagem};
muitas vezes aplica PDI para renderização realista. \\
\textbf{Neurociência} & Inspira modelos de PDI (ex.: filtros semelhantes
a células ganglionares da retina). \\
\end{longtable}

\section{Etapas do PDI}\label{etapas-do-pdi}

As etapas do PDI são apresentadas na Figura~\ref{fig-1-etapas}, que
podem ser compreendidas como uma cadeia de transformações que reduz a
redundância dos dados em busca de significado:

\begin{itemize}
\tightlist
\item
  \textbf{Baixo Nível:} Atua diretamente sobre os pixels da imagem
  ruidosa para realizar melhorias e filtragens, gerando como saída uma
  imagem limpa ou realçada.
\item
  \textbf{Médio Nível:} Recebe a imagem tratada e realiza a segmentação
  e descrição, transformando a matriz de pixels em atributos
  estruturados (forma, tamanho e textura).
\item
  \textbf{Alto Nível:} Utiliza a tabela de atributos para alimentar
  processos de lógica e inteligência artificial, resultando na decisão
  ou reconhecimento final (como o diagnóstico médico).
\end{itemize}

\begin{figure}

\centering{

\includegraphics[width=1\linewidth,height=\textheight,keepaspectratio]{cap01/imagens/fig-1-etapas.png}

}

\caption{\label{fig-1-etapas}Representação do fluxo sequencial de
processamento: a saída de cada nível torna-se a entrada do nível
subsequente.}

\end{figure}%

Detalhando mais essas etapas, a Figura~\ref{fig-1-etapas2} ilustra a
jornada completa da informação visual, desde a sua captura até a
interpretação final. O fluxo inicia-se na \textbf{Aquisição da Imagem
(1)} e avança por processos de tratamento como o \textbf{Melhoramento
(2)} e a \textbf{Restauração (3)}. Em seguida, o conteúdo é isolado
através da \textbf{Segmentação (4)} e refinado pela \textbf{Morfologia
(5)}. O estágio crucial de transição ocorre na \textbf{Representação e
Descrição (6)}, onde objetos visuais são convertidos em dados
matemáticos (como área e perímetro), permitindo que o sistema execute o
\textbf{Reconhecimento do Objeto (7)}. Este ecossistema é suportado por
pilares como o \textbf{Processamento de Imagem Colorida} e a
\textbf{Compressão}, que garantem a eficiência no armazenamento e
análise dos dados.

\begin{figure}

\centering{

\includegraphics[width=1\linewidth,height=\textheight,keepaspectratio]{cap01/imagens/fig-1-etapas2.png}

}

\caption{\label{fig-1-etapas2}Fluxo detalhado do PDI: da aquisição
sensorial à extração de atributos e reconhecimento automatizado,
incluindo em processamento colorido e compressão.}

\end{figure}%

\section{Formação da Imagem e o
Espectro}\label{formauxe7uxe3o-da-imagem-e-o-espectro}

O processo de formação de uma imagem é fundamentado na interação entre a
matéria e a energia radiante. Essencialmente, uma imagem é concebida
quando um \textbf{sensor registra a radiação} resultante da interação
com um objeto físico. No contexto da visão humana e da fotografia
convencional, este fenômeno depende de uma fonte de luz que ilumine a
cena; as características dos objetos são então codificadas através das
\textbf{variações de intensidade e cor} da luz que atinge o sensor,
conforme ilustrado na Figura~\ref{fig-1-luz}.

\begin{figure}

\centering{

\includegraphics[width=1\linewidth,height=\textheight,keepaspectratio]{cap01/imagens/fig-1-luz.png}

}

\caption{\label{fig-1-luz}Representação do espectro visível e sua
posição em relação às demais radiações eletromagnéticas, destacando a
variação dos comprimentos de onda de 380 nm a 750 nm.}

\end{figure}%

A \textbf{luz visível} ocupa apenas uma pequena faixa do espectro
eletromagnético --- entre \textbf{380 nm (violeta)} e \textbf{750 nm
(vermelho)} ---, conforme ilustrado na Figura~\ref{fig-1-espectro}.
Sensores digitais convencionais operam nessa mesma janela, mas
equipamentos especializados podem captar radiações invisíveis ao olho
humano, como o infravermelho e os raios X. Em PDI, a imagem formada
depende diretamente da sensibilidade espectral do sensor utilizado.

\begin{figure}

\centering{

\pandocbounded{\includegraphics[keepaspectratio]{cap01/cap01.py.pt_files/figure-pdf/fig-1-espectro-output-1.png}}

}

\caption{\label{fig-1-espectro}(A) Espectro eletromagnético completo em
escala logarítmica, com destaque para a faixa visível. (B) Detalhe da
luz visível (380--750 nm) e suas cores componentes. (C) Decomposição da
luz branca pelo prisma: menor comprimento de onda sofre maior refração,
separando UV, visível e infravermelho.}

\end{figure}%

\section{O que é uma Imagem
Digital?}\label{o-que-uxe9-uma-imagem-digital}

Uma \textbf{imagem digital} é formada por uma grade de \textbf{pixels}
(\emph{Picture Elements}), onde cada pixel é a menor unidade elementar
da imagem.

\begin{tcolorbox}[enhanced jigsaw, titlerule=0mm, left=2mm, opacityback=0, toprule=.15mm, arc=.35mm, opacitybacktitle=0.6, colbacktitle=quarto-callout-tip-color!10!white, colframe=quarto-callout-tip-color-frame, rightrule=.15mm, breakable, coltitle=black, bottomtitle=1mm, toptitle=1mm, title=\textcolor{quarto-callout-tip-color}{\faLightbulb}\hspace{0.5em}{O que é um Pixel?}, colback=white, bottomrule=.15mm, leftrule=.75mm]

Um \textbf{pixel} é a menor unidade que compõe uma imagem digital. Cada
pixel ocupa uma posição única na grade e armazena um ou mais valores
numéricos que representam sua intensidade ou cor.

\end{tcolorbox}

\subsection*{Representação
Matemática}\label{representauxe7uxe3o-matemuxe1tica}
\addcontentsline{toc}{subsection}{Representação Matemática}

Diferente de uma função contínua, o domínio de uma imagem digital é um
plano retangular finito \(\mathbb{E} \subset \mathbb{Z}^2\), que
representa a grade de amostragem. Este domínio é indexado por
coordenadas inteiras:

\begin{equation}\phantomsection\label{eq-1-dominio}{
\mathbb{E}  = \{ (x, y) \in \mathbb{Z}^2 \mid 0 \le x < L,\; 0 \le y < H \}
}\end{equation}

Onde:

\begin{itemize}
\tightlist
\item
  \textbf{\(L\)}: representa a largura da imagem (número de colunas).
\item
  \textbf{\(H\)}: representa a altura da imagem (número de linhas).
\end{itemize}

A imagem digital é uma função que associa cada par de coordenadas
\((x,y)\) a um ou mais valores que descrevem a aparência do pixel.

\begin{equation}\phantomsection\label{eq-1-image}{
f: \mathbb{E}  \to \mathcal{V}
}\end{equation}

O conjunto \(\mathcal{V}\) define os valores possíveis para o pixel
(codomínio), variando conforme o tipo de imagem, como demonstrado na
Tabela~\ref{tbl-1-tipos-imagem}.

\begin{longtable}[]{@{}llc@{}}
\caption{Principais tipos de imagem digital e seus respectivos conjuntos
de valores possíveis para cada
pixel.}\label{tbl-1-tipos-imagem}\tabularnewline
\toprule\noalign{}
Tipo de imagem & \(\mathcal{V}\) (valores do pixel) & Representação \\
\midrule\noalign{}
\endfirsthead
\toprule\noalign{}
Tipo de imagem & \(\mathcal{V}\) (valores do pixel) & Representação \\
\midrule\noalign{}
\endhead
\bottomrule\noalign{}
\endlastfoot
\textbf{Binária} & \(\{0, 1\}\) ou \(\{0, 255\}\) &
$\blacksquare$$\square$ \\
\textbf{Tons de cinza} & \(\{0, 1, \dots, 255\}\) &
\texttt{[~]}\texttt{[=]}\texttt{[\#]}$\blacksquare$ \\
\textbf{Colorida (RGB)} & \(\{0, \dots, 255\}^3\) (triplas) &
\textcolor{red}{$\blacksquare$}\textcolor{green}{$\blacksquare$}\textcolor{blue}{$\blacksquare$} \\
\end{longtable}

\textbf{Exemplo prático:} Uma imagem colorida (RGB) é representada
matematicamente por uma função que retorna três valores para cada pixel.
No computador, isso resulta em três matrizes sobrepostas (canais R, G e
B), onde cada célula contém um valor de intensidade para aquele canal
específico no pixel.

\section{Configuração do Ambiente}\label{configurauxe7uxe3o-do-ambiente}

Utilizaremos o ecossistema científico do Python, com destaque para o
\textbf{NumPy} (matemática matricial), \textbf{OpenCV} (visão
computacional padrão de mercado) e a biblioteca \textbf{morph.py}
(Zampirolli et al., 2025), desenvolvida especificamente para fins
didáticos neste livro.

\subsection{\texorpdfstring{\emoji{package} Bibliotecas
utilizadas}{ Bibliotecas utilizadas}}\label{bibliotecas-utilizadas}

\begin{longtable}[]{@{}
  >{\raggedright\arraybackslash}p{(\linewidth - 2\tabcolsep) * \real{0.4000}}
  >{\raggedright\arraybackslash}p{(\linewidth - 2\tabcolsep) * \real{0.6000}}@{}}
\toprule\noalign{}
\begin{minipage}[b]{\linewidth}\raggedright
Biblioteca
\end{minipage} & \begin{minipage}[b]{\linewidth}\raggedright
Função principal
\end{minipage} \\
\midrule\noalign{}
\endhead
\bottomrule\noalign{}
\endlastfoot
\texttt{numpy} & Representação matricial de imagens \\
\texttt{opencv-python} & Leitura, escrita e operações de visão
computacional \\
\texttt{matplotlib} & Visualização de imagens e gráficos \\
\texttt{morph.py} & Abstração didática das operações de PDI \\
\end{longtable}

\begin{tcolorbox}[enhanced jigsaw, titlerule=0mm, left=2mm, opacityback=0, toprule=.15mm, arc=.35mm, opacitybacktitle=0.6, colbacktitle=quarto-callout-note-color!10!white, colframe=quarto-callout-note-color-frame, rightrule=.15mm, breakable, coltitle=black, bottomtitle=1mm, toptitle=1mm, title=\textcolor{quarto-callout-note-color}{\faInfo}\hspace{0.5em}{Versões do \texttt{morph.py}}, colback=white, bottomrule=.15mm, leftrule=.75mm]

A biblioteca \texttt{morph.py} possui duas versões públicas:

\begin{itemize}
\item
  \textbf{Versão 1.0} --- versão original, publicada junto ao artigo
  apresentado no EduComp2024: \url{https://github.com/fzampirolli/morph}
\item
  \textbf{Versão 1.1} --- versão compacta, utilizada nas atividades
  deste livro:
  \url{https://github.com/fzampirolli/pdi-vc/blob/master/morph/morph.py}
\end{itemize}

A versão 1.1 foi necessária para uso nas atividades do
\textbf{Moodle/VPL} (Laboratório Virtual de Programação). Na versão 1.0,
bibliotecas como \texttt{matplotlib}, \texttt{requests} e
\texttt{skimage} são carregadas no momento do \texttt{import\ morph}, o
que causava erro de memória (\texttt{Jail:\ out\ of\ memory,\ 128MiB})
no ambiente restrito do VPL.

Na versão 1.1, esses imports foram convertidos para \textbf{carregamento
lazy}: cada biblioteca é importada apenas dentro do método que a
utiliza, e somente quando esse método é de fato chamado. Assim, um
simples \texttt{import\ morph} carrega apenas \texttt{numpy} e
\texttt{cv2}, que são suficientes para a maioria dos exercícios
práticos.

\end{tcolorbox}

\begin{Shaded}
\begin{Highlighting}[]
\ImportTok{import}\NormalTok{ os, importlib, urllib.request}

\CommentTok{\# URLs do repositório}
\NormalTok{BASE\_URL }\OperatorTok{=} \StringTok{"https://raw.githubusercontent.com/fzampirolli/pdi{-}vc/master/morph"}
\ControlFlowTok{for}\NormalTok{ f }\KeywordTok{in}\NormalTok{ [}\StringTok{"morph.py"}\NormalTok{]:}
    \ControlFlowTok{if} \KeywordTok{not}\NormalTok{ os.path.exists(f):}
\NormalTok{        urllib.request.urlretrieve(}\SpecialStringTok{f"}\SpecialCharTok{\{}\NormalTok{BASE\_URL}\SpecialCharTok{\}}\SpecialStringTok{/}\SpecialCharTok{\{}\NormalTok{f}\SpecialCharTok{\}}\SpecialStringTok{"}\NormalTok{, f)}

\ImportTok{import}\NormalTok{ morph}
\NormalTok{importlib.}\BuiltInTok{reload}\NormalTok{(morph)}
\ImportTok{from}\NormalTok{ morph }\ImportTok{import}\NormalTok{ mm}

\CommentTok{\# Correção: Verifica se o atributo existe antes de imprimir, }
\CommentTok{\# ou simplesmente confirma o carregamento do módulo.}
\NormalTok{version }\OperatorTok{=} \BuiltInTok{getattr}\NormalTok{(morph, }\StringTok{"\_\_version\_\_"}\NormalTok{, }\StringTok{"local\_file"}\NormalTok{)}
\BuiltInTok{print}\NormalTok{(}\SpecialStringTok{f"✅ Ambiente pronto. Módulo \textquotesingle{}morph\textquotesingle{} carregado com sucesso."}\NormalTok{)}
\BuiltInTok{print}\NormalTok{(}\SpecialStringTok{f"📦 Funções disponíveis no mm: }\SpecialCharTok{\{}\BuiltInTok{dir}\NormalTok{(mm)[}\OperatorTok{{-}}\DecValTok{5}\NormalTok{:]}\SpecialCharTok{\}}\SpecialStringTok{..."}\NormalTok{) }\CommentTok{\# Mostra as últimas funções para confirmar}
\end{Highlighting}
\end{Shaded}

\begin{verbatim}
✅ Ambiente pronto. Módulo 'morph' carregado com sucesso.
📦 Funções disponíveis no mm: ['union', 'verifyBoundBox', 'water0', 'waterB', 'watershed']...
\end{verbatim}

\section{Fundamentos de Matrizes
(Cuidado!)}\label{fundamentos-de-matrizes-cuidado}

Como uma imagem digital é uma matriz, precisamos saber criá-las
corretamente. Em Python, existe uma armadilha comum ao usar o operador *
em listas:

\begin{tcolorbox}[enhanced jigsaw, titlerule=0mm, left=2mm, opacityback=0, toprule=.15mm, arc=.35mm, opacitybacktitle=0.6, colbacktitle=quarto-callout-warning-color!10!white, colframe=quarto-callout-warning-color-frame, rightrule=.15mm, breakable, coltitle=black, bottomtitle=1mm, toptitle=1mm, title=\textcolor{quarto-callout-warning-color}{\faExclamationTriangle}\hspace{0.5em}{Atenção à cópia de referências}, colback=white, bottomrule=.15mm, leftrule=.75mm]

Ao fazer \texttt{m\ =\ {[}{[}0{]}*2{]}*3}, você não cria 3 linhas
independentes, mas sim 3 referências para a \textbf{mesma} linha.
Alterar um valor em uma linha alterará todas as outras!

\end{tcolorbox}

Para visualizar esse comportamento, você pode testar o código no
\href{https://pythontutor.com/}{Python Tutor} e comparar com a forma
correta de criar matriz com listas:
\texttt{m\ =\ {[}{[}0{]}*2\ for\ \_\ in\ range(3){]}}.

A forma recomendada e padrão em Processamento Digital de Imagens é usar
\textbf{NumPy}, que cria matrizes com dados independentes e oferece
eficiência computacional. O código a seguir mostra diferentes formas de
criação de imagens sintéticas --- o resultado é exibido na
Figura~\ref{fig-imagens-sinteticas}.

\begin{Shaded}
\begin{Highlighting}[]
\CommentTok{\# Criando uma imagem preta (zeros) de 5x5 pixels}
\NormalTok{img\_preta }\OperatorTok{=}\NormalTok{ np.zeros((}\DecValTok{5}\NormalTok{, }\DecValTok{5}\NormalTok{), dtype}\OperatorTok{=}\StringTok{\textquotesingle{}uint8\textquotesingle{}}\NormalTok{)}

\CommentTok{\# Criando uma imagem branca (255) de 5x5 pixels}
\NormalTok{img\_branca }\OperatorTok{=}\NormalTok{ np.ones((}\DecValTok{5}\NormalTok{, }\DecValTok{5}\NormalTok{), dtype}\OperatorTok{=}\StringTok{\textquotesingle{}uint8\textquotesingle{}}\NormalTok{) }\OperatorTok{*} \DecValTok{255}

\CommentTok{\# Criando uma imagem aleatória para testes (Ruído)}
\NormalTok{img\_random }\OperatorTok{=}\NormalTok{ mm.randomImage(}\DecValTok{4}\NormalTok{, }\DecValTok{6}\NormalTok{, maxValue}\OperatorTok{=}\DecValTok{255}\NormalTok{)}

\BuiltInTok{print}\NormalTok{(}\StringTok{"Matriz Aleatória Gerada:"}\NormalTok{)}
\BuiltInTok{print}\NormalTok{(mm.drawImage(img\_random))}

\NormalTok{mm.show(img\_random, title}\OperatorTok{=}\StringTok{"Exemplo: Captura RGB"}\NormalTok{)}
\end{Highlighting}
\end{Shaded}

\begin{verbatim}
Matriz Aleatória Gerada:
138 151  52 128  13 137 
189  69 156 106  77 223 
194 165  49  82 178 226 
217 173 220 164 224 191 
\end{verbatim}

\begin{figure}[H]

\centering{

\pandocbounded{\includegraphics[keepaspectratio]{cap01/cap01.py.pt_files/figure-pdf/fig-imagens-sinteticas-output-2.png}}

}

\caption{\label{fig-imagens-sinteticas}Exemplo de imagem sintética
gerada por matriz aleatória 4x6.}

\end{figure}%

\section{Lendo e Exibindo Imagens}\label{lendo-e-exibindo-imagens}

Em Python, as imagens são lidas como \emph{arrays} do NumPy. Isso
significa que toda a potência da álgebra linear está disponível para o
processamento de imagens.

A operação mais básica em PDI é a leitura de uma imagem. A função
\texttt{mm.read()} da \texttt{morph.py} aceita caminhos locais e URLs,
ver o resultado na Figura~\ref{fig-1-lena}.

\begin{Shaded}
\begin{Highlighting}[]
\CommentTok{\# Carregando uma imagem de exemplo via URL}
\NormalTok{url }\OperatorTok{=} \StringTok{"https://upload.wikimedia.org/wikipedia/en/7/7d/Lenna\_\%28test\_image\%29.png"}
\NormalTok{img }\OperatorTok{=}\NormalTok{ mm.read(url)}

\CommentTok{\# Exibindo informações básicas}
\BuiltInTok{print}\NormalTok{(}\SpecialStringTok{f"Dimensões (H, W, Canais): }\SpecialCharTok{\{}\NormalTok{img}\SpecialCharTok{.}\NormalTok{shape}\SpecialCharTok{\}}\SpecialStringTok{"}\NormalTok{)}
\BuiltInTok{print}\NormalTok{(}\SpecialStringTok{f"Tipo de dado: }\SpecialCharTok{\{}\NormalTok{img}\SpecialCharTok{.}\NormalTok{dtype}\SpecialCharTok{\}}\SpecialStringTok{"}\NormalTok{)}

\CommentTok{\# Exibição estética}
\NormalTok{mm.show(img, title}\OperatorTok{=}\StringTok{"Exemplo: Captura RGB"}\NormalTok{)}
\end{Highlighting}
\end{Shaded}

\begin{verbatim}
Dimensões (H, W, Canais): (512, 512, 3)
Tipo de dado: uint8
\end{verbatim}

\begin{figure}[H]

\centering{

\pandocbounded{\includegraphics[keepaspectratio]{cap01/cap01.py.pt_files/figure-pdf/fig-1-lena-output-2.png}}

}

\caption{\label{fig-1-lena}Imagem clássica de processamento de imagens:
`Lena'.}

\end{figure}%

\subsection*{Alternativa: Baixando a imagem
localmente}\label{alternativa-baixando-a-imagem-localmente}
\addcontentsline{toc}{subsection}{Alternativa: Baixando a imagem
localmente}

Se por algum motivo a leitura direta da URL não funcionar (firewall,
restrições de rede, ou se você preferir trabalhar com arquivos locais),
uma abordagem simples é \textbf{baixar a imagem} usando o comando
\texttt{wget} (disponível em ambientes Linux, macOS e no Google Colab) e
depois carregá-la com \texttt{mm.read()} a partir do arquivo salvo.

\begin{Shaded}
\begin{Highlighting}[]
\ExtensionTok{!wget} \AttributeTok{{-}O}\NormalTok{ lena.png https://upload.wikimedia.org/wikipedia/en/7/7d/Lenna\_\%28test\_image\%29.png}
\end{Highlighting}
\end{Shaded}

\begin{Shaded}
\begin{Highlighting}[]
\NormalTok{img }\OperatorTok{=}\NormalTok{ mm.read(}\StringTok{\textquotesingle{}lena.png\textquotesingle{}}\NormalTok{)}
\NormalTok{mm.show(img, title}\OperatorTok{=}\StringTok{"Exemplo: Captura RGB (arquivo local)"}\NormalTok{)}
\end{Highlighting}
\end{Shaded}

\subsubsection*{Explicação:}\label{explicauxe7uxe3o}
\addcontentsline{toc}{subsubsection}{Explicação:}

\begin{itemize}
\tightlist
\item
  \texttt{!wget\ -O\ lena.png\ \textless{}URL\textgreater{}} baixa a
  imagem da Wikimedia e a salva com o nome \texttt{lena.png}.
\item
  \texttt{mm.read(\textquotesingle{}lena.png\textquotesingle{})} lê o
  arquivo local (sem necessidade de cabeçalhos HTTP ou tratamento
  especial).
\item
  Essa estratégia é útil quando você quer \textbf{evitar dependências de
  rede} durante a execução ou \textbf{reutilizar a mesma imagem} várias
  vezes sem novo download.
\end{itemize}

\begin{tcolorbox}[enhanced jigsaw, titlerule=0mm, left=2mm, opacityback=0, toprule=.15mm, arc=.35mm, opacitybacktitle=0.6, colbacktitle=quarto-callout-tip-color!10!white, colframe=quarto-callout-tip-color-frame, rightrule=.15mm, breakable, coltitle=black, bottomtitle=1mm, toptitle=1mm, title=\textcolor{quarto-callout-tip-color}{\faLightbulb}\hspace{0.5em}{Nota}, colback=white, bottomrule=.15mm, leftrule=.75mm]

O prefixo \texttt{!} funciona no Jupyter Notebook, JupyterLab e Google
Colab. Em um terminal comum, basta remover o \texttt{!}. Se o
\texttt{wget} não estiver instalado, instale-o com
\texttt{apt\ install\ wget} (Linux) ou use
\texttt{curl\ -o\ lena.png\ \textless{}URL\textgreater{}} como
alternativa.

\end{tcolorbox}

\section{Conversão de Tipos e
Limiarização}\label{conversuxe3o-de-tipos-e-limiarizauxe7uxe3o}

Conforme vimos, uma imagem colorida no espaço RGB é representada pela
função:

\[
f: \mathbb{E} \to \{0,1,\dots,255\}^3
\]

Ou seja, para cada pixel \((x,y)\), temos três valores \((R,G,B)\) que
definem sua cor.

\subsection*{Conversão para Tons de
Cinza}\label{conversuxe3o-para-tons-de-cinza}
\addcontentsline{toc}{subsection}{Conversão para Tons de Cinza}

Para converter uma imagem RGB em tons de cinza (\emph{grayscale}), é
necessário combinar os três canais em um único valor de intensidade
\(g\), que representa o brilho percebido. Como o olho humano não é
igualmente sensível ao vermelho, verde e azul, utiliza-se uma
\textbf{média ponderada}. O padrão
\href{https://www.itu.int/rec/R-REC-BT.601/}{ITU-R BT.601} (ITU-R, 2011)
define os seguintes pesos:

\begin{equation}\phantomsection\label{eq-1-cinza}{
g = 0.299\,R + 0.587\,G + 0.114\,B 
}\end{equation}

Após o cálculo, o valor \(g\) é arredondado para o inteiro mais próximo
e ajustado ao intervalo \([0, 255]\). O resultado é uma nova imagem,
agora em tons de cinza, representada por:

\[
f_{\text{cinza}}: \mathbb{E} \to \{0,1,\dots,255\}
\]

\subsection*{\texorpdfstring{Limiarização
(\emph{Thresholding})}{Limiarização (Thresholding)}}\label{limiarizauxe7uxe3o-thresholding}
\addcontentsline{toc}{subsection}{Limiarização (\emph{Thresholding})}

A partir da imagem em tons de cinza \(f_{\text{cinza}}(x,y)\), uma
operação fundamental é a \textbf{limiarização}, que produz uma imagem
\textbf{binária} (apenas preto e branco). Para isso, escolhe-se um valor
de corte \(T\) (geralmente no intervalo \([0,255]\)) e define-se:

\begin{equation}\phantomsection\label{eq-1-limiar}{
f_{\text{bin}}(x,y) = 
\begin{cases} 
255 & \text{se } f_{\text{cinza}}(x,y) > T \\[4pt]
0 & \text{caso contrário}
\end{cases}
}\end{equation}

\textbf{Exemplo:} Com \(T = 128\), pixels com intensidade acima de 128
tornam-se brancos (255); os demais tornam-se pretos (0).

A limiarização é amplamente usada para \textbf{segmentar objetos do
fundo}, extrair bordas ou criar máscaras binárias para processamento
posterior.

\begin{quote}
\textbf{Nota:} O valor 255 representa o branco máximo em imagens de 8
bits, enquanto 0 representa o preto absoluto.
\end{quote}

\subsection*{Exemplo prático de conversão e
limiarização}\label{exemplo-pruxe1tico-de-conversuxe3o-e-limiarizauxe7uxe3o}
\addcontentsline{toc}{subsection}{Exemplo prático de conversão e
limiarização}

A Figura~\ref{fig-1-processamento-basico} ilustra os principais passos
para transformar uma imagem colorida em tons de cinza e, em seguida,
binarizá‑la por limiarização. O código a seguir implementa essas etapas:

\begin{Shaded}
\begin{Highlighting}[]
\ImportTok{import}\NormalTok{ matplotlib.pyplot }\ImportTok{as}\NormalTok{ plt}

\CommentTok{\# 1. Converter para Tons de Cinza}
\NormalTok{img\_gray }\OperatorTok{=}\NormalTok{ mm.gray(img)}

\CommentTok{\# 2. Aplicar limiar (Pixels \textgreater{} 128 tornam{-}se 255, outros 0)}
\NormalTok{limiar }\OperatorTok{=} \DecValTok{128}
\NormalTok{img\_binaria }\OperatorTok{=}\NormalTok{ mm.threshold(img\_gray, limiar)}

\CommentTok{\# Uso da nova função}
\NormalTok{mm.show(}
\NormalTok{    [img, img\_gray, img\_binaria], }
\NormalTok{    titles}\OperatorTok{=}\NormalTok{[}\StringTok{"Original"}\NormalTok{, }\StringTok{"Tons de Cinza"}\NormalTok{, }\SpecialStringTok{f"Binária (T=}\SpecialCharTok{\{}\NormalTok{limiar}\SpecialCharTok{\}}\SpecialStringTok{)"}\NormalTok{],}
\NormalTok{    cols}\OperatorTok{=}\DecValTok{3}
\NormalTok{)}
\end{Highlighting}
\end{Shaded}

\begin{figure}[H]

\centering{

\pandocbounded{\includegraphics[keepaspectratio]{cap01/cap01.py.pt_files/figure-pdf/fig-1-processamento-basico-output-1.png}}

}

\caption{\label{fig-1-processamento-basico}Processamento básico de
imagens: (a) imagem original, (b) imagem em tons de cinza, (c) imagem
binarizada por limiar (T=128).}

\end{figure}%

\section{Limiarização pelo método de
Otsu}\label{limiarizauxe7uxe3o-pelo-muxe9todo-de-otsu}

Conforme apresentado na Equação~\ref{eq-1-limiar}, a limiarização
converte uma imagem em tons de cinza para binária usando um valor de
corte \(T\). Até agora fixamos \(T = 128\) manualmente.

\begin{Shaded}
\begin{Highlighting}[]
\NormalTok{img\_bin\_fixo }\OperatorTok{=}\NormalTok{ mm.threshold(img\_gray, T}\OperatorTok{=}\DecValTok{128}\NormalTok{)}
\end{Highlighting}
\end{Shaded}

Entretanto, a escolha manual de \(T\) nem sempre é trivial. A biblioteca
\texttt{mm} oferece uma alternativa automática: quando o parâmetro
\texttt{limiar} \textbf{não é fornecido}, a função
\texttt{mm.threshold(img\_gray)} calcula o valor de \(T\) pelo
\textbf{método de Otsu} (Otsu, 1979). Esse método, que será detalhado em
capítulos futuros, maximiza a variância inter‑classes do histograma,
separando automaticamente os pixels de objeto e fundo.

O código abaixo compara a limiarização manual (\(T=128\)) com a
automática (Otsu), mostrando também o valor de \(T\) calculado:

\begin{Shaded}
\begin{Highlighting}[]
\ImportTok{import}\NormalTok{ cv2}
\CommentTok{\# Limiarização com T fixo (manual)}
\NormalTok{T\_fixo }\OperatorTok{=} \DecValTok{128}
\NormalTok{img\_bin\_fixo }\OperatorTok{=}\NormalTok{ mm.threshold(img\_gray, T\_fixo)}

\CommentTok{\# Limiarização pelo método de Otsu (T automático)}
\NormalTok{T\_otsu, img\_bin\_otsu }\OperatorTok{=}\NormalTok{ cv2.threshold(img\_gray, }\DecValTok{0}\NormalTok{, }\DecValTok{255}\NormalTok{,}
\NormalTok{                                      cv2.THRESH\_BINARY }\OperatorTok{+}\NormalTok{ cv2.THRESH\_OTSU)}
\BuiltInTok{print}\NormalTok{(}\SpecialStringTok{f\textquotesingle{}Limiar calculado por Otsu: T = }\SpecialCharTok{\{}\NormalTok{T\_otsu}\SpecialCharTok{\}}\SpecialStringTok{\textquotesingle{}}\NormalTok{)}
\CommentTok{\# ou simplesmente:}
\CommentTok{\# img\_bin\_otsu = mm.threshold(img\_gray)}

\CommentTok{\# Exibição lado a lado}
\NormalTok{mm.show(}
\NormalTok{    [img\_gray, img\_bin\_fixo, img\_bin\_otsu], }
\NormalTok{    titles}\OperatorTok{=}\NormalTok{[}\StringTok{"Tons de Cinza"}\NormalTok{, }\SpecialStringTok{f"Binária (T=}\SpecialCharTok{\{}\NormalTok{T\_fixo}\SpecialCharTok{\}}\SpecialStringTok{)"}\NormalTok{, }\SpecialStringTok{f"Binária (Otsu, T=}\SpecialCharTok{\{}\NormalTok{T\_otsu}\SpecialCharTok{\}}\SpecialStringTok{)"}\NormalTok{],}
\NormalTok{    cols}\OperatorTok{=}\DecValTok{3}
\NormalTok{)}
\end{Highlighting}
\end{Shaded}

\begin{verbatim}
Limiar calculado por Otsu: T = 117.0
\end{verbatim}

\begin{figure}[H]

\centering{

\pandocbounded{\includegraphics[keepaspectratio]{cap01/cap01.py.pt_files/figure-pdf/fig-1-otsu-output-2.png}}

}

\caption{\label{fig-1-otsu}Comparação entre limiarização manual (T=128)
e automática (Otsu) sobre a imagem em tons de cinza.}

\end{figure}%

A Figura~\ref{fig-1-otsu} mostra que o limiar obtido por Otsu se adapta
automaticamente à imagem, resultando em uma binarização mais eficiente
do que um valor fixo, especialmente quando as intensidades do objeto e
do fundo são bem separadas no histograma. Essa técnica é amplamente
utilizada em sistemas de visão computacional para binarização de
documentos, detecção de objetos e pré‑processamento de imagens.

\begin{tcolorbox}[enhanced jigsaw, titlerule=0mm, left=2mm, opacityback=0, toprule=.15mm, arc=.35mm, opacitybacktitle=0.6, colbacktitle=quarto-callout-tip-color!10!white, colframe=quarto-callout-tip-color-frame, rightrule=.15mm, breakable, coltitle=black, bottomtitle=1mm, toptitle=1mm, title=\textcolor{quarto-callout-tip-color}{\faLightbulb}\hspace{0.5em}{Simplicidade da biblioteca \texttt{morph.py}}, colback=white, bottomrule=.15mm, leftrule=.75mm]

Enquanto o OpenCV exige a chamada completa:

\begin{Shaded}
\begin{Highlighting}[]
\NormalTok{T\_otsu, img\_bin }\OperatorTok{=}\NormalTok{ cv2.threshold(img\_gray, }\DecValTok{0}\NormalTok{, }\DecValTok{255}\NormalTok{,}
\NormalTok{                                cv2.THRESH\_BINARY }\OperatorTok{+}\NormalTok{ cv2.THRESH\_OTSU)}
\end{Highlighting}
\end{Shaded}

a biblioteca \texttt{mm} abstrai toda essa complexidade: basta chamar
\texttt{mm.threshold(img\_gray)}. O limiar de Otsu é calculado
automaticamente e a imagem binária é retornada diretamente. Essa
abordagem permite que você se concentre no conceito, não nos detalhes de
implementação.

\end{tcolorbox}

\section{Acesso a Pixels}\label{acesso-a-pixels}

Em Python com NumPy, uma imagem é representada como um \textbf{array
multidimensional}. Para acessar um pixel específico, utilizam-se as
coordenadas \textbf{linha} (eixo Y) e \textbf{coluna} (eixo X):
\texttt{img{[}linha,\ coluna{]}}.

O código abaixo demonstra como obter o valor de um pixel em uma imagem
RGB e em sua versão em tons de cinza, além de criar uma pequena imagem
sintética para visualizar a estrutura de uma matriz de pixels.

\begin{Shaded}
\begin{Highlighting}[]
\ImportTok{import}\NormalTok{ numpy }\ImportTok{as}\NormalTok{ np}
\ImportTok{import}\NormalTok{ matplotlib.pyplot }\ImportTok{as}\NormalTok{ plt}

\CommentTok{\# Coordenadas do pixel que queremos examinar}
\NormalTok{r, c }\OperatorTok{=} \DecValTok{100}\NormalTok{, }\DecValTok{100}

\CommentTok{\# Acessa o pixel na imagem em tons de cinza (um valor escalar)}
\NormalTok{pixel\_cinza }\OperatorTok{=}\NormalTok{ img\_gray[r, c]}

\CommentTok{\# Acessa o pixel na imagem colorida RGB (vetor de 3 valores)}
\NormalTok{pixel\_rgb }\OperatorTok{=}\NormalTok{ img[r, c]}

\BuiltInTok{print}\NormalTok{(}\SpecialStringTok{f\textquotesingle{}Pixel na posição (}\SpecialCharTok{\{}\NormalTok{r}\SpecialCharTok{\}}\SpecialStringTok{,}\SpecialCharTok{\{}\NormalTok{c}\SpecialCharTok{\}}\SpecialStringTok{):\textquotesingle{}}\NormalTok{)}
\BuiltInTok{print}\NormalTok{(}\SpecialStringTok{f\textquotesingle{}  {-} Tons de cinza : }\SpecialCharTok{\{}\NormalTok{pixel\_cinza}\SpecialCharTok{\}}\SpecialStringTok{\textquotesingle{}}\NormalTok{)}
\BuiltInTok{print}\NormalTok{(}\SpecialStringTok{f\textquotesingle{}  {-} RGB           : R=}\SpecialCharTok{\{}\NormalTok{pixel\_rgb[}\DecValTok{0}\NormalTok{]}\SpecialCharTok{\}}\SpecialStringTok{, G=}\SpecialCharTok{\{}\NormalTok{pixel\_rgb[}\DecValTok{1}\NormalTok{]}\SpecialCharTok{\}}\SpecialStringTok{, B=}\SpecialCharTok{\{}\NormalTok{pixel\_rgb[}\DecValTok{2}\NormalTok{]}\SpecialCharTok{\}}\SpecialStringTok{\textquotesingle{}}\NormalTok{)}

\CommentTok{\# Cria uma imagem sintética 5×5 com todos os pixels pretos (0)}
\NormalTok{syn }\OperatorTok{=}\NormalTok{ np.zeros((}\DecValTok{5}\NormalTok{, }\DecValTok{5}\NormalTok{), dtype}\OperatorTok{=}\NormalTok{np.uint8)}
\CommentTok{\# Torna o pixel central (linha 2, coluna 2) branco (255)}
\NormalTok{syn[}\DecValTok{2}\NormalTok{, }\DecValTok{2}\NormalTok{] }\OperatorTok{=} \DecValTok{255}

\BuiltInTok{print}\NormalTok{(}\StringTok{"}\CharTok{\textbackslash{}n}\StringTok{Matriz da imagem sintética 5×5:"}\NormalTok{)}
\BuiltInTok{print}\NormalTok{(syn)}

\CommentTok{\# Exibe a imagem sintética}
\CommentTok{\# plt.figure(figsize=(3, 3))}
\CommentTok{\# plt.imshow(syn, cmap=\textquotesingle{}gray\textquotesingle{}, vmin=0, vmax=255)}
\CommentTok{\# plt.title("Imagem sintética 5×5\textbackslash{}n(pixel central branco)")}
\CommentTok{\# plt.axis(\textquotesingle{}off\textquotesingle{})}
\CommentTok{\# plt.tight\_layout()}
\CommentTok{\# plt.show()}
\CommentTok{\# ou simplesmente:}

\NormalTok{mm.show(syn, title}\OperatorTok{=}\StringTok{"Exemplo: Imagem Sintética 5x5"}\NormalTok{)}
\end{Highlighting}
\end{Shaded}

\begin{verbatim}
Pixel na posição (100,100):
  - Tons de cinza : 102
  - RGB           : R=178, G=68, B=78

Matriz da imagem sintética 5×5:
[[  0   0   0   0   0]
 [  0   0   0   0   0]
 [  0   0 255   0   0]
 [  0   0   0   0   0]
 [  0   0   0   0   0]]
\end{verbatim}

\begin{figure}[H]

\centering{

\pandocbounded{\includegraphics[keepaspectratio]{cap01/cap01.py.pt_files/figure-pdf/fig-1-pixel-acesso-output-2.png}}

}

\caption{\label{fig-1-pixel-acesso}Exemplo de acesso a pixels: (a)
imagem original com pixel destacado, (b) imagem em tons de cinza, (c)
matriz sintética 5×5 com pixel central branco.}

\end{figure}%

\textbf{Explicação linha a linha:}

\begin{enumerate}
\def\labelenumi{\arabic{enumi}.}
\tightlist
\item
  \texttt{img\_gray{[}r,\ c{]}} -- retorna um único número inteiro entre
  0 e 255, correspondente ao nível de cinza naquela posição.\\
\item
  \texttt{img{[}r,\ c{]}} -- retorna uma tupla ou array com três valores
  (R, G, B).\\
\item
  \texttt{np.zeros((5,5),\ dtype=np.uint8)} -- cria uma matriz
  \(5\times 5\) preenchida com zeros (preto).\\
\item
  \texttt{syn{[}2,\ 2{]}\ =\ 255} -- altera o elemento central da matriz
  para 255 (branco), demonstrando como modificar um pixel.\\
\item
  A exibição com \texttt{plt.imshow} revela um pequeno quadrado com um
  ponto branco no meio, ilustrando visualmente a estrutura discreta da
  imagem.
\end{enumerate}

A Figura~\ref{fig-1-pixel-acesso} mostra os valores impressos e a imagem
sintética. Perceba que a indexação em Python é \textbf{zero‑based}: o
canto superior esquerdo é \texttt{(0,0)}. A ordem \textbf{linha ×
coluna} corresponde à estrutura matricial da imagem: primeira dimensão →
altura (Y), segunda dimensão → largura (X).

\section{Resumo}\label{resumo}

Neste capítulo foram apresentados os fundamentos de representação de
imagens digitais:

\begin{itemize}
\tightlist
\item
  \textbf{Imagem digital} = função \(f(x,y)\) que mapeia coordenadas
  para intensidades (escalares ou vetoriais).
\item
  \textbf{Domínio}: conjunto finito
  \(\mathbb{E} = \{(x,y) \in \mathbb{Z}^2 \mid 0 \le x < L,\; 0 \le y < H\}\).
\item
  \textbf{Tipos principais}: binária (\(\mathcal{V} = \{0, 255\}\)),
  tons de cinza (\(\mathcal{V} = [0,255]\)) e RGB
  (\(\mathcal{V} = [0,255]^3\)).
\item
  A biblioteca \texttt{morph.py} (ou \texttt{mm}) oferece funções
  didáticas para operações básicas de PDI, como \texttt{mm.gray()},
  \texttt{mm.threshold()}, \texttt{mm.show\_multiple()}.
\item
  \textbf{Armadilha NumPy}: nunca use \texttt{{[}{[}0{]}*n{]}*m} para
  criar matrizes --- use sempre \texttt{np.zeros()} ou
  \texttt{np.ones()}.
\item
  \textbf{Limiarização} converte tons de cinza em binária; o
  \textbf{método de Otsu} determina o valor de corte automaticamente
  maximizando a variância inter-classes.
\item
  Acesso a pixels via \texttt{img{[}linha,\ coluna{]}}, com indexação
  \textbf{zero-based}.
\end{itemize}

O Capítulo 2 abordará \textbf{histogramas e equalização de contraste}.

\section{\texorpdfstring{\emoji{robot} Uso do NotebookLM como Tutor
Complementar}{ Uso do NotebookLM como Tutor Complementar}}\label{uso-do-notebooklm-como-tutor-complementar}

Nesta edição, além dos notebooks interativos no Google Colab,
incentivamos o uso do \textbf{NotebookLM} como ferramenta complementar
de aprendizagem. Essa ferramenta de IA utiliza exclusivamente os
documentos fornecidos pelo autor como base de conhecimento, garantindo
respostas alinhadas ao conteúdo do livro.

Para cada capítulo, preparamos um projeto específico na plataforma. Para
uma experiência de estudo ampliada, utilize o acesso abaixo:

\begin{tcolorbox}[enhanced jigsaw, titlerule=0mm, left=2mm, opacityback=0, toprule=.15mm, arc=.35mm, opacitybacktitle=0.6, colbacktitle=quarto-callout-important-color!10!white, colframe=quarto-callout-important-color-frame, rightrule=.15mm, breakable, coltitle=black, bottomtitle=1mm, toptitle=1mm, title={\emoji{graduation-cap} Estude com o Tutor Inteligente}, colback=white, bottomrule=.15mm, leftrule=.75mm]

Para interagir com o conteúdo deste capítulo, acesse o link a seguir. O
ambiente contém materiais didáticos em diferentes formatos, gerados a
partir do \textbf{PDF} do capítulo. Na plataforma, explore especialmente
as opções \textbf{Guia de Estudo} e \textbf{Conversa} para aprofundar
sua compreensão.

\href{https://notebooklm.google.com/notebook/aca1138a-02aa-4f98-b777-1e25795ca635}{\emoji{rocket}
ACESSAR NOTEBOOKLM: CAPÍTULO 01}

\end{tcolorbox}

\subsection*{Funcionalidades Disponíveis na
Plataforma}\label{funcionalidades-disponuxedveis-na-plataforma}
\addcontentsline{toc}{subsection}{Funcionalidades Disponíveis na
Plataforma}

O \textbf{NotebookLM} oferece uma suíte avançada de ferramentas baseadas
em IA para transformar o conteúdo estático do livro em uma experiência
de aprendizado dinâmica e multimídia. Conforme ilustrado na
Figura~\ref{fig-notebooklm}, a plataforma utiliza técnicas de
\textbf{RAG} (\emph{Retrieval-Augmented Generation}), fundamentadas no
trabalho de Lewis et al. (2020), para basear as respostas estritamente
nos documentos fornecidos e minimizar a ocorrência de alucinações.

As principais funcionalidades incluem:

\begin{itemize}
\tightlist
\item
  \textbf{Resumos Multimodais (Áudio e Vídeo)}: Geração de conversas
  naturais entre especialistas no formato de \textbf{Resumo em Áudio}
  (estilo \emph{podcast}) e \textbf{Resumo em Vídeo}, discutindo os
  temas centrais do capítulo, como as diferenças entre processamento de
  imagens e visão computacional, ou a interpretação de transformações
  como a limiarização e o método de Otsu.
\item
  \textbf{Visualização de Estruturas (Mapa Mental e Infográfico)}:
  Criação automática de diagramas que conectam visualmente os conceitos,
  por exemplo, o fluxo de processamento desde a captura da imagem
  digital, passando pela conversão para tons de cinza, limiarização e
  segmentação binária.
\item
  \textbf{Ferramentas de Avaliação (Teste e Cartões Didáticos)}: Geração
  de \textbf{Testes} de múltipla escolha e \textbf{Cartões Didáticos}
  (\emph{flashcards}) para fixação de conhecimento, baseados no texto
  autoral (ex.: perguntas sobre a fórmula de conversão RGB→cinza ou
  sobre o funcionamento do limiar global e de Otsu).
\item
  \textbf{Apoio à Apresentação (\emph{Slides} e Relatórios)}: Auxílio na
  estruturação de \textbf{Apresentações de \emph{Slides}} e na redação
  de \textbf{Relatórios} técnicos e \textbf{\emph{Briefings}},
  facilitando a comunicação de resultados de experimentos com imagens.
\item
  \textbf{Análise de Dados (Tabela de Dados)}: Organização de dados
  extraídos do texto em tabelas estruturadas, auxiliando na compreensão
  de exemplos práticos, como a comparação entre diferentes valores de
  limiar.
\item
  \textbf{\emph{Chat} Contextualizado}: Permite o questionamento direto
  sobre o código e a teoria, como: \emph{``Como implementar a conversão
  de RGB para tons de cinza usando os pesos do padrão ITU‑R BT.601?''}
  ou \emph{``O que acontece com a imagem binária se eu escolher um
  limiar T=200 em vez de T=128?''}.
\end{itemize}

\begin{figure}

\centering{

\includegraphics[width=0.8\linewidth,height=\textheight,keepaspectratio]{cap01/imagens/fig-1-notebooklm.png}

}

\caption{\label{fig-notebooklm}Estúdio do NotebookLM.}

\end{figure}%

\section{Lista de Exercícios}\label{lista-de-exercuxedcios}

\begin{enumerate}
\def\labelenumi{\arabic{enumi}.}
\item
  (15\%) Com suas próprias palavras, defina \textbf{imagem digital} e
  \textbf{pixel}. Dê um exemplo concreto de como uma imagem colorida
  (RGB) é representada matricialmente no computador.
\item
  (15\%) Explique as diferenças entre \textbf{imagem binária},
  \textbf{tons de cinza (8 bits)} e \textbf{colorida RGB}, indicando a
  faixa de valores possíveis para cada pixel em cada tipo.
\item
  (20\%) Considerando a fórmula de conversão RGB → tons de cinza do
  padrão ITU‑R BT.601: \[g = 0.299\,R + 0.587\,G + 0.114\,B\] Calcule o
  valor do pixel em tons de cinza para \((R,G,B) = (80, 180, 30)\).
  Arredonde para o inteiro mais próximo.
\item
  (20\%) O que é \textbf{limiarização} (\emph{thresholding})? Explique a
  diferença entre escolher um limiar \(T\) fixo (ex.: \(T=128\)) e
  utilizar o \textbf{método de Otsu} para determinação automática do
  limiar. Em poucas palavras, como o método de Otsu escolhe o limiar?
\item
  (15\%) No contexto da biblioteca didática \texttt{mm} discutida no
  capítulo, responda:

  \begin{itemize}
  \tightlist
  \item
    \begin{enumerate}
    \def\labelenumii{\alph{enumii})}
    \tightlist
    \item
      (7,5\%) Como se acessa o valor do pixel na posição (linha=50,
      coluna=60) de uma imagem em tons de cinza \texttt{img\_gray}?
    \end{enumerate}
  \item
    \begin{enumerate}
    \def\labelenumii{\alph{enumii})}
    \setcounter{enumii}{1}
    \tightlist
    \item
      (7,5\%) Qual a vantagem de usar \texttt{mm.threshold(img\_gray)}
      sem passar o limiar? Compare com a chamada equivalente no OpenCV.
    \end{enumerate}
  \end{itemize}
\item
  (15\%) O que a propriedade \texttt{img.shape} retorna para uma imagem
  NumPy no formato RGB? Dê um exemplo concreto com uma imagem de 640×480
  pixels.
\end{enumerate}

\section{Referências do Capítulo}\label{referuxeancias-do-capuxedtulo}

Este capítulo baseou‑se principalmente nas seguintes obras de
processamento de imagens e visão computacional:

\begin{itemize}
\tightlist
\item
  Gonzalez; Woods (2018) para fundamentos de \textbf{Processamento
  Digital de Imagens} (PDI).
\item
  Russ; Russ (2023) para aspectos práticos de \textbf{processamento e
  análise de imagens} (The Image Processing Handbook, 7ª ed., 2016).
\item
  Szeliski (2022) para \textbf{visão computacional} e algoritmos
  fundamentais (Computer Vision: Algorithms and Applications, 2ª ed.,
  2022).
\item
  Bradski; Kaehler (2016) para a biblioteca \textbf{OpenCV} e exemplos
  práticos em Python (Learning OpenCV 4, 2ª ed., 2016).
\item
  Lewis et al. (2020) para o conceito de \textbf{geração aumentada por
  recuperação (RAG)}, aplicado ao NotebookLM.
\end{itemize}

\begin{center}\rule{0.5\linewidth}{0.5pt}\end{center}

\href{https://colab.research.google.com/github/fzampirolli/pdi-vc/blob/master/notebooks_alunos/cap01/cap01.EPs_aluno.ipynb}{
\includegraphics[width=0.16\linewidth,height=\textheight,keepaspectratio]{cap01/imagens/colab-badge.png}}
\href{https://github.com/fzampirolli/pdi-vc}{
\includegraphics[width=0.2\linewidth,height=\textheight,keepaspectratio]{cap01/imagens/github-badge.png}}

\section{\texorpdfstring{\emoji{laptop} \textbf{Parte Prática com
Exercícios de
Programação}}{ Parte Prática com Exercícios de Programação}}\label{parte-pruxe1tica-com-exercuxedcios-de-programauxe7uxe3o}

\begin{center}\rule{0.5\linewidth}{0.5pt}\end{center}

\subsection*{\texorpdfstring{\emoji{bullseye} Objetivo deste
Caderno}{ Objetivo deste Caderno}}\label{objetivo-deste-caderno}
\addcontentsline{toc}{subsection}{\emoji{bullseye} Objetivo deste
Caderno}

Os \textbf{Exercícios de Programação (EPs)} apresentados a seguir podem
ser submetidos também em atividades do Moodle (atividades VPL) que
fornecem \emph{feedback} automático.

\textbf{Este caderno foi desenvolvido para superar limitações de uso do
Moodle.} Com ele, você pode:

\begin{enumerate}
\def\labelenumi{\arabic{enumi}.}
\tightlist
\item
  \textbf{Desenvolver:} Escrever e editar sua solução diretamente no
  ambiente Colab.
\item
  \textbf{Validar:} Testar seu código localmente utilizando os
  \textbf{mesmos casos de teste} do Moodle.
\item
  \textbf{Organizar:} Salvar seus códigos das atividades VPL de forma
  segura.
\item
  \textbf{Avaliar:} Quando estiver conectado ao Moodle, basta copiar sua
  solução e clicar em \textbf{Avaliar} no Moodle para registrar sua nota
  oficial.
\end{enumerate}

\begin{center}\rule{0.5\linewidth}{0.5pt}\end{center}

\subsection*{\texorpdfstring{\S{} Instruções Passo a
Passo}{ Instruções Passo a Passo}}\label{instruuxe7uxf5es-passo-a-passo}
\addcontentsline{toc}{subsection}{\S{} Instruções Passo a Passo}

Em um ambiente de execução (como VSCode, Jupyter ou Colab), siga a ordem
abaixo para configurar o ambiente e validar seus exercícios:

\subsubsection*{Preparação do
Ambiente}\label{preparauxe7uxe3o-do-ambiente}
\addcontentsline{toc}{subsubsection}{Preparação do Ambiente}

Execute a célula de código abaixo para baixar \texttt{morph.py} e
\texttt{testsuite.py} do repositório da disciplina --- apenas se ainda
não existirem no diretório local. Com ambos em \texttt{./}, o notebook e
os subprocessos do \texttt{TestSuite} encontram o módulo sem
configurações extras de caminho.

\emph{Nota:} O script \texttt{testsuite.py} buscará automaticamente os
casos de teste em \texttt{all/\{cap\}/cases} no
\href{https://github.com/fzampirolli/pdi-vc/tree/master/all}{GitHub}.

\subsubsection*{Escrevendo o Código}\label{escrevendo-o-cuxf3digo}
\addcontentsline{toc}{subsubsection}{Escrevendo o Código}

Salve sua solução em uma célula de código usando o comando mágico
\texttt{\%\%writefile}. O nome do arquivo deve seguir o padrão
\texttt{EPX\_Y.*}, onde \texttt{X} é o capítulo, \texttt{Y} é o
exercício e \texttt{*} é a extensão da linguagem.

\begin{quote}
\textbf{Exemplo:} \texttt{\%\%writefile\ EP01\_01.py}
\end{quote}

\subsubsection*{Download}\label{download}
\addcontentsline{toc}{subsubsection}{Download}

Baixe \texttt{morph.py} e \texttt{testsuite.py} executando a célula
abaixo:

\begin{Shaded}
\begin{Highlighting}[]
\ImportTok{import}\NormalTok{ os, sys, importlib, inspect, urllib.request}

\CommentTok{\# URLs do repositório}
\NormalTok{BASE\_URL }\OperatorTok{=} \StringTok{"https://raw.githubusercontent.com/fzampirolli/pdi{-}vc/master/morph"}
\ControlFlowTok{for}\NormalTok{ f }\KeywordTok{in}\NormalTok{ [}\StringTok{"morph.py"}\NormalTok{, }\StringTok{"testsuite.py"}\NormalTok{]:}
    \ControlFlowTok{if} \KeywordTok{not}\NormalTok{ os.path.exists(f):}
\NormalTok{        urllib.request.urlretrieve(}\SpecialStringTok{f"}\SpecialCharTok{\{}\NormalTok{BASE\_URL}\SpecialCharTok{\}}\SpecialStringTok{/}\SpecialCharTok{\{}\NormalTok{f}\SpecialCharTok{\}}\SpecialStringTok{"}\NormalTok{, f)}

\ImportTok{import}\NormalTok{ morph, testsuite}
\NormalTok{importlib.}\BuiltInTok{reload}\NormalTok{(morph)}\OperatorTok{;}\NormalTok{ importlib.}\BuiltInTok{reload}\NormalTok{(testsuite)}
\ImportTok{from}\NormalTok{ morph }\ImportTok{import}\NormalTok{ mm}
\ImportTok{from}\NormalTok{ testsuite }\ImportTok{import}\NormalTok{ TestSuite}

\BuiltInTok{print}\NormalTok{(}\SpecialStringTok{f"✅ Ambiente pronto. Morph: }\SpecialCharTok{\{}\NormalTok{morph}\SpecialCharTok{.}\NormalTok{\_\_version\_\_}\SpecialCharTok{\}}\SpecialStringTok{ | TestSuite: }\SpecialCharTok{\{}\NormalTok{testsuite}\SpecialCharTok{.}\NormalTok{\_\_version\_\_}\SpecialCharTok{\}}\SpecialStringTok{"}\NormalTok{)}
\end{Highlighting}
\end{Shaded}

\begin{verbatim}
✅ Ambiente pronto. Morph: 1.1.0 | TestSuite: 1.1.0
\end{verbatim}

\begin{center}\rule{0.5\linewidth}{0.5pt}\end{center}

\subsubsection*{Executando os Testes}\label{executando-os-testes}
\addcontentsline{toc}{subsubsection}{Executando os Testes}

Após salvar o arquivo com sua solução, execute o comando abaixo (em uma
nova célula) para avaliar os testes automáticos:

\begin{Shaded}
\begin{Highlighting}[]
\NormalTok{TestSuite(}\StringTok{"EP01\_01.extensão"}\NormalTok{).run()}
\end{Highlighting}
\end{Shaded}

Substitua a extensão conforme a linguagem usada:

\begin{longtable}[]{@{}ll@{}}
\toprule\noalign{}
Linguagem & Extensão \\
\midrule\noalign{}
\endhead
\bottomrule\noalign{}
\endlastfoot
Python & \texttt{.py} \\
Java & \texttt{.java} \\
C & \texttt{.c} \\
C++ & \texttt{.cpp} \\
JavaScript & \texttt{.js} \\
R & \texttt{.r} \\
\end{longtable}

\begin{quote}
\textbf{Como funciona:} O \texttt{TestSuite} baixa os casos de teste do
GitHub, executa seu programa com cada entrada e compara a saída com a
esperada -- calculando automaticamente sua nota.
\end{quote}

\subsection*{\texorpdfstring{\textbf{[!]} Importante: Regras e Boas
Práticas}{ Importante: Regras e Boas Práticas}}\label{importante-regras-e-boas-pruxe1ticas}
\addcontentsline{toc}{subsection}{\textbf{[!]} Importante: Regras e Boas
Práticas}

\subsubsection*{\texorpdfstring{\emoji{small-blue-diamond} Sobre a
Entrada de
Dados}{ Sobre a Entrada de Dados}}\label{sobre-a-entrada-de-dados}
\addcontentsline{toc}{subsubsection}{\emoji{small-blue-diamond} Sobre a
Entrada de Dados}

O seu programa deve ler a entrada padrão (teclado).

\begin{itemize}
\tightlist
\item
  \textbf{Python:} Utilize \texttt{input()}.
\item
  \textbf{Outras linguagens:} Utilize o comando de leitura padrão
  equivalente (\texttt{cin}, \texttt{Scanner}, etc.).
\end{itemize}

\subsubsection*{\texorpdfstring{\emoji{small-blue-diamond} Configuração
de IA no
Colab}{ Configuração de IA no Colab}}\label{configurauxe7uxe3o-de-ia-no-colab}
\addcontentsline{toc}{subsubsection}{\emoji{small-blue-diamond}
Configuração de IA no Colab}

Para um melhor aprendizado, recomenda-se desativar o preenchimento
automático de código por IA, pois não estará disponível durante as
avaliações. Por exemplo no navegador Chrome:

\begin{itemize}
\tightlist
\item
  Vá em: \textbf{Ferramentas \textgreater{} Configurações \textgreater{}
  IA generativa}
\item
  Desmarque: \textbf{Habilitar geração de código}
\end{itemize}

\subsubsection*{\texorpdfstring{\emoji{small-blue-diamond} Integridade
Acadêmica
(Plágio)}{ Integridade Acadêmica (Plágio)}}\label{integridade-acaduxeamica-pluxe1gio}
\addcontentsline{toc}{subsubsection}{\emoji{small-blue-diamond}
Integridade Acadêmica (Plágio)}

Este recurso de testes locais aplica-se a EPs \textbf{sem variações}. No
entanto:

\begin{itemize}
\tightlist
\item
  \textbf{Individualidade:} Cada aluno deve desenvolver sua própria
  solução.
\item
  \textbf{Detecção de Similaridade:} O professor utiliza ferramentas que
  detectam cópias, \textbf{mesmo com alteração de nomes de variáveis ou
  espaços em branco}.
\end{itemize}

\begin{center}\rule{0.5\linewidth}{0.5pt}\end{center}

\subsection{\texorpdfstring{EP01\_01 \emoji{straight-ruler} Três
métricas de distância em
PDI}{EP01\_01  Três métricas de distância em PDI}}\label{ep01_01-truxeas-muxe9tricas-de-distuxe2ncia-em-pdi}

Nesta atividade, você deve escrever um programa que calcule as três
distâncias mais usadas em PDI: \textbf{Euclidiana (L2)},
\textbf{City‑Block (L1)} e \textbf{Chessboard (L$\infty$)}.

\begin{itemize}
\tightlist
\item
  Leia \textbf{4 números reais} que representam as coordenadas:
  \(A_x, A_y, B_x, B_y\).
\item
  Calcule as três distâncias utilizando as fórmulas:
\end{itemize}

\[d_{\text{Euclidiana}} = \sqrt{(B_x - A_x)^2 + (B_y - A_y)^2}\]

\[d_{\text{City-block}} = |B_x - A_x| + |B_y - A_y|\]

\[d_{\text{Chessboard}} = \max\big(|B_x - A_x|,\; |B_y - A_y|\big)\]

\begin{itemize}
\tightlist
\item
  Imprima os três resultados, cada um em uma linha, formatados com
  \textbf{duas casas decimais}, na ordem: \textbf{Euclidiana},
  \textbf{City‑block}, \textbf{Chessboard}.
\end{itemize}

\emoji{pushpin} \textbf{Importante}:

\begin{itemize}
\tightlist
\item
  Utilize as funções matemáticas padrão da sua linguagem:
  \texttt{math.sqrt}, \texttt{abs} (ou \texttt{fabs}) e \texttt{max}.
\item
  A saída deve conter \textbf{apenas os números} (um por linha), sem
  textos adicionais.
\item
  Ver um simulador interativo para esta questão na
  \textbf{Figura~\ref{fig-01-metricas-distancia}} (gráfico com arrasto
  dos pontos e visualização das três métricas).
\end{itemize}

\begin{center}\rule{0.5\linewidth}{0.5pt}\end{center}

\subsubsection{\texorpdfstring{\emoji{framed-picture} Por que isso
importa? -- Custo
computacional}{ Por que isso importa? -- Custo computacional}}\label{por-que-isso-importa-custo-computacional}

Em uma imagem \textbf{1000×1000 pixels} (1 milhão de pixels), calcular a
distância de cada pixel a um ponto de referência exige \textbf{1 milhão
de operações}. A escolha da métrica afeta o desempenho:

\begin{longtable}[]{@{}
  >{\raggedright\arraybackslash}p{(\linewidth - 6\tabcolsep) * \real{0.2500}}
  >{\raggedright\arraybackslash}p{(\linewidth - 6\tabcolsep) * \real{0.2500}}
  >{\raggedright\arraybackslash}p{(\linewidth - 6\tabcolsep) * \real{0.2500}}
  >{\raggedright\arraybackslash}p{(\linewidth - 6\tabcolsep) * \real{0.2500}}@{}}
\toprule\noalign{}
\begin{minipage}[b]{\linewidth}\raggedright
Métrica
\end{minipage} & \begin{minipage}[b]{\linewidth}\raggedright
Operações por pixel
\end{minipage} & \begin{minipage}[b]{\linewidth}\raggedright
Custo relativo (1M pixels)
\end{minipage} & \begin{minipage}[b]{\linewidth}\raggedright
Quando usar
\end{minipage} \\
\midrule\noalign{}
\endhead
\bottomrule\noalign{}
\endlastfoot
\textbf{Euclidiana (L2)} & 2 subtrações, 2 multiplicações, 1 soma, 1
\texttt{sqrt} & \emoji{red-circle} Mais custosa -- \texttt{sqrt} é cara
& Distância ``real'' no espaço contínuo \\
\textbf{City‑block (L1)} & 2 subtrações, 2 \texttt{abs}, 1 soma &
\emoji{yellow-circle} Moderada -- sem raiz quadrada & Grids, robótica,
imagens binárias \\
\textbf{Chessboard (L$\infty$)} & 2 subtrações, 2 \texttt{abs}, 1
\texttt{max} & \emoji{green-circle} Mais eficiente & Movimentos de
peças, morfologia \\
\end{longtable}

A função \texttt{sqrt} é computacionalmente mais cara que operações como
adição, subtração, multiplicação e valor absoluto. Em CPUs modernas, a
diferença pode ser pequena (cerca de 1,5× a 3×), mas em sistemas
embarcados ou em laços de milhões de iterações, qualquer ganho importa.
Por isso, \textbf{quando o objetivo é apenas comparar distâncias} (ex.:
encontrar o ponto mais próximo), use a \textbf{distância euclidiana ao
quadrado}.

\begin{center}\rule{0.5\linewidth}{0.5pt}\end{center}

\subsubsection{\texorpdfstring{\emoji{clipboard} Tarefa (especificação
para
VPL)}{ Tarefa (especificação para VPL)}}\label{tarefa-especificauxe7uxe3o-para-vpl}

\textbf{Entrada:}\\
Uma única linha com quatro números reais: \texttt{Ax\ Ay\ Bx\ By}

\textbf{Saída:}\\
Três linhas, cada uma com um número real de duas casas decimais
(Euclidiana, City‑block, Chessboard).

\begin{center}\rule{0.5\linewidth}{0.5pt}\end{center}

\subsubsection{\texorpdfstring{\emoji{pushpin}
Exemplos}{ Exemplos}}\label{exemplos}

\begin{longtable}[]{@{}
  >{\raggedright\arraybackslash}p{(\linewidth - 4\tabcolsep) * \real{0.3333}}
  >{\raggedright\arraybackslash}p{(\linewidth - 4\tabcolsep) * \real{0.3333}}
  >{\raggedright\arraybackslash}p{(\linewidth - 4\tabcolsep) * \real{0.3333}}@{}}
\toprule\noalign{}
\begin{minipage}[b]{\linewidth}\raggedright
Entrada
\end{minipage} & \begin{minipage}[b]{\linewidth}\raggedright
Saída
\end{minipage} & \begin{minipage}[b]{\linewidth}\raggedright
Observação
\end{minipage} \\
\midrule\noalign{}
\endhead
\bottomrule\noalign{}
\endlastfoot
\texttt{0}\texttt{0}\texttt{3}\texttt{4} &
\texttt{5.00}\texttt{7.00}\texttt{4.00} & Triângulo 3‑4‑5 \\
\texttt{0}\texttt{0}\texttt{1}\texttt{1} &
\texttt{1.41}\texttt{2.00}\texttt{1.00} & Diagonal unitária \\
\end{longtable}

\begin{center}\rule{0.5\linewidth}{0.5pt}\end{center}

\textbf{Exemplo de teste de \texttt{sqrt} em Python, com \texttt{timeit}
isolando cada operação:}

\begin{Shaded}
\begin{Highlighting}[]
\ImportTok{import}\NormalTok{ math}
\ImportTok{import}\NormalTok{ timeit}

\NormalTok{N }\OperatorTok{=} \DecValTok{50\_000\_000}

\KeywordTok{def}\NormalTok{ apenas\_soma():}
\NormalTok{    a, b }\OperatorTok{=} \FloatTok{3.0}\NormalTok{, }\FloatTok{4.0}
    \ControlFlowTok{return}\NormalTok{ a }\OperatorTok{+}\NormalTok{ b}

\KeywordTok{def}\NormalTok{ soma\_e\_sqrt():}
\NormalTok{    a, b }\OperatorTok{=} \FloatTok{3.0}\NormalTok{, }\FloatTok{4.0}
    \ControlFlowTok{return}\NormalTok{ math.sqrt(a}\OperatorTok{*}\NormalTok{a }\OperatorTok{+}\NormalTok{ b}\OperatorTok{*}\NormalTok{b)}

\NormalTok{t\_soma }\OperatorTok{=}\NormalTok{ timeit.timeit(apenas\_soma, number}\OperatorTok{=}\NormalTok{N)}
\NormalTok{t\_sqrt }\OperatorTok{=}\NormalTok{ timeit.timeit(soma\_e\_sqrt, number}\OperatorTok{=}\NormalTok{N)}

\BuiltInTok{print}\NormalTok{(}\SpecialStringTok{f"Soma simples       : }\SpecialCharTok{\{}\NormalTok{t\_soma}\SpecialCharTok{:.3f\}}\SpecialStringTok{ s"}\NormalTok{)}
\BuiltInTok{print}\NormalTok{(}\SpecialStringTok{f"Soma + sqrt        : }\SpecialCharTok{\{}\NormalTok{t\_sqrt}\SpecialCharTok{:.3f\}}\SpecialStringTok{ s"}\NormalTok{)}
\BuiltInTok{print}\NormalTok{(}\SpecialStringTok{f"Razão (sqrt/soma)  : }\SpecialCharTok{\{}\NormalTok{t\_sqrt}\OperatorTok{/}\NormalTok{t\_soma}\SpecialCharTok{:.2f\}}\SpecialStringTok{x"}\NormalTok{)}
\end{Highlighting}
\end{Shaded}

\begin{verbatim}
Soma simples       : 4.866 s
Soma + sqrt        : 9.350 s
Razão (sqrt/soma)  : 1.92x
\end{verbatim}

\begin{figure}

\centering{

\begin{verbatim}
<IPython.core.display.HTML object>
\end{verbatim}

}

\caption{\label{fig-01-metricas-distancia}Simulador Compacto:
Euclidiana, City-block e Chessboard}

\end{figure}%

\subsubsection{\texorpdfstring{\emoji{snake}
Python}{ Python}}\label{python}

Basta criar uma célula de código normal e inserir o código Python. A
entrada pode ser simulada usando \texttt{input()}, que funciona
normalmente.

Exemplo de célula:

\begin{Shaded}
\begin{Highlighting}[]
\OperatorTok{\%\%}\NormalTok{writefile EP01\_01.py}
\CommentTok{\# Código Python}
\NormalTok{x1,y1,x2,y2 }\OperatorTok{=} \BuiltInTok{int}\NormalTok{(}\BuiltInTok{input}\NormalTok{()), }\BuiltInTok{int}\NormalTok{(}\BuiltInTok{input}\NormalTok{()), }\BuiltInTok{int}\NormalTok{(}\BuiltInTok{input}\NormalTok{()), }\BuiltInTok{int}\NormalTok{(}\BuiltInTok{input}\NormalTok{())}
\CommentTok{\# Cálculo das diferenças}
\NormalTok{dx }\OperatorTok{=} \BuiltInTok{abs}\NormalTok{(x2 }\OperatorTok{{-}}\NormalTok{ x1)}
\NormalTok{dy }\OperatorTok{=} \BuiltInTok{abs}\NormalTok{(y2 }\OperatorTok{{-}}\NormalTok{ y1)}

\CommentTok{\# 1. Distância Euclidiana (L2)}
\NormalTok{dist\_euclidiana }\OperatorTok{=}\NormalTok{ (dx}\OperatorTok{**}\DecValTok{2} \OperatorTok{+}\NormalTok{ dy}\OperatorTok{**}\DecValTok{2}\NormalTok{)}\OperatorTok{**}\FloatTok{0.5}

\CommentTok{\# 2. Distância City{-}block / Manhattan (L1)}
\NormalTok{dist\_city\_block }\OperatorTok{=}\NormalTok{ dx }\OperatorTok{+}\NormalTok{ dy}

\CommentTok{\# 3. Distância Chessboard / Chebyshev (Linf)}
\NormalTok{dist\_chessboard }\OperatorTok{=} \BuiltInTok{max}\NormalTok{(dx, dy)}

\CommentTok{\# Saída formatada conforme os casos de teste}
\BuiltInTok{print}\NormalTok{(}\SpecialStringTok{f"}\SpecialCharTok{\{}\NormalTok{dist\_euclidiana}\SpecialCharTok{:.2f\}}\SpecialStringTok{"}\NormalTok{)}
\BuiltInTok{print}\NormalTok{(}\SpecialStringTok{f"}\SpecialCharTok{\{}\NormalTok{dist\_city\_block}\SpecialCharTok{:.2f\}}\SpecialStringTok{"}\NormalTok{)}
\BuiltInTok{print}\NormalTok{(}\SpecialStringTok{f"}\SpecialCharTok{\{}\NormalTok{dist\_chessboard}\SpecialCharTok{:.2f\}}\SpecialStringTok{"}\NormalTok{)}
\end{Highlighting}
\end{Shaded}

\begin{verbatim}
Writing EP01_01.py
\end{verbatim}

\begin{Shaded}
\begin{Highlighting}[]
\CommentTok{\# Espera que você digite 4 números inteiros ao executar esta célula.}
\CommentTok{\# No Jupyter ou Google Colab, a mágica \%run {-}i permite que o script leia do teclado.}
\CommentTok{\# No terminal comum, você usaria: python3 EP01\_01.py (sem o \textquotesingle{}!\textquotesingle{} e sem \textquotesingle{}\%run\textquotesingle{}).}

\CommentTok{\# \%run {-}i EP01\_01.py}
\end{Highlighting}
\end{Shaded}

\begin{Shaded}
\begin{Highlighting}[]
\CommentTok{\# Envia 4 inteiros como entrada padrão (stdin) para o script EP01\_01.py usando um pipe}
\OperatorTok{!}\NormalTok{echo }\OperatorTok{{-}}\NormalTok{e }\StringTok{"0}\CharTok{\textbackslash{}n}\StringTok{0}\CharTok{\textbackslash{}n}\StringTok{4}\CharTok{\textbackslash{}n}\StringTok{4"} \OperatorTok{|}\NormalTok{ python3 EP01\_01.py}
\end{Highlighting}
\end{Shaded}

\begin{verbatim}
5.66
8.00
4.00
\end{verbatim}

\begin{Shaded}
\begin{Highlighting}[]
\NormalTok{TestSuite(}\StringTok{"EP01\_01.py"}\NormalTok{).run()}
\end{Highlighting}
\end{Shaded}

\begin{verbatim}
<IPython.core.display.HTML object>
\end{verbatim}

\begin{center}\rule{0.5\linewidth}{0.5pt}\end{center}

\subsubsection{\texorpdfstring{\emoji{hot-beverage}
Java}{ Java}}\label{java}

Para executar Java no Colab, você precisa usar uma célula com o prefixo
\texttt{\%\%writefile} para salvar o código em arquivo, compilar e
executar.

\begin{Shaded}
\begin{Highlighting}[]
\OperatorTok{\%\%}\NormalTok{writefile EP01\_01.java}
\ImportTok{import}\NormalTok{ java.util.Scanner}\OperatorTok{;}

\KeywordTok{class}\NormalTok{ EP01\_01 \{}
\NormalTok{    public static void main(String[] args) \{}
\NormalTok{        Scanner s }\OperatorTok{=}\NormalTok{ new Scanner(System.}\KeywordTok{in}\NormalTok{)}\OperatorTok{;}
        
\NormalTok{        double x1 }\OperatorTok{=}\NormalTok{ s.nextDouble(), y1 }\OperatorTok{=}\NormalTok{ s.nextDouble()}\OperatorTok{;}
\NormalTok{        double x2 }\OperatorTok{=}\NormalTok{ s.nextDouble(), y2 }\OperatorTok{=}\NormalTok{ s.nextDouble()}\OperatorTok{;}
        
\NormalTok{        double dx }\OperatorTok{=}\NormalTok{ Math.}\BuiltInTok{abs}\NormalTok{(x2 }\OperatorTok{{-}}\NormalTok{ x1)}\OperatorTok{;}
\NormalTok{        double dy }\OperatorTok{=}\NormalTok{ Math.}\BuiltInTok{abs}\NormalTok{(y2 }\OperatorTok{{-}}\NormalTok{ y1)}\OperatorTok{;}
        
\NormalTok{        System.out.printf(}\StringTok{"}\SpecialCharTok{\%.2f}\CharTok{\textbackslash{}n}\StringTok{"}\NormalTok{, Math.sqrt(dx}\OperatorTok{*}\NormalTok{dx }\OperatorTok{+}\NormalTok{ dy}\OperatorTok{*}\NormalTok{dy))}\OperatorTok{;}
\NormalTok{        System.out.printf(}\StringTok{"}\SpecialCharTok{\%.2f}\CharTok{\textbackslash{}n}\StringTok{"}\NormalTok{, dx }\OperatorTok{+}\NormalTok{ dy)}\OperatorTok{;}
\NormalTok{        System.out.printf(}\StringTok{"}\SpecialCharTok{\%.2f}\CharTok{\textbackslash{}n}\StringTok{"}\NormalTok{, Math.}\BuiltInTok{max}\NormalTok{(dx, dy))}\OperatorTok{;}
\NormalTok{    \}}
\NormalTok{\}}
\end{Highlighting}
\end{Shaded}

\begin{verbatim}
Writing EP01_01.java
\end{verbatim}

\begin{Shaded}
\begin{Highlighting}[]
\OperatorTok{!}\NormalTok{javac EP01\_01.java}
\OperatorTok{!}\NormalTok{echo }\OperatorTok{{-}}\NormalTok{e }\StringTok{"0}\CharTok{\textbackslash{}n}\StringTok{0}\CharTok{\textbackslash{}n}\StringTok{4}\CharTok{\textbackslash{}n}\StringTok{4"} \OperatorTok{|}\NormalTok{ java EP01\_01}
\end{Highlighting}
\end{Shaded}

\begin{verbatim}
5,66
8,00
4,00
\end{verbatim}

\begin{Shaded}
\begin{Highlighting}[]
\NormalTok{TestSuite(}\StringTok{"EP01\_01.java"}\NormalTok{).run()}
\end{Highlighting}
\end{Shaded}

\begin{verbatim}
<IPython.core.display.HTML object>
\end{verbatim}

\begin{center}\rule{0.5\linewidth}{0.5pt}\end{center}

\subsubsection{\texorpdfstring{\emoji{laptop} C}{ C}}\label{c}

De forma similar, use \texttt{\%\%writefile} para salvar o código,
depois compile e execute. Para C, utilizamos o compilador \textbf{GCC}.

\begin{Shaded}
\begin{Highlighting}[]
\CommentTok{\# Instalar o compilador GCC e ferramentas de build}
\CommentTok{\# O build{-}essential inclui gcc, g++, make, etc.}

\ImportTok{import}\NormalTok{ platform, subprocess}
\NormalTok{s }\OperatorTok{=}\NormalTok{ platform.system()}
\ControlFlowTok{if}\NormalTok{ s }\OperatorTok{==} \StringTok{"Linux"}\NormalTok{:}
    \OperatorTok{!}\NormalTok{sudo apt}\OperatorTok{{-}}\NormalTok{get update }\OperatorTok{{-}}\NormalTok{qq }\OperatorTok{\&\&}\NormalTok{ sudo apt}\OperatorTok{{-}}\NormalTok{get install }\OperatorTok{{-}}\NormalTok{y build}\OperatorTok{{-}}\NormalTok{essential }\OperatorTok{{-}}\NormalTok{qq}
\ControlFlowTok{elif}\NormalTok{ s }\OperatorTok{==} \StringTok{"Darwin"}\NormalTok{:}
    \ControlFlowTok{if}\NormalTok{ subprocess.run([}\StringTok{"which"}\NormalTok{,}\StringTok{"gcc"}\NormalTok{],capture\_output}\OperatorTok{=}\VariableTok{True}\NormalTok{).returncode }\OperatorTok{!=} \DecValTok{0}\NormalTok{:}
        \BuiltInTok{print}\NormalTok{(}\StringTok{"⚠️ Instale o Xcode Command Line Tools: xcode{-}select {-}{-}install"}\NormalTok{)}
\ControlFlowTok{else}\NormalTok{:}
    \BuiltInTok{print}\NormalTok{(}\StringTok{"⚠️ Windows: use WSL (recomendado) ou MinGW (https://www.mingw{-}w64.org)"}\NormalTok{)}
\BuiltInTok{print}\NormalTok{(}\StringTok{"✅ Ambiente C/C++ pronto."}\NormalTok{)}
\end{Highlighting}
\end{Shaded}

\begin{verbatim}
✅ Ambiente C/C++ pronto.
\end{verbatim}

\begin{Shaded}
\begin{Highlighting}[]
\OperatorTok{\%\%}\NormalTok{writefile EP01\_01.c}
\CommentTok{\#include \textless{}stdio.h\textgreater{}}
\CommentTok{\#include \textless{}math.h\textgreater{}}

\BuiltInTok{int}\NormalTok{ main() \{}
\NormalTok{    double x1, y1, x2, y2}\OperatorTok{;}
    \ControlFlowTok{if}\NormalTok{ (scanf(}\StringTok{"}\SpecialCharTok{\%lf}\StringTok{ }\SpecialCharTok{\%lf}\StringTok{ }\SpecialCharTok{\%lf}\StringTok{ }\SpecialCharTok{\%lf}\StringTok{"}\NormalTok{, }\OperatorTok{\&}\NormalTok{x1, }\OperatorTok{\&}\NormalTok{y1, }\OperatorTok{\&}\NormalTok{x2, }\OperatorTok{\&}\NormalTok{y2) }\OperatorTok{!=} \DecValTok{4}\NormalTok{) }\ControlFlowTok{return} \DecValTok{0}\OperatorTok{;}

\NormalTok{    double dx }\OperatorTok{=}\NormalTok{ fabs(x2 }\OperatorTok{{-}}\NormalTok{ x1)}\OperatorTok{;}
\NormalTok{    double dy }\OperatorTok{=}\NormalTok{ fabs(y2 }\OperatorTok{{-}}\NormalTok{ y1)}\OperatorTok{;}

    \OperatorTok{//}\NormalTok{ Euclidiana, City}\OperatorTok{{-}}\NormalTok{block e Chessboard}
\NormalTok{    printf(}\StringTok{"}\SpecialCharTok{\%.2f}\CharTok{\textbackslash{}n}\StringTok{"}\NormalTok{, sqrt(dx }\OperatorTok{*}\NormalTok{ dx }\OperatorTok{+}\NormalTok{ dy }\OperatorTok{*}\NormalTok{ dy))}\OperatorTok{;}
\NormalTok{    printf(}\StringTok{"}\SpecialCharTok{\%.2f}\CharTok{\textbackslash{}n}\StringTok{"}\NormalTok{, dx }\OperatorTok{+}\NormalTok{ dy)}\OperatorTok{;}
\NormalTok{    printf(}\StringTok{"}\SpecialCharTok{\%.2f}\CharTok{\textbackslash{}n}\StringTok{"}\NormalTok{, fmax(dx, dy))}\OperatorTok{;}

    \ControlFlowTok{return} \DecValTok{0}\OperatorTok{;}
\NormalTok{\}}
\end{Highlighting}
\end{Shaded}

\begin{verbatim}
Writing EP01_01.c
\end{verbatim}

\begin{Shaded}
\begin{Highlighting}[]
\CommentTok{\# Compila o arquivo .c gerando o executável EP01\_01}
\CommentTok{\# {-}lm é usado para linkar a biblioteca matemática (math.h) se necessário}

\OperatorTok{!}\NormalTok{gcc EP01\_01.c }\OperatorTok{{-}}\NormalTok{o EP01\_01 }\OperatorTok{{-}}\NormalTok{lm}
\OperatorTok{!}\NormalTok{echo }\OperatorTok{{-}}\NormalTok{e }\StringTok{"0}\CharTok{\textbackslash{}n}\StringTok{0}\CharTok{\textbackslash{}n}\StringTok{4}\CharTok{\textbackslash{}n}\StringTok{4"} \OperatorTok{|}\NormalTok{ .}\OperatorTok{/}\NormalTok{EP01\_01}
\end{Highlighting}
\end{Shaded}

\begin{verbatim}
5.66
8.00
4.00
\end{verbatim}

\begin{Shaded}
\begin{Highlighting}[]
\NormalTok{TestSuite(}\StringTok{"EP01\_01.c"}\NormalTok{).run()}
\end{Highlighting}
\end{Shaded}

\begin{verbatim}
<IPython.core.display.HTML object>
\end{verbatim}

\begin{center}\rule{0.5\linewidth}{0.5pt}\end{center}

\subsubsection{\texorpdfstring{\emoji{laptop} C++}{ C++}}\label{c-1}

De forma similar ao Java, use \texttt{\%\%writefile} para salvar o
código, depois compile e execute. Lembre-se de que, no Colab, também é
necessário instalar o seguinte:

\begin{Shaded}
\begin{Highlighting}[]
\CommentTok{\# Instalar o compilador G++ para C++}
\CommentTok{\# O build{-}essential inclui o g++, make, etc.}

\ImportTok{import}\NormalTok{ platform, subprocess}
\NormalTok{s }\OperatorTok{=}\NormalTok{ platform.system()}
\ControlFlowTok{if}\NormalTok{ s }\OperatorTok{==} \StringTok{"Linux"}\NormalTok{:}
    \OperatorTok{!}\NormalTok{sudo apt}\OperatorTok{{-}}\NormalTok{get update }\OperatorTok{{-}}\NormalTok{qq }\OperatorTok{\&\&}\NormalTok{ sudo apt}\OperatorTok{{-}}\NormalTok{get install }\OperatorTok{{-}}\NormalTok{y build}\OperatorTok{{-}}\NormalTok{essential }\OperatorTok{{-}}\NormalTok{qq}
\ControlFlowTok{elif}\NormalTok{ s }\OperatorTok{==} \StringTok{"Darwin"} \KeywordTok{and}\NormalTok{ subprocess.run([}\StringTok{"which"}\NormalTok{,}\StringTok{"g++"}\NormalTok{],capture\_output}\OperatorTok{=}\VariableTok{True}\NormalTok{).returncode }\OperatorTok{!=} \DecValTok{0}\NormalTok{:}
    \BuiltInTok{print}\NormalTok{(}\StringTok{"⚠️ Instale o Xcode Command Line Tools: xcode{-}select {-}{-}install"}\NormalTok{)}
\ControlFlowTok{elif}\NormalTok{ s }\OperatorTok{==} \StringTok{"Windows"}\NormalTok{:}
    \BuiltInTok{print}\NormalTok{(}\StringTok{"⚠️ Windows: use WSL (recomendado) ou MinGW (https://www.mingw{-}w64.org)"}\NormalTok{)}
\BuiltInTok{print}\NormalTok{(}\StringTok{"✅ Compilador C++ pronto."}\NormalTok{)}
\end{Highlighting}
\end{Shaded}

\begin{verbatim}
✅ Compilador C++ pronto.
\end{verbatim}

\begin{Shaded}
\begin{Highlighting}[]
\OperatorTok{\%\%}\NormalTok{writefile EP01\_01.cpp}
\CommentTok{\#include \textless{}iostream\textgreater{}}
\CommentTok{\#include \textless{}iomanip\textgreater{}}
\CommentTok{\#include \textless{}cmath\textgreater{}}
\CommentTok{\#include \textless{}algorithm\textgreater{}}

\BuiltInTok{int}\NormalTok{ main() \{}
\NormalTok{    double x1, y1, x2, y2}\OperatorTok{;}
    \ControlFlowTok{if}\NormalTok{ (}\OperatorTok{!}\NormalTok{(std::cin }\OperatorTok{\textgreater{}\textgreater{}}\NormalTok{ x1 }\OperatorTok{\textgreater{}\textgreater{}}\NormalTok{ y1 }\OperatorTok{\textgreater{}\textgreater{}}\NormalTok{ x2 }\OperatorTok{\textgreater{}\textgreater{}}\NormalTok{ y2)) }\ControlFlowTok{return} \DecValTok{0}\OperatorTok{;}

\NormalTok{    double dx }\OperatorTok{=}\NormalTok{ std::}\BuiltInTok{abs}\NormalTok{(x2 }\OperatorTok{{-}}\NormalTok{ x1)}\OperatorTok{;}
\NormalTok{    double dy }\OperatorTok{=}\NormalTok{ std::}\BuiltInTok{abs}\NormalTok{(y2 }\OperatorTok{{-}}\NormalTok{ y1)}\OperatorTok{;}

\NormalTok{    std::cout }\OperatorTok{\textless{}\textless{}}\NormalTok{ std::fixed }\OperatorTok{\textless{}\textless{}}\NormalTok{ std::setprecision(}\DecValTok{2}\NormalTok{)}\OperatorTok{;}
    
    \OperatorTok{//}\NormalTok{ Euclidiana, City}\OperatorTok{{-}}\NormalTok{block e Chessboard}
\NormalTok{    std::cout }\OperatorTok{\textless{}\textless{}}\NormalTok{ std::sqrt(dx}\OperatorTok{*}\NormalTok{dx }\OperatorTok{+}\NormalTok{ dy}\OperatorTok{*}\NormalTok{dy) }\OperatorTok{\textless{}\textless{}}\NormalTok{ std::endl}\OperatorTok{;}
\NormalTok{    std::cout }\OperatorTok{\textless{}\textless{}}\NormalTok{ (dx }\OperatorTok{+}\NormalTok{ dy) }\OperatorTok{\textless{}\textless{}}\NormalTok{ std::endl}\OperatorTok{;}
\NormalTok{    std::cout }\OperatorTok{\textless{}\textless{}}\NormalTok{ std::}\BuiltInTok{max}\NormalTok{(dx, dy) }\OperatorTok{\textless{}\textless{}}\NormalTok{ std::endl}\OperatorTok{;}

    \ControlFlowTok{return} \DecValTok{0}\OperatorTok{;}
\NormalTok{\}}
\end{Highlighting}
\end{Shaded}

\begin{verbatim}
Writing EP01_01.cpp
\end{verbatim}

\begin{Shaded}
\begin{Highlighting}[]
\OperatorTok{!}\NormalTok{g}\OperatorTok{++}\NormalTok{ EP01\_01.cpp }\OperatorTok{{-}}\NormalTok{o EP01\_01}
\OperatorTok{!}\NormalTok{echo }\OperatorTok{{-}}\NormalTok{e }\StringTok{"0}\CharTok{\textbackslash{}n}\StringTok{0}\CharTok{\textbackslash{}n}\StringTok{4}\CharTok{\textbackslash{}n}\StringTok{4"} \OperatorTok{|}\NormalTok{ .}\OperatorTok{/}\NormalTok{EP01\_01}
\end{Highlighting}
\end{Shaded}

\begin{verbatim}
5.66
8.00
4.00
\end{verbatim}

\begin{Shaded}
\begin{Highlighting}[]
\NormalTok{TestSuite(}\StringTok{"EP01\_01.cpp"}\NormalTok{).run()}
\end{Highlighting}
\end{Shaded}

\begin{verbatim}
<IPython.core.display.HTML object>
\end{verbatim}

\begin{center}\rule{0.5\linewidth}{0.5pt}\end{center}

\subsubsection{\texorpdfstring{\emoji{globe-with-meridians} JavaScript
(Node.js)}{ JavaScript (Node.js)}}\label{javascript-node.js}

Para JavaScript, use \texttt{\%\%writefile} para criar o arquivo e
execute com Node:

\begin{Shaded}
\begin{Highlighting}[]
\OperatorTok{\%\%}\NormalTok{writefile EP01\_01.js}
\NormalTok{function escreva(s) \{ }\ControlFlowTok{try}\NormalTok{ \{ document.write(s }\OperatorTok{+} \StringTok{"\textless{}br\textgreater{}"}\NormalTok{)}\OperatorTok{;}\NormalTok{ \} catch(e) \{ console.log(s)}\OperatorTok{;}\NormalTok{ \} \}}

\NormalTok{process.stdin.once(}\StringTok{\textquotesingle{}data\textquotesingle{}}\NormalTok{, data }\OperatorTok{=\textgreater{}}\NormalTok{ \{}
\NormalTok{  const valores }\OperatorTok{=}\NormalTok{ data.toString().trim().split(}\OperatorTok{/}\NormalTok{\textbackslash{}s}\OperatorTok{+/}\NormalTok{)}\OperatorTok{;}
\NormalTok{  const x1 }\OperatorTok{=}\NormalTok{ parseFloat(valores[}\DecValTok{0}\NormalTok{])}\OperatorTok{;}
\NormalTok{  const y1 }\OperatorTok{=}\NormalTok{ parseFloat(valores[}\DecValTok{1}\NormalTok{])}\OperatorTok{;}
\NormalTok{  const x2 }\OperatorTok{=}\NormalTok{ parseFloat(valores[}\DecValTok{2}\NormalTok{])}\OperatorTok{;}
\NormalTok{  const y2 }\OperatorTok{=}\NormalTok{ parseFloat(valores[}\DecValTok{3}\NormalTok{])}\OperatorTok{;}

\NormalTok{  const dx }\OperatorTok{=}\NormalTok{ Math.}\BuiltInTok{abs}\NormalTok{(x2 }\OperatorTok{{-}}\NormalTok{ x1)}\OperatorTok{;}
\NormalTok{  const dy }\OperatorTok{=}\NormalTok{ Math.}\BuiltInTok{abs}\NormalTok{(y2 }\OperatorTok{{-}}\NormalTok{ y1)}\OperatorTok{;}

  \OperatorTok{//}\NormalTok{ Euclidiana, City}\OperatorTok{{-}}\NormalTok{block e Chessboard}
\NormalTok{  escreva(Math.sqrt(dx }\OperatorTok{*}\NormalTok{ dx }\OperatorTok{+}\NormalTok{ dy }\OperatorTok{*}\NormalTok{ dy).toFixed(}\DecValTok{2}\NormalTok{))}\OperatorTok{;}
\NormalTok{  escreva((dx }\OperatorTok{+}\NormalTok{ dy).toFixed(}\DecValTok{2}\NormalTok{))}\OperatorTok{;}
\NormalTok{  escreva(Math.}\BuiltInTok{max}\NormalTok{(dx, dy).toFixed(}\DecValTok{2}\NormalTok{))}\OperatorTok{;}
\NormalTok{\})}\OperatorTok{;}
\end{Highlighting}
\end{Shaded}

\begin{verbatim}
Writing EP01_01.js
\end{verbatim}

\begin{Shaded}
\begin{Highlighting}[]
\NormalTok{TestSuite(}\StringTok{"EP01\_01.js"}\NormalTok{).run()}
\end{Highlighting}
\end{Shaded}

\begin{verbatim}
<IPython.core.display.HTML object>
\end{verbatim}

\begin{center}\rule{0.5\linewidth}{0.5pt}\end{center}

\subsubsection{\texorpdfstring{\emoji{bar-chart} R}{ R}}\label{r}

No Colab, você pode executar código R diretamente utilizando a mágica
\texttt{\%\%R}.\\
O programa deve ler \textbf{quatro números} (x1, y1, x2, y2) e exibir a
distância euclidiana com \textbf{duas casas decimais}.

Exemplo de célula:

\begin{Shaded}
\begin{Highlighting}[]
\SpecialCharTok{\%\%}\NormalTok{R}
\NormalTok{dados }\OtherTok{\textless{}{-}} \FunctionTok{scan}\NormalTok{(}\AttributeTok{file =} \StringTok{"stdin"}\NormalTok{, }\AttributeTok{n =} \DecValTok{4}\NormalTok{, }\AttributeTok{quiet =} \ConstantTok{TRUE}\NormalTok{)}
\NormalTok{x1 }\OtherTok{\textless{}{-}}\NormalTok{ dados[}\DecValTok{1}\NormalTok{]; y1 }\OtherTok{\textless{}{-}}\NormalTok{ dados[}\DecValTok{2}\NormalTok{]; x2 }\OtherTok{\textless{}{-}}\NormalTok{ dados[}\DecValTok{3}\NormalTok{]; y2 }\OtherTok{\textless{}{-}}\NormalTok{ dados[}\DecValTok{4}\NormalTok{]}
\NormalTok{dist }\OtherTok{\textless{}{-}} \FunctionTok{sqrt}\NormalTok{((x2 }\SpecialCharTok{{-}}\NormalTok{ x1)}\SpecialCharTok{\^{}}\DecValTok{2} \SpecialCharTok{+}\NormalTok{ (y2 }\SpecialCharTok{{-}}\NormalTok{ y1)}\SpecialCharTok{\^{}}\DecValTok{2}\NormalTok{)}
\FunctionTok{cat}\NormalTok{(}\FunctionTok{sprintf}\NormalTok{(}\StringTok{"\%.2f"}\NormalTok{, dist))}
\end{Highlighting}
\end{Shaded}

\begin{Shaded}
\begin{Highlighting}[]
\CommentTok{\# Instalar o R e o Rscript}
\CommentTok{\# O r{-}base instala o ambiente R completo, incluindo o Rscript}

\ImportTok{import}\NormalTok{ platform, subprocess}
\NormalTok{s }\OperatorTok{=}\NormalTok{ platform.system()}
\ControlFlowTok{if}\NormalTok{ s }\OperatorTok{==} \StringTok{"Linux"}\NormalTok{:}
    \OperatorTok{!}\NormalTok{sudo apt}\OperatorTok{{-}}\NormalTok{get update }\OperatorTok{{-}}\NormalTok{qq }\OperatorTok{\&\&}\NormalTok{ sudo apt}\OperatorTok{{-}}\NormalTok{get install }\OperatorTok{{-}}\NormalTok{y r}\OperatorTok{{-}}\NormalTok{base }\OperatorTok{{-}}\NormalTok{qq}
\ControlFlowTok{elif}\NormalTok{ s }\OperatorTok{==} \StringTok{"Darwin"}\NormalTok{:}
    \ControlFlowTok{if}\NormalTok{ subprocess.run([}\StringTok{"which"}\NormalTok{,}\StringTok{"R"}\NormalTok{],capture\_output}\OperatorTok{=}\VariableTok{True}\NormalTok{).returncode }\OperatorTok{!=} \DecValTok{0}\NormalTok{:}
        \BuiltInTok{print}\NormalTok{(}\StringTok{"⚠️ Instale o R no macOS: https://cran.r{-}project.org/bin/macosx/"}\NormalTok{)}
\ControlFlowTok{else}\NormalTok{:}
    \BuiltInTok{print}\NormalTok{(}\StringTok{"⚠️ Windows: baixe o R em https://cran.r{-}project.org/bin/windows/base/"}\NormalTok{)}
\BuiltInTok{print}\NormalTok{(}\StringTok{"✅ Ambiente R pronto."}\NormalTok{)}
\end{Highlighting}
\end{Shaded}

\begin{verbatim}
✅ Ambiente R pronto.
\end{verbatim}

Para testar da mesma forma que nos exemplos anteriores, você pode usar a
entrada padrão (stdin) no terminal ou adaptar o código conforme abaixo:

\begin{Shaded}
\begin{Highlighting}[]
\OperatorTok{\%\%}\NormalTok{writefile EP01\_01.r}
\NormalTok{dados }\OperatorTok{\textless{}{-}}\NormalTok{ scan(}\BuiltInTok{file} \OperatorTok{=} \StringTok{"stdin"}\NormalTok{, n }\OperatorTok{=} \DecValTok{4}\NormalTok{, quiet }\OperatorTok{=}\NormalTok{ TRUE)}
\NormalTok{dx }\OperatorTok{\textless{}{-}} \BuiltInTok{abs}\NormalTok{(dados[}\DecValTok{3}\NormalTok{] }\OperatorTok{{-}}\NormalTok{ dados[}\DecValTok{1}\NormalTok{])}
\NormalTok{dy }\OperatorTok{\textless{}{-}} \BuiltInTok{abs}\NormalTok{(dados[}\DecValTok{4}\NormalTok{] }\OperatorTok{{-}}\NormalTok{ dados[}\DecValTok{2}\NormalTok{])}

\CommentTok{\# Saídas: Euclidiana, City{-}block e Chessboard}
\NormalTok{cat(sprintf(}\StringTok{"}\SpecialCharTok{\%.2f}\CharTok{\textbackslash{}n}\SpecialCharTok{\%.2f}\CharTok{\textbackslash{}n}\SpecialCharTok{\%.2f}\CharTok{\textbackslash{}n}\StringTok{"}\NormalTok{, }
\NormalTok{    sqrt(dx}\OperatorTok{\^{}}\DecValTok{2} \OperatorTok{+}\NormalTok{ dy}\OperatorTok{\^{}}\DecValTok{2}\NormalTok{), }
\NormalTok{    dx }\OperatorTok{+}\NormalTok{ dy, }
    \BuiltInTok{max}\NormalTok{(dx, dy)))}
\end{Highlighting}
\end{Shaded}

\begin{verbatim}
Overwriting EP01_01.r
\end{verbatim}

\begin{Shaded}
\begin{Highlighting}[]
\OperatorTok{!}\NormalTok{echo }\OperatorTok{{-}}\NormalTok{e }\StringTok{"0}\CharTok{\textbackslash{}n}\StringTok{0}\CharTok{\textbackslash{}n}\StringTok{4}\CharTok{\textbackslash{}n}\StringTok{4"} \OperatorTok{|}\NormalTok{ Rscript EP01\_01.r}
\end{Highlighting}
\end{Shaded}

\begin{verbatim}
5.66
8.00
4.00
\end{verbatim}

\begin{Shaded}
\begin{Highlighting}[]
\NormalTok{TestSuite(}\StringTok{"EP01\_01.r"}\NormalTok{).run()}
\end{Highlighting}
\end{Shaded}

\begin{verbatim}
<IPython.core.display.HTML object>
\end{verbatim}

\subsection{\texorpdfstring{EP01\_02 \emoji{bar-chart} Desempenho
Preditivo --- Métricas de ML em
VC}{EP01\_02  Desempenho Preditivo --- Métricas de ML em VC}}\label{ep01_02-desempenho-preditivo-muxe9tricas-de-ml-em-vc}

Nesta atividade, você entrará no mundo do \textbf{Aprendizado de
Máquina} (\emph{Machine Learning}). Seu objetivo é avaliar o desempenho
de um classificador binário calculando métricas a partir de uma
\textbf{Matriz de Confusão}.

\subsubsection{\texorpdfstring{\emoji{brain} Por que a métrica certa
importa?}{ Por que a métrica certa importa?}}\label{por-que-a-muxe9trica-certa-importa}

Imagine um detector de \textbf{moedas de R\$ 1,00}. O impacto do erro
define a métrica prioritária:

\begin{longtable}[]{@{}
  >{\raggedright\arraybackslash}p{(\linewidth - 4\tabcolsep) * \real{0.3333}}
  >{\raggedright\arraybackslash}p{(\linewidth - 4\tabcolsep) * \real{0.3333}}
  >{\raggedright\arraybackslash}p{(\linewidth - 4\tabcolsep) * \real{0.3333}}@{}}
\toprule\noalign{}
\begin{minipage}[b]{\linewidth}\raggedright
Métrica
\end{minipage} & \begin{minipage}[b]{\linewidth}\raggedright
Exemplo Prático
\end{minipage} & \begin{minipage}[b]{\linewidth}\raggedright
Importância em PDI/VC
\end{minipage} \\
\midrule\noalign{}
\endhead
\bottomrule\noalign{}
\endlastfoot
\textbf{Acurácia} & \textbf{Contagem de Grãos} & Útil quando as classes
são equilibradas (ex: metade dos grãos com defeito, metade
saudáveis). \\
\textbf{Precisão} & \textbf{Segurança/Biometria} & Crucial para evitar
\textbf{Falsos Positivos} (ex: não permitir que um impostor acesse um
sistema por erro de reconhecimento). \\
\textbf{Sensibilidade} & \textbf{Saúde (Tumores)} & Crucial para evitar
\textbf{Falsos Negativos} (ex: não deixar um tumor passar despercebido
em um exame de Raio-X). \\
\textbf{F1-score} & \textbf{Cédulas de Dinheiro} & Ideal para um
\textbf{equilíbrio} entre não rejeitar notas verdadeiras e não aceitar
notas falsas. \\
\end{longtable}

\subsubsection{\texorpdfstring{\emoji{bar-chart} A Matriz de
Confusão}{ A Matriz de Confusão}}\label{a-matriz-de-confusuxe3o}

\begin{longtable}[]{@{}
  >{\raggedright\arraybackslash}p{(\linewidth - 4\tabcolsep) * \real{0.3333}}
  >{\raggedright\arraybackslash}p{(\linewidth - 4\tabcolsep) * \real{0.3333}}
  >{\raggedright\arraybackslash}p{(\linewidth - 4\tabcolsep) * \real{0.3333}}@{}}
\toprule\noalign{}
\begin{minipage}[b]{\linewidth}\raggedright
\end{minipage} & \begin{minipage}[b]{\linewidth}\raggedright
\textbf{Predita Positiva}
\end{minipage} & \begin{minipage}[b]{\linewidth}\raggedright
\textbf{Predita Negativa}
\end{minipage} \\
\midrule\noalign{}
\endhead
\bottomrule\noalign{}
\endlastfoot
\textbf{Real Positiva} & \textbf{VP} (Verdadeiro Positivo) & \textbf{FN}
(Falso Negativo) \\
\textbf{Real Negativa} & \textbf{FP} (Falso Positivo) & \textbf{VN}
(Verdadeiro Negativo) \\
\end{longtable}

\textbf{Sua tarefa:}

\begin{enumerate}
\def\labelenumi{\arabic{enumi}.}
\tightlist
\item
  Leia \textbf{4 valores inteiros} na ordem: \texttt{VP}, \texttt{FN},
  \texttt{FP}, \texttt{VN}.
\item
  Calcule as métricas utilizando as fórmulas:
\end{enumerate}

\[\text{Acurácia} = \frac{VP + VN}{VP + VN + FP + FN}\]

\[\text{Precisão} = \frac{VP}{VP + FP}\]

\[\text{Sensibilidade (Recall)} = \frac{VP}{VP + FN}\]

\[\text{F1-score} = \frac{2 \times \text{Precisão} \times \text{Sensibilidade}}{\text{Precisão} + \text{Sensibilidade}}\]

\begin{enumerate}
\def\labelenumi{\arabic{enumi}.}
\setcounter{enumi}{2}
\tightlist
\item
  Imprima os resultados formatados com \textbf{duas casas decimais}, um
  por linha.
\end{enumerate}

\emoji{pushpin} \textbf{Importante}:

\begin{itemize}
\tightlist
\item
  Use divisão em ponto flutuante para evitar resultados truncados.
\item
  A ordem de saída deve ser: \textbf{Acurácia}, \textbf{Precisão},
  \textbf{Sensibilidade} e \textbf{F1-score}.
\item
  Veja a simulação interativa deste EP na
  Figura~\ref{fig-01-desempenho-preditivo}.
\end{itemize}

\begin{center}\rule{0.5\linewidth}{0.5pt}\end{center}

\subsubsection{\texorpdfstring{\emoji{pushpin} Exemplo de
Execução}{ Exemplo de Execução}}\label{exemplo-de-execuuxe7uxe3o}

\begin{longtable}[]{@{}lll@{}}
\toprule\noalign{}
Entrada & Saída & Observação \\
\midrule\noalign{}
\endhead
\bottomrule\noalign{}
\endlastfoot
\texttt{40} & \texttt{0.75} & Acurácia \\
\texttt{10} & \texttt{0.73} & Precisão \\
\texttt{15} & \texttt{0.80} & Sensibilidade \\
\texttt{35} & \texttt{0.76} & F1-score \\
\end{longtable}

\emph{(Nota: VP=40, FN=10, FP=15, VN=35. Total de casos = 100)}

\begin{center}\rule{0.5\linewidth}{0.5pt}\end{center}

\begin{figure}

\centering{

\begin{verbatim}
<IPython.core.display.HTML object>
\end{verbatim}

}

\caption{\label{fig-01-desempenho-preditivo}Simulador: Desempenho
Preditivo --- Métricas de ML}

\end{figure}%

\begin{Shaded}
\begin{Highlighting}[]
\OperatorTok{\%\%}\NormalTok{writefile EP01\_02.py}
\CommentTok{\# sua solução}
\end{Highlighting}
\end{Shaded}

\begin{verbatim}
Writing EP01_02.py
\end{verbatim}

\begin{Shaded}
\begin{Highlighting}[]
\NormalTok{TestSuite(}\StringTok{"EP01\_02.py"}\NormalTok{).run()}
\end{Highlighting}
\end{Shaded}

\begin{verbatim}
<IPython.core.display.HTML object>
\end{verbatim}

\begin{center}\rule{0.5\linewidth}{0.5pt}\end{center}

\subsection{\texorpdfstring{EP01\_03 $\nearrow$ \emph{Mean Average
Precision} (mAP) --- Curva
Precisão‑Sensibilidade}{EP01\_03  Mean Average Precision (mAP) --- Curva Precisão‑Sensibilidade}}\label{ep01_03-mean-average-precision-map-curva-precisuxe3osensibilidade}

Nesta atividade, você avaliará um classificador binário (ex.: detecção
de desmatamento em imagens de satélite) através da \textbf{curva
Precisão‑Sensibilidade} e da métrica \textbf{mAP (\emph{Mean Average
Precision})}. O mAP é padrão em competições como
\href{https://cocodataset.org/}{COCO (\emph{Common Objects in Context})}
e \href{http://host.robots.ox.ac.uk/pascal/VOC/}{PASCAL VOC
(\emph{Visual Object Classes})} e dos modelos
\href{https://github.com/ultralytics/ultralytics}{YOLO (\emph{You Only
Look Once})}.

\subsubsection{\texorpdfstring{\emoji{brain} Por que o mAP é a
métrica‑padrão?}{ Por que o mAP é a métrica‑padrão?}}\label{por-que-o-map-uxe9-a-muxe9tricapadruxe3o}

Na EP01\_02 você viu que a escolha do limiar muda totalmente a Precisão
e a Sensibilidade. O \textbf{mAP (Mean Average Precision)} resolve isso:
avalia o modelo em \textbf{vários limiares} (cada limiar deve gerar uma
matriz de confusão diferente) e resume o desempenho pela \textbf{área
sob a curva Precisão‑Sensibilidade (P‑S)}.

Enquanto o F1‑\emph{Score} olha para um único ponto de equilíbrio, o mAP
considera a curva inteira. Quanto mais próximo de \textbf{1,0}, melhor o
detector em todos os limiares e classes (ex.: moedas de 25, 50 e 1
real).

\begin{longtable}[]{@{}
  >{\raggedright\arraybackslash}p{(\linewidth - 4\tabcolsep) * \real{0.2571}}
  >{\raggedright\arraybackslash}p{(\linewidth - 4\tabcolsep) * \real{0.4000}}
  >{\raggedright\arraybackslash}p{(\linewidth - 4\tabcolsep) * \real{0.3429}}@{}}
\toprule\noalign{}
\begin{minipage}[b]{\linewidth}\raggedright
Métrica
\end{minipage} & \begin{minipage}[b]{\linewidth}\raggedright
O que resume
\end{minipage} & \begin{minipage}[b]{\linewidth}\raggedright
Limitação
\end{minipage} \\
\midrule\noalign{}
\endhead
\bottomrule\noalign{}
\endlastfoot
\textbf{F1‑\emph{Score}} & Equilíbrio P × S num único limiar & Depende
do limiar escolhido \\
\textbf{AP} & Área sob a curva P‑S de uma classe & Válida apenas para
uma classe \\
\textbf{mAP} & Média das APs sobre todas as classes & Mais complexo de
implementar \\
\end{longtable}

\begin{quote}
Referências:
\href{https://blog.roboflow.com/mean-average-precision/}{Roboflow ---
mAP} · \href{https://youtu.be/COQPqUH-HG8}{Vídeo explicativo}
\end{quote}

\subsubsection{🔢 Como o mAP é calculado --- passo a
passo}\label{como-o-map-uxe9-calculado-passo-a-passo}

\begin{enumerate}
\def\labelenumi{\arabic{enumi}.}
\item
  \textbf{Limiares fixos} (use sempre esta lista):

\begin{verbatim}
limiares = [0.00, 0.09, 0.21, 0.31, 0.39, 0.52, 0.60, 0.71, 0.81, 0.89, 1.00]
\end{verbatim}
\item
  Para cada limiar (t), classifique as amostras:
  \texttt{predito\ =\ 1\ se\ confiança\ ≥\ t,\ senão\ 0}.\\
  Calcule VP, FP, FN, VN e obtenha Precisão((t)) e Sensibilidade((t)).
\item
  Monte a curva P‑S: pares (Sensibilidade((t)), Precisão((t))),
  ordenados por Sensibilidade crescente.
\item
  \textbf{Monotonize a Precisão}:\\
  \[P_{\text{mono}}[i] = \max_{j \ge i} P[j]\]
\item
  Calcule a \textbf{AP} (área sob a curva monotônica) usando a
  \textbf{regra do trapézio} (aproximação mais precisa que a simples
  soma de Riemann):\\
  \[AP = \sum_{i=1}^{m-1} \frac{P_{\text{mono}}[i-1] + P_{\text{mono}}[i]}{2} \cdot (S[i] - S[i-1])\]
\item
  \textbf{mAP} = média das APs de todas as classes. Neste EP há apenas
  \textbf{1 classe}, portanto mAP = AP.
\end{enumerate}

\begin{tcolorbox}[enhanced jigsaw, titlerule=0mm, left=2mm, opacityback=0, toprule=.15mm, arc=.35mm, opacitybacktitle=0.6, colbacktitle=quarto-callout-note-color!10!white, colframe=quarto-callout-note-color-frame, rightrule=.15mm, breakable, coltitle=black, bottomtitle=1mm, toptitle=1mm, title=\textcolor{quarto-callout-note-color}{\faInfo}\hspace{0.5em}{Nota}, colback=white, bottomrule=.15mm, leftrule=.75mm]

\textbf{\emoji{triangular-ruler} Diferença resumida:}\\
A \emph{soma de Riemann} aproxima a área por retângulos, podendo
subestimar ou superestimar. A \textbf{regra do trapézio} usa trapézios,
reduzindo o erro ao considerar a média entre os valores nos extremos do
intervalo, sendo geralmente mais precisa para funções suaves por partes,
como a curva Precisão‑Sensibilidade.

\end{tcolorbox}

\subsubsection{\texorpdfstring{\emoji{clipboard}
Tarefa}{ Tarefa}}\label{tarefa}

Leia um inteiro \texttt{n} (quantidade de amostras). Em seguida leia
\texttt{n} linhas, cada uma com: \texttt{verdade} (0 ou 1) e
\texttt{confiança} (float 0.0--1.0).

Calcule e imprima, para o \textbf{limiar 0.85} (índice 9 da lista):

\begin{itemize}
\tightlist
\item
  Matriz de Confusão (VP, FN, FP, VN)
\item
  Acurácia, Precisão, Sensibilidade e F1‑Score
\end{itemize}

Em seguida, para \textbf{todos os limiares}, imprima:

\begin{itemize}
\tightlist
\item
  Precisões cruas, Precisões monotônicas e Sensibilidades, separadas por
  \texttt{,}
\item
  mAP final
\end{itemize}

\subsubsection{\texorpdfstring{\emoji{pushpin}
Importante}{ Importante}}\label{importante}

\begin{itemize}
\tightlist
\item
  Limiar fixo para as métricas individuais: \textbf{0.85}
\item
  Divisão segura: se denominador for zero, use 0
\item
  Formatação: duas casas decimais
\item
  Monotonize de trás para frente
\item
  A \textbf{?@fig-01-mAP-interativo} apresenta uma simulação desta
  questão
\end{itemize}

\subsubsection{\texorpdfstring{\emoji{pushpin} Exemplo de
Execução}{ Exemplo de Execução}}\label{exemplo-de-execuuxe7uxe3o-1}

\begin{longtable}[]{@{}
  >{\raggedright\arraybackslash}p{(\linewidth - 2\tabcolsep) * \real{0.3600}}
  >{\raggedright\arraybackslash}p{(\linewidth - 2\tabcolsep) * \real{0.6400}}@{}}
\toprule\noalign{}
\begin{minipage}[b]{\linewidth}\raggedright
Entrada
\end{minipage} & \begin{minipage}[b]{\linewidth}\raggedright
Saída Esperada
\end{minipage} \\
\midrule\noalign{}
\endhead
\bottomrule\noalign{}
\endlastfoot
70 0.941 0.801 0.690 0.671 0.301 0.151 0.15 & \# MÉTRICAS PARA O LIMIAR
0.85 \#Matriz de Confusão:VP = 0, FN = 5FP = 1, VN = 1Métricas de
Avaliação:Acurácia: 0.14Precisão: 0.00Sensibilidade: 0.00F1-Score:
0.00\# MÉTRICAS PARA TODOS OS LIMIARES \#Precisões: 0.00, 0.00, 0.00,
0.50, 0.50, 0.50, 0.50, 0.50, 0.60, 0.71, 0.71Precisões mon.: 0.71,
0.71, 0.71, 0.71, 0.71, 0.71, 0.71, 0.71, 0.71, 0.71,
0.71Sensibilidades: 0.00, 0.00, 0.00, 0.20, 0.40, 0.40, 0.40, 0.40,
0.60, 1.00, 1.00mAP: 0.71 \\
\end{longtable}

\subsubsection{\texorpdfstring{\emoji{snake} Dica para calcular o AP
(com regra do
trapézio)}{ Dica para calcular o AP (com regra do trapézio)}}\label{dica-para-calcular-o-ap-com-regra-do-trapuxe9zio}

\begin{Shaded}
\begin{Highlighting}[]
\KeywordTok{def}\NormalTok{ calcular\_AP(verdades, confiancas, limiares):}
\NormalTok{    m }\OperatorTok{=} \BuiltInTok{len}\NormalTok{(limiares)}
\NormalTok{    precisoes }\OperatorTok{=}\NormalTok{ [}\FloatTok{0.0}\NormalTok{] }\OperatorTok{*}\NormalTok{ m}
\NormalTok{    sensibilidades }\OperatorTok{=}\NormalTok{ [}\FloatTok{0.0}\NormalTok{] }\OperatorTok{*}\NormalTok{ m}
    \ControlFlowTok{for}\NormalTok{ i }\KeywordTok{in} \BuiltInTok{range}\NormalTok{(m):}
\NormalTok{        p, s }\OperatorTok{=}\NormalTok{ calcular\_metricas(verdades, confiancas, limiares[i])}
\NormalTok{        precisoes[m}\OperatorTok{{-}}\DecValTok{1}\OperatorTok{{-}}\NormalTok{i] }\OperatorTok{=}\NormalTok{ p}
\NormalTok{        sensibilidades[m}\OperatorTok{{-}}\DecValTok{1}\OperatorTok{{-}}\NormalTok{i] }\OperatorTok{=}\NormalTok{ s}
\NormalTok{    prec\_mono }\OperatorTok{=}\NormalTok{ precisoes.copy()}
    \ControlFlowTok{for}\NormalTok{ i }\KeywordTok{in} \BuiltInTok{range}\NormalTok{(m}\OperatorTok{{-}}\DecValTok{2}\NormalTok{, }\OperatorTok{{-}}\DecValTok{1}\NormalTok{, }\OperatorTok{{-}}\DecValTok{1}\NormalTok{):}
        \ControlFlowTok{if}\NormalTok{ prec\_mono[i] }\OperatorTok{\textless{}}\NormalTok{ prec\_mono[i}\OperatorTok{+}\DecValTok{1}\NormalTok{]:}
\NormalTok{            prec\_mono[i] }\OperatorTok{=}\NormalTok{ prec\_mono[i}\OperatorTok{+}\DecValTok{1}\NormalTok{]}
\NormalTok{    AP }\OperatorTok{=} \FloatTok{0.0}
    \ControlFlowTok{for}\NormalTok{ i }\KeywordTok{in} \BuiltInTok{range}\NormalTok{(}\DecValTok{1}\NormalTok{, m):}
        \CommentTok{\# Regra do trapézio: média das alturas vezes a base}
\NormalTok{        area\_trapezio }\OperatorTok{=}\NormalTok{ (prec\_mono[i}\OperatorTok{{-}}\DecValTok{1}\NormalTok{] }\OperatorTok{+}\NormalTok{ prec\_mono[i]) }\OperatorTok{/} \FloatTok{2.0}
\NormalTok{        AP }\OperatorTok{+=}\NormalTok{ area\_trapezio }\OperatorTok{*}\NormalTok{ (sensibilidades[i] }\OperatorTok{{-}}\NormalTok{ sensibilidades[i}\OperatorTok{{-}}\DecValTok{1}\NormalTok{])}
    \ControlFlowTok{return}\NormalTok{ precisoes, prec\_mono, sensibilidades, AP}
\end{Highlighting}
\end{Shaded}

\begin{Shaded}
\begin{Highlighting}[]
\OperatorTok{\%\%}\NormalTok{writefile EP01\_03.py}
\CommentTok{\# sua solução}
\end{Highlighting}
\end{Shaded}

\begin{verbatim}
Overwriting EP01_03.py
\end{verbatim}

\begin{Shaded}
\begin{Highlighting}[]
\NormalTok{TestSuite(}\StringTok{"EP01\_03.py"}\NormalTok{).run()}
\end{Highlighting}
\end{Shaded}

\begin{verbatim}
<IPython.core.display.HTML object>
\end{verbatim}

\begin{center}\rule{0.5\linewidth}{0.5pt}\end{center}

\subsection{\texorpdfstring{EP01\_04 \emoji{framed-picture} Leitura e
Informações de uma Imagem em
Matriz}{EP01\_04  Leitura e Informações de uma Imagem em Matriz}}\label{ep01_04-leitura-e-informauxe7uxf5es-de-uma-imagem-em-matriz}

Nesta atividade, você deve escrever um programa que processe uma imagem
digital representada como uma matriz de pixels em tons de cinza.

\begin{itemize}
\tightlist
\item
  Leia dois inteiros \textbf{L} e \textbf{C}, representando o número de
  linhas e colunas.
\item
  Leia os \textbf{L * C} valores inteiros que compõem a matriz da imagem
  (cada valor entre 0 e 255).
\item
  Calcule e imprima as seguintes informações:
\end{itemize}

\begin{enumerate}
\def\labelenumi{\arabic{enumi}.}
\tightlist
\item
  O número de linhas.
\item
  O número de colunas.
\item
  O valor do maior pixel (\textbf{Máximo}).
\item
  O valor do menor pixel (\textbf{Mínimo}).
\item
  A \textbf{Média} aritmética de todos os pixels.
\end{enumerate}

\emoji{pushpin} \textbf{Importante}:

\begin{itemize}
\tightlist
\item
  A saída deve seguir exatamente o formato rotulado (ex:
  \texttt{Linhas:\ X}).
\item
  O valor da média deve ser formatado com \textbf{duas casas decimais}.
\item
  Ver um simulador interativo para esta questão na
  \textbf{Figura~\ref{fig-01-estatisticas-pixels}} (grade interativa
  para visualização de intensidades e cálculos em tempo real).
\end{itemize}

\begin{center}\rule{0.5\linewidth}{0.5pt}\end{center}

\subsubsection{\texorpdfstring{\emoji{brain} Por que isso importa? -- A
Imagem como
Dados}{ Por que isso importa? -- A Imagem como Dados}}\label{por-que-isso-importa-a-imagem-como-dados}

Toda imagem digital é, no fundo, uma estrutura de dados. Em tons de
cinza de 8 bits, cada pixel é um valor escalar. Extrair estatísticas
básicas é o primeiro passo para:

\begin{longtable}[]{@{}
  >{\raggedright\arraybackslash}p{(\linewidth - 2\tabcolsep) * \real{0.5000}}
  >{\raggedright\arraybackslash}p{(\linewidth - 2\tabcolsep) * \real{0.5000}}@{}}
\toprule\noalign{}
\begin{minipage}[b]{\linewidth}\raggedright
Operação
\end{minipage} & \begin{minipage}[b]{\linewidth}\raggedright
Utilidade Prática
\end{minipage} \\
\midrule\noalign{}
\endhead
\bottomrule\noalign{}
\endlastfoot
\textbf{Máximo/Mínimo} & Identificar se a imagem está ``lavada'' (baixo
contraste) ou saturada. \\
\textbf{Média} & Calcular o brilho global da cena para ajustes de
exposição. \\
\textbf{Normalização} & Redimensionar os valores para intervalos como
\([0, 1]\) em redes neurais. \\
\end{longtable}

\begin{center}\rule{0.5\linewidth}{0.5pt}\end{center}

\subsubsection{\texorpdfstring{\emoji{clipboard} Tarefa (especificação
para
VPL)}{ Tarefa (especificação para VPL)}}\label{tarefa-especificauxe7uxe3o-para-vpl-1}

\textbf{Entrada:}

A primeira linha contém o inteiro \texttt{L} (linhas).

A segunda linha contém o inteiro \texttt{C} (colunas).

As linhas seguintes contêm os elementos da matriz.

\textbf{Saída:}

Cinco linhas formatadas conforme o exemplo:

\texttt{Linhas:\ L}

\texttt{Colunas:\ C}

\texttt{Max:\ V}

\texttt{Min:\ V}

\texttt{Media:\ V.VV}

\begin{center}\rule{0.5\linewidth}{0.5pt}\end{center}

\subsubsection{\texorpdfstring{\emoji{pushpin}
Exemplos}{ Exemplos}}\label{exemplos-1}

\begin{longtable}[]{@{}
  >{\raggedright\arraybackslash}p{(\linewidth - 4\tabcolsep) * \real{0.3333}}
  >{\raggedright\arraybackslash}p{(\linewidth - 4\tabcolsep) * \real{0.3333}}
  >{\raggedright\arraybackslash}p{(\linewidth - 4\tabcolsep) * \real{0.3333}}@{}}
\toprule\noalign{}
\begin{minipage}[b]{\linewidth}\raggedright
Entrada
\end{minipage} & \begin{minipage}[b]{\linewidth}\raggedright
Saída
\end{minipage} & \begin{minipage}[b]{\linewidth}\raggedright
Observação
\end{minipage} \\
\midrule\noalign{}
\endhead
\bottomrule\noalign{}
\endlastfoot
230 128 25550 100 200 & Linhas: 2Colunas: 3Max: 255Min: 0Media: 122.17 &
Imagem pequena de alto contraste \\
\end{longtable}

\begin{figure}

\centering{

\begin{verbatim}
<IPython.core.display.HTML object>
\end{verbatim}

}

\caption{\label{fig-01-estatisticas-pixels}Simulador: Estatísticas de
Pixels}

\end{figure}%

\begin{Shaded}
\begin{Highlighting}[]
\OperatorTok{\%\%}\NormalTok{writefile EP01\_04.py}
\CommentTok{\# sua solução}
\end{Highlighting}
\end{Shaded}

\begin{verbatim}
Writing EP01_04.py
\end{verbatim}

\begin{Shaded}
\begin{Highlighting}[]
\NormalTok{TestSuite(}\StringTok{"EP01\_04.py"}\NormalTok{).run()}
\end{Highlighting}
\end{Shaded}

\begin{verbatim}
<IPython.core.display.HTML object>
\end{verbatim}

\begin{center}\rule{0.5\linewidth}{0.5pt}\end{center}

\subsection{\texorpdfstring{EP01\_05
\emoji{counterclockwise-arrows-button} Negativo de uma Imagem em Tons de
Cinza}{EP01\_05  Negativo de uma Imagem em Tons de Cinza}}\label{ep01_05-negativo-de-uma-imagem-em-tons-de-cinza}

Nesta atividade, você deve escrever um programa que calcule o negativo
de uma imagem digital.

\begin{itemize}
\tightlist
\item
  Leia dois inteiros \textbf{L} e \textbf{C}, representando o número de
  linhas e colunas.
\item
  Leia os valores inteiros que compõem a matriz da imagem.
\item
  Para cada pixel, aplique a transformação de inversão:
\end{itemize}

\[pixel_{negativo} = 255 - pixel_{original}\]

\begin{itemize}
\tightlist
\item
  Imprima a matriz resultante, mantendo o formato original (L linhas e C
  colunas).
\end{itemize}

\emoji{pushpin} \textbf{Importante}:

\begin{itemize}
\tightlist
\item
  Os valores de cada linha na saída devem ser separados por um
  \textbf{espaço em branco}.
\item
  A saída deve conter \textbf{apenas os números} da matriz resultante.
\item
  Ver um simulador interativo para esta questão na
  \textbf{Figura~\ref{fig-01-negativo-imagem}} (comparação em tempo real
  entre a matriz original e seu negativo).
\end{itemize}

\begin{center}\rule{0.5\linewidth}{0.5pt}\end{center}

\subsubsection{\texorpdfstring{\emoji{brain} Por que isso importa? --
Inversão de
Intensidade}{ Por que isso importa? -- Inversão de Intensidade}}\label{por-que-isso-importa-inversuxe3o-de-intensidade}

O \textbf{negativo} é uma transformação linear básica que inverte a
escala de brilho. É uma ferramenta essencial para o olho humano
identificar detalhes claros que estão ``escondidos'' em fundos muito
escuros, sendo amplamente utilizada em:

\begin{longtable}[]{@{}
  >{\raggedright\arraybackslash}p{(\linewidth - 2\tabcolsep) * \real{0.5000}}
  >{\raggedright\arraybackslash}p{(\linewidth - 2\tabcolsep) * \real{0.5000}}@{}}
\toprule\noalign{}
\begin{minipage}[b]{\linewidth}\raggedright
Aplicação
\end{minipage} & \begin{minipage}[b]{\linewidth}\raggedright
Utilidade
\end{minipage} \\
\midrule\noalign{}
\endhead
\bottomrule\noalign{}
\endlastfoot
\textbf{Imagens Médicas} & Melhora a visualização de anomalias em
tecidos densos (ex: Raios-X). \\
\textbf{Astronomia} & Destacar galáxias e nebulosas tênues contra o
vazio do espaço. \\
\textbf{Artes Digitais} & Efeitos estéticos e preparação de máscaras de
seleção. \\
\end{longtable}

\begin{center}\rule{0.5\linewidth}{0.5pt}\end{center}

\subsubsection{\texorpdfstring{\emoji{clipboard} Tarefa (especificação
para
VPL)}{ Tarefa (especificação para VPL)}}\label{tarefa-especificauxe7uxe3o-para-vpl-2}

\textbf{Entrada:}

A primeira linha contém o inteiro \texttt{L}.

A segunda linha contém o inteiro \texttt{C}.

As linhas seguintes contêm os elementos da matriz.

\textbf{Saída:}

A matriz invertida com \texttt{L} linhas e \texttt{C} colunas.

\begin{center}\rule{0.5\linewidth}{0.5pt}\end{center}

\subsubsection{\texorpdfstring{\emoji{pushpin}
Exemplos}{ Exemplos}}\label{exemplos-2}

\begin{longtable}[]{@{}
  >{\raggedright\arraybackslash}p{(\linewidth - 4\tabcolsep) * \real{0.3333}}
  >{\raggedright\arraybackslash}p{(\linewidth - 4\tabcolsep) * \real{0.3333}}
  >{\raggedright\arraybackslash}p{(\linewidth - 4\tabcolsep) * \real{0.3333}}@{}}
\toprule\noalign{}
\begin{minipage}[b]{\linewidth}\raggedright
Entrada
\end{minipage} & \begin{minipage}[b]{\linewidth}\raggedright
Saída
\end{minipage} & \begin{minipage}[b]{\linewidth}\raggedright
Observação
\end{minipage} \\
\midrule\noalign{}
\endhead
\bottomrule\noalign{}
\endlastfoot
230 128 25550 100 255 & 255 127 0205 155 55 & Onde era 0 (preto) vira
255 (branco) \\
\end{longtable}

\begin{figure}

\centering{

\begin{verbatim}
<IPython.core.display.HTML object>
\end{verbatim}

}

\caption{\label{fig-01-negativo-imagem}Simulador: Transformação de
Negativo}

\end{figure}%

\begin{Shaded}
\begin{Highlighting}[]
\OperatorTok{\%\%}\NormalTok{writefile EP01\_05.py}
\CommentTok{\# sua solução}
\end{Highlighting}
\end{Shaded}

\begin{verbatim}
Writing EP01_05.py
\end{verbatim}

\begin{Shaded}
\begin{Highlighting}[]
\NormalTok{TestSuite(}\StringTok{"EP01\_05.py"}\NormalTok{).run()}
\end{Highlighting}
\end{Shaded}

\begin{verbatim}
<IPython.core.display.HTML object>
\end{verbatim}

\begin{center}\rule{0.5\linewidth}{0.5pt}\end{center}

\subsection{\texorpdfstring{EP01\_06 \emoji{artist-palette} Conversão
RGB → Tons de Cinza (ITU-R
BT.601)}{EP01\_06  Conversão RGB → Tons de Cinza (ITU-R BT.601)}}\label{ep01_06-conversuxe3o-rgb-tons-de-cinza-itu-r-bt.601}

Nesta atividade, você deve escrever um programa que converta pixels
coloridos (RGB) para tons de cinza utilizando a ponderação fisiológica
do padrão ITU-R BT.601.

\begin{itemize}
\tightlist
\item
  Leia dois inteiros \textbf{L} e \textbf{C}, representando o número de
  linhas e colunas.
\item
  Leia \textbf{L × C} triplas de inteiros, onde cada tripla representa
  os canais \textbf{R} (Red), \textbf{G} (Green) e \textbf{B} (Blue) de
  um pixel.
\item
  Para cada pixel, calcule o valor de cinza (\(g\)) usando a fórmula:
\end{itemize}

\[g = \text{round}(0.299 \times R + 0.587 \times G + 0.114 \times B)\]

\begin{itemize}
\tightlist
\item
  Imprima a matriz resultante (L linhas e C colunas) contendo os valores
  inteiros convertidos.
\end{itemize}

\emoji{pushpin} \textbf{Importante}:

\begin{itemize}
\tightlist
\item
  Utilize a função \texttt{round()} da sua linguagem para garantir o
  arredondamento correto para o inteiro mais próximo.
\item
  A saída deve conter apenas os valores de cinza, mantendo a estrutura
  de matriz (separados por espaço na linha).
\item
  Ver um simulador interativo para esta questão na
  \textbf{Figura~\ref{fig-01-rgb-gray}} (ajuste os sliders para ver como
  cada cor contribui para o brilho final).
\end{itemize}

\begin{center}\rule{0.5\linewidth}{0.5pt}\end{center}

\subsubsection{\texorpdfstring{\emoji{brain} Por que não usar apenas a
média?}{ Por que não usar apenas a média?}}\label{por-que-nuxe3o-usar-apenas-a-muxe9dia}

O olho humano não percebe todas as cores com a mesma intensidade. Somos
muito mais sensíveis ao \textbf{Verde} do que ao \textbf{Azul} devido à
nossa evolução biológica. O padrão ITU-R BT.601 utiliza pesos
específicos para criar uma imagem em cinza que pareça naturalmente
correta para nossa visão:

\begin{longtable}[]{@{}
  >{\raggedright\arraybackslash}p{(\linewidth - 4\tabcolsep) * \real{0.3333}}
  >{\raggedright\arraybackslash}p{(\linewidth - 4\tabcolsep) * \real{0.3333}}
  >{\raggedright\arraybackslash}p{(\linewidth - 4\tabcolsep) * \real{0.3333}}@{}}
\toprule\noalign{}
\begin{minipage}[b]{\linewidth}\raggedright
Canal
\end{minipage} & \begin{minipage}[b]{\linewidth}\raggedright
Peso
\end{minipage} & \begin{minipage}[b]{\linewidth}\raggedright
Percepção Humana
\end{minipage} \\
\midrule\noalign{}
\endhead
\bottomrule\noalign{}
\endlastfoot
\emoji{green-circle} \textbf{Verde} & \textbf{58.7\%} & Máxima
sensibilidade (distinção de folhagens). \\
\emoji{red-circle} \textbf{Vermelho} & \textbf{29.9\%} & Sensibilidade
média. \\
\emoji{blue-circle} \textbf{Azul} & \textbf{11.4\%} & Baixa
sensibilidade (tons mais escuros). \\
\end{longtable}

\begin{center}\rule{0.5\linewidth}{0.5pt}\end{center}

\subsubsection{\texorpdfstring{\emoji{clipboard} Tarefa (especificação
para
VPL)}{ Tarefa (especificação para VPL)}}\label{tarefa-especificauxe7uxe3o-para-vpl-3}

\textbf{Entrada:}

A primeira linha contém o inteiro \texttt{L}.

A segunda linha contém o inteiro \texttt{C}.

As linhas seguintes contêm triplas de inteiros \texttt{R\ G\ B} para
cada pixel.

\textbf{Saída:}

A matriz de cinzas com \texttt{L} linhas e \texttt{C} colunas.

\begin{center}\rule{0.5\linewidth}{0.5pt}\end{center}

\subsubsection{\texorpdfstring{\emoji{pushpin}
Exemplos}{ Exemplos}}\label{exemplos-3}

\begin{longtable}[]{@{}
  >{\raggedright\arraybackslash}p{(\linewidth - 4\tabcolsep) * \real{0.3333}}
  >{\raggedright\arraybackslash}p{(\linewidth - 4\tabcolsep) * \real{0.3333}}
  >{\raggedright\arraybackslash}p{(\linewidth - 4\tabcolsep) * \real{0.3333}}@{}}
\toprule\noalign{}
\begin{minipage}[b]{\linewidth}\raggedright
Entrada
\end{minipage} & \begin{minipage}[b]{\linewidth}\raggedright
Saída
\end{minipage} & \begin{minipage}[b]{\linewidth}\raggedright
Observação
\end{minipage} \\
\midrule\noalign{}
\endhead
\bottomrule\noalign{}
\endlastfoot
13255 0 00 255 00 0 255 & 76 150 29 & Note como o Verde (150) é mais
brilhante que o Azul (29) \\
\end{longtable}

\begin{figure}

\centering{

\begin{verbatim}
<IPython.core.display.HTML object>
\end{verbatim}

}

\caption{\label{fig-01-rgb-gray}Simulador: Percepção de Cor e Pesos
ITU-R}

\end{figure}%

\begin{Shaded}
\begin{Highlighting}[]
\OperatorTok{\%\%}\NormalTok{writefile EP01\_06.py}
\CommentTok{\# sua solução}
\end{Highlighting}
\end{Shaded}

\begin{verbatim}
Writing EP01_06.py
\end{verbatim}

\begin{Shaded}
\begin{Highlighting}[]
\NormalTok{TestSuite(}\StringTok{"EP01\_06.py"}\NormalTok{).run()}
\end{Highlighting}
\end{Shaded}

\begin{verbatim}
<IPython.core.display.HTML object>
\end{verbatim}

\begin{center}\rule{0.5\linewidth}{0.5pt}\end{center}

\subsection{\texorpdfstring{EP01\_07 \emoji{black-circle} Limiarização
Manual: Imagem
Binária}{EP01\_07  Limiarização Manual: Imagem Binária}}\label{ep01_07-limiarizauxe7uxe3o-manual-imagem-binuxe1ria}

Nesta atividade, você deve escrever um programa que realize a
segmentação de uma imagem através da limiarização (\emph{thresholding}).

\begin{itemize}
\tightlist
\item
  Leia dois inteiros \textbf{L} e \textbf{C}, representando as dimensões
  da matriz.
\item
  Leia um inteiro \textbf{T}, que será o valor do limiar (corte).
\item
  Leia os valores inteiros da matriz.
\item
  Para cada pixel \(p\), aplique a seguinte regra de binarização:
\end{itemize}

\[\text{resultado} = \begin{cases} 255 & \text{se } p > T \\ 0 & \text{se } p \le T \end{cases}\]

\begin{itemize}
\tightlist
\item
  Imprima a matriz resultante contendo apenas os valores 0 ou 255.
\end{itemize}

\emoji{pushpin} \textbf{Importante}:

\begin{itemize}
\tightlist
\item
  Preste atenção ao operador: o pixel só vira branco (255) se for
  \textbf{estritamente maior} que \(T\).
\item
  A saída deve manter a estrutura de matriz (L linhas e C colunas).
\item
  Ver um simulador interativo para esta questão na
  \textbf{Figura~\ref{fig-01-limiarizacao}} (ajuste o slider de \(T\)
  para observar como os objetos são isolados do fundo).
\end{itemize}

\begin{center}\rule{0.5\linewidth}{0.5pt}\end{center}

\subsubsection{\texorpdfstring{\emoji{brain} O que é
Segmentação?}{ O que é Segmentação?}}\label{o-que-uxe9-segmentauxe7uxe3o}

A limiarização é o método mais simples de separar objetos de interesse
do fundo da imagem. Ao transformar tons de cinza em preto e branco puro,
criamos um mapa binário que facilita a contagem de objetos ou a
identificação de formas:

\begin{longtable}[]{@{}lll@{}}
\toprule\noalign{}
Valor do Pixel (\(p\)) & Condição & Resultado Final \\
\midrule\noalign{}
\endhead
\bottomrule\noalign{}
\endlastfoot
Escuro (\(p \le T\)) & Fundo/Ruído & \textbf{0} (Preto) \\
Claro (\(p > T\)) & Objeto/Destaque & \textbf{255} (Branco) \\
\end{longtable}

\begin{center}\rule{0.5\linewidth}{0.5pt}\end{center}

\subsubsection{\texorpdfstring{\emoji{clipboard} Tarefa (especificação
para
VPL)}{ Tarefa (especificação para VPL)}}\label{tarefa-especificauxe7uxe3o-para-vpl-4}

\textbf{Entrada:}

A primeira linha contém o inteiro \texttt{L}.

A segunda linha contém o inteiro \texttt{C}.

A terceira linha contém o inteiro \texttt{T} (limiar).

As linhas seguintes contêm os elementos da matriz.

\textbf{Saída:}

A matriz binarizada (0 ou 255) com \texttt{L} linhas e \texttt{C}
colunas.

\begin{center}\rule{0.5\linewidth}{0.5pt}\end{center}

\subsubsection{\texorpdfstring{\emoji{pushpin}
Exemplos}{ Exemplos}}\label{exemplos-4}

\begin{longtable}[]{@{}
  >{\raggedright\arraybackslash}p{(\linewidth - 4\tabcolsep) * \real{0.3333}}
  >{\raggedright\arraybackslash}p{(\linewidth - 4\tabcolsep) * \real{0.3333}}
  >{\raggedright\arraybackslash}p{(\linewidth - 4\tabcolsep) * \real{0.3333}}@{}}
\toprule\noalign{}
\begin{minipage}[b]{\linewidth}\raggedright
Entrada
\end{minipage} & \begin{minipage}[b]{\linewidth}\raggedright
Saída
\end{minipage} & \begin{minipage}[b]{\linewidth}\raggedright
Observação
\end{minipage} \\
\midrule\noalign{}
\endhead
\bottomrule\noalign{}
\endlastfoot
241280 100 128 20050 129 255 64 & 0 0 0 2550 255 255 0 & Note que o
valor 128 virou 0 (pois \(128 \le 128\)) \\
\end{longtable}

\begin{figure}

\centering{

\begin{verbatim}
<IPython.core.display.HTML object>
\end{verbatim}

}

\caption{\label{fig-01-limiarizacao}Simulador: Limiarização Interativa}

\end{figure}%

\begin{Shaded}
\begin{Highlighting}[]
\OperatorTok{\%\%}\NormalTok{writefile EP01\_07.py}
\CommentTok{\# sua solução}
\end{Highlighting}
\end{Shaded}

\begin{verbatim}
Writing EP01_07.py
\end{verbatim}

\begin{Shaded}
\begin{Highlighting}[]
\NormalTok{TestSuite(}\StringTok{"EP01\_07.py"}\NormalTok{).run()}
\end{Highlighting}
\end{Shaded}

\begin{verbatim}
<IPython.core.display.HTML object>
\end{verbatim}

\begin{center}\rule{0.5\linewidth}{0.5pt}\end{center}

\subsection{\texorpdfstring{EP01\_08 $\odot$ Ajuste de Brilho e
Contraste}{EP01\_08  Ajuste de Brilho e Contraste}}\label{ep01_08-ajuste-de-brilho-e-contraste}

Nesta atividade, você deve implementar uma transformação linear de
intensidade para ajustar o brilho e o contraste de uma imagem.

\begin{itemize}
\tightlist
\item
  Leia dois inteiros \textbf{L} e \textbf{C}, representando as dimensões
  da matriz.
\item
  Leia um valor real \textbf{\(\alpha\)} (contraste) e um inteiro
  \textbf{\(\beta\)} (brilho).
\item
  Leia os valores inteiros da matriz.
\item
  Para cada pixel \(p\), calcule o novo valor \(p'\) usando a fórmula:
\end{itemize}

\[p' = \text{clip}(\text{round}(\alpha \times p + \beta))\]

\begin{itemize}
\tightlist
\item
  Imprima a matriz resultante com os valores finais.
\end{itemize}

\emoji{pushpin} \textbf{Importante}:

\begin{itemize}
\tightlist
\item
  \textbf{Saturação (Clip):} O resultado final deve ser obrigatoriamente
  um inteiro no intervalo \([0, 255]\). Valores menores que 0 tornam-se
  0, e valores maiores que 255 tornam-se 255.
\item
  \textbf{Arredondamento:} Realize o arredondamento matemático antes de
  converter para inteiro e aplicar o clip.
\item
  Ver um simulador interativo para esta questão na
  \textbf{Figura~\ref{fig-01-brilho-contraste}} (manipule os parâmetros
  \(\alpha\) e \(\beta\) para visualizar o efeito de realce e correção).
\end{itemize}

\begin{center}\rule{0.5\linewidth}{0.5pt}\end{center}

\subsubsection{\texorpdfstring{\emoji{brain} Transformação Linear de
Intensidade}{ Transformação Linear de Intensidade}}\label{transformauxe7uxe3o-linear-de-intensidade}

Ajustar o brilho e o contraste são operações fundamentais em PDI. Elas
permitem corrigir imagens subexpostas (escuras demais) ou melhorar a
visibilidade de detalhes sutis:

\begin{longtable}[]{@{}
  >{\raggedright\arraybackslash}p{(\linewidth - 4\tabcolsep) * \real{0.3333}}
  >{\raggedright\arraybackslash}p{(\linewidth - 4\tabcolsep) * \real{0.3333}}
  >{\raggedright\arraybackslash}p{(\linewidth - 4\tabcolsep) * \real{0.3333}}@{}}
\toprule\noalign{}
\begin{minipage}[b]{\linewidth}\raggedright
Parâmetro
\end{minipage} & \begin{minipage}[b]{\linewidth}\raggedright
Função
\end{minipage} & \begin{minipage}[b]{\linewidth}\raggedright
Efeito
\end{minipage} \\
\midrule\noalign{}
\endhead
\bottomrule\noalign{}
\endlastfoot
\textbf{\(\alpha\) (Alpha)} & Multiplicador & Ajusta o
\textbf{Contraste}. \(>1\) aumenta a diferença entre tons; \(<1\)
diminui. \\
\textbf{\(\beta\) (Beta)} & Aditivo & Ajusta o \textbf{Brilho}. Valores
positivos clareiam; negativos escurecem. \\
\textbf{Clip} & Limitador & Garante que o pixel permaneça dentro do
padrão de 8 bits (\(0-255\)). \\
\end{longtable}

\begin{center}\rule{0.5\linewidth}{0.5pt}\end{center}

\subsubsection{\texorpdfstring{\emoji{clipboard} Tarefa (especificação
para
VPL)}{ Tarefa (especificação para VPL)}}\label{tarefa-especificauxe7uxe3o-para-vpl-5}

\textbf{Entrada:}

A primeira linha contém \texttt{L}.

A segunda linha contém \texttt{C}.

A terceira linha contém \texttt{alpha} (\(\alpha\)) e \texttt{beta}
(\(\beta\)).

As linhas seguintes contêm os elementos da matriz.

\textbf{Saída:}

A matriz transformada com \texttt{L} linhas e \texttt{C} colunas.

\begin{center}\rule{0.5\linewidth}{0.5pt}\end{center}

\subsubsection{\texorpdfstring{\emoji{pushpin}
Exemplos}{ Exemplos}}\label{exemplos-5}

\begin{longtable}[]{@{}
  >{\raggedright\arraybackslash}p{(\linewidth - 4\tabcolsep) * \real{0.3333}}
  >{\raggedright\arraybackslash}p{(\linewidth - 4\tabcolsep) * \real{0.3333}}
  >{\raggedright\arraybackslash}p{(\linewidth - 4\tabcolsep) * \real{0.3333}}@{}}
\toprule\noalign{}
\begin{minipage}[b]{\linewidth}\raggedright
Entrada
\end{minipage} & \begin{minipage}[b]{\linewidth}\raggedright
Saída
\end{minipage} & \begin{minipage}[b]{\linewidth}\raggedright
Observação
\end{minipage} \\
\midrule\noalign{}
\endhead
\bottomrule\noalign{}
\endlastfoot
141.5 -300 100 180 255 & 0 120 240 255 & Note o efeito de saturação no
último pixel \\
\end{longtable}

\begin{figure}

\centering{

\begin{verbatim}
<IPython.core.display.HTML object>
\end{verbatim}

}

\caption{\label{fig-01-brilho-contraste}Simulador: Ajuste de Brilho e
Contraste}

\end{figure}%

\begin{Shaded}
\begin{Highlighting}[]
\OperatorTok{\%\%}\NormalTok{writefile EP01\_08.py}
\CommentTok{\# sua solução}
\end{Highlighting}
\end{Shaded}

\begin{verbatim}
Writing EP01_08.py
\end{verbatim}

\begin{Shaded}
\begin{Highlighting}[]
\NormalTok{TestSuite(}\StringTok{"EP01\_08.py"}\NormalTok{).run()}
\end{Highlighting}
\end{Shaded}

\begin{verbatim}
<IPython.core.display.HTML object>
\end{verbatim}

\begin{center}\rule{0.5\linewidth}{0.5pt}\end{center}

\subsection{\texorpdfstring{EP01\_09 \emoji{chequered-flag} Padrão de
Tabuleiro: Matriz de
Xadrez}{EP01\_09  Padrão de Tabuleiro: Matriz de Xadrez}}\label{ep01_09-padruxe3o-de-tabuleiro-matriz-de-xadrez}

Nesta atividade, você deve escrever um programa que gere uma imagem
sintética no padrão de um tabuleiro de xadrez.

\begin{itemize}
\tightlist
\item
  Leia dois inteiros \textbf{L} (linhas) e \textbf{C} (colunas).
\item
  Gere uma matriz onde os valores alternam entre \textbf{0} (preto) e
  \textbf{1} (branco).
\item
  A lógica de preenchimento deve seguir a regra da paridade:
\item
  O elemento na posição \((0,0)\) é sempre \textbf{0}.
\item
  Um pixel na posição \((i, j)\) será \textbf{1} se a soma dos índices
  \((i + j)\) for \textbf{ímpar}.
\item
  Um pixel na posição \((i, j)\) será \textbf{0} se a soma dos índices
  \((i + j)\) for \textbf{par}.
\end{itemize}

\emoji{pushpin} \textbf{Importante}:

\begin{itemize}
\tightlist
\item
  As cores devem alternar corretamente tanto horizontalmente quanto
  verticalmente.
\item
  A saída deve ser a matriz impressa linha por linha, com os elementos
  separados por um espaço.
\item
  Ver um simulador interativo para esta questão na
  \textbf{Figura~\ref{fig-01-gerador-xadrez}} (ajuste as dimensões para
  visualizar a construção da malha e a saída textual correspondente).
\end{itemize}

\begin{center}\rule{0.5\linewidth}{0.5pt}\end{center}

\subsubsection{\texorpdfstring{\emoji{brain} Padrões
Sintéticos}{ Padrões Sintéticos}}\label{padruxf5es-sintuxe9ticos}

Criar padrões geométricos é um exercício fundamental para dominar a
\textbf{lógica de índices} em matrizes. Em Processamento Digital de
Imagens, o padrão de xadrez não é apenas estético; ele é amplamente
utilizado para:

\begin{longtable}[]{@{}
  >{\raggedright\arraybackslash}p{(\linewidth - 2\tabcolsep) * \real{0.5000}}
  >{\raggedright\arraybackslash}p{(\linewidth - 2\tabcolsep) * \real{0.5000}}@{}}
\toprule\noalign{}
\begin{minipage}[b]{\linewidth}\raggedright
Aplicação
\end{minipage} & \begin{minipage}[b]{\linewidth}\raggedright
Utilidade
\end{minipage} \\
\midrule\noalign{}
\endhead
\bottomrule\noalign{}
\endlastfoot
\textbf{Calibração de Câmera} & Estimar parâmetros intrínsecos e
extrínsecos da lente. \\
\textbf{Correção de Distorção} & Identificar e corrigir o efeito de
``barril'' ou ``almofada'' em lentes grande-angulares. \\
\textbf{Mapeamento 3D} & Projetar padrões conhecidos para reconstruir
superfícies em sistemas de luz estruturada. \\
\end{longtable}

\begin{center}\rule{0.5\linewidth}{0.5pt}\end{center}

\subsubsection{\texorpdfstring{\emoji{clipboard} Tarefa (especificação
para
VPL)}{ Tarefa (especificação para VPL)}}\label{tarefa-especificauxe7uxe3o-para-vpl-6}

\textbf{Entrada:}

Uma linha contendo o inteiro \texttt{L} (linhas).

Uma linha contendo o inteiro \texttt{C} (colunas).

\textbf{Saída:}

A matriz de xadrez com \texttt{L} linhas e \texttt{C} colunas, impressa
com espaços entre os elementos.

\begin{center}\rule{0.5\linewidth}{0.5pt}\end{center}

\subsubsection{\texorpdfstring{\emoji{pushpin}
Exemplos}{ Exemplos}}\label{exemplos-6}

\begin{longtable}[]{@{}
  >{\raggedright\arraybackslash}p{(\linewidth - 4\tabcolsep) * \real{0.3333}}
  >{\raggedright\arraybackslash}p{(\linewidth - 4\tabcolsep) * \real{0.3333}}
  >{\raggedright\arraybackslash}p{(\linewidth - 4\tabcolsep) * \real{0.3333}}@{}}
\toprule\noalign{}
\begin{minipage}[b]{\linewidth}\raggedright
Entrada
\end{minipage} & \begin{minipage}[b]{\linewidth}\raggedright
Saída
\end{minipage} & \begin{minipage}[b]{\linewidth}\raggedright
Observação
\end{minipage} \\
\midrule\noalign{}
\endhead
\bottomrule\noalign{}
\endlastfoot
34 & 0 1 0 11 0 1 00 1 0 1 & Note que cada linha começa com o inverso da
anterior \\
\end{longtable}

\begin{figure}

\centering{

\begin{verbatim}
<IPython.core.display.HTML object>
\end{verbatim}

}

\caption{\label{fig-01-gerador-xadrez}Simulador: Gerador de Xadrez}

\end{figure}%

\begin{Shaded}
\begin{Highlighting}[]
\OperatorTok{\%\%}\NormalTok{writefile EP01\_09.py}
\CommentTok{\# sua solução}
\end{Highlighting}
\end{Shaded}

\begin{verbatim}
Writing EP01_09.py
\end{verbatim}

\begin{Shaded}
\begin{Highlighting}[]
\NormalTok{TestSuite(}\StringTok{"EP01\_09.py"}\NormalTok{).run()}
\end{Highlighting}
\end{Shaded}

\begin{verbatim}
<IPython.core.display.HTML object>
\end{verbatim}

\begin{center}\rule{0.5\linewidth}{0.5pt}\end{center}

\subsection{\texorpdfstring{EP01\_10 \emoji{page-facing-up} Metadados:
Leitura de Arquivo
PGM}{EP01\_10  Metadados: Leitura de Arquivo PGM}}\label{ep01_10-metadados-leitura-de-arquivo-pgm}

Nesta atividade, você deve ler um arquivo de imagem no formato
\textbf{PGM (Portable Gray Map)} e extrair suas dimensões a partir do
cabeçalho.

\begin{itemize}
\tightlist
\item
  O formato \textbf{PGM (P2)} é um arquivo de texto simples (ASCII) que
  armazena imagens em tons de cinza.
\item
  O arquivo possui um \textbf{cabeçalho} estruturado da seguinte forma:
\end{itemize}

\begin{enumerate}
\def\labelenumi{\arabic{enumi}.}
\tightlist
\item
  \textbf{Versão:} O identificador \texttt{P2}.
\item
  \textbf{Comentários:} Linhas opcionais que começam com \texttt{\#}
  (devem ser ignoradas).
\item
  \textbf{Dimensões:} Dois inteiros representando \textbf{Largura} e
  \textbf{Altura}.
\item
  \textbf{Máximo:} Um inteiro representando a intensidade máxima
  (geralmente 255).
\end{enumerate}

\begin{itemize}
\tightlist
\item
  Após o cabeçalho, seguem-se os dados dos pixels.
\end{itemize}

\emoji{pushpin} \textbf{Importante}:

\begin{itemize}
\tightlist
\item
  \textbf{Leitura de Arquivo:} Você deve abrir o arquivo indicado no
  exemplo utilizando a função \texttt{open()} do Python.
\item
  \textbf{Ordem de Saída:} Diferente do arquivo, a saída esperada deve
  estar no formato de tupla: \texttt{(Altura,\ Largura,\ Canais)}.
\item
  Como arquivos PGM são em tons de cinza, o número de \textbf{Canais} é
  sempre \textbf{1}.
\item
  Ver um simulador interativo para esta questão na
  \textbf{Figura~\ref{fig-01-estrutura-pgm}} (ajuste as dimensões para
  ver como o cabeçalho ASCII é gerado).
\end{itemize}

\begin{center}\rule{0.5\linewidth}{0.5pt}\end{center}

\subsubsection{\texorpdfstring{\emoji{brain} Entendendo o formato
PGM}{ Entendendo o formato PGM}}\label{entendendo-o-formato-pgm}

O formato PGM é um dos mais simples para processamento de imagens. Por
ser em texto puro, ele permite visualizar os metadados e até os valores
dos pixels abrindo o arquivo em um bloco de notas:

\begin{longtable}[]{@{}
  >{\raggedright\arraybackslash}p{(\linewidth - 4\tabcolsep) * \real{0.3333}}
  >{\raggedright\arraybackslash}p{(\linewidth - 4\tabcolsep) * \real{0.3333}}
  >{\raggedright\arraybackslash}p{(\linewidth - 4\tabcolsep) * \real{0.3333}}@{}}
\toprule\noalign{}
\begin{minipage}[b]{\linewidth}\raggedright
Componente
\end{minipage} & \begin{minipage}[b]{\linewidth}\raggedright
Exemplo
\end{minipage} & \begin{minipage}[b]{\linewidth}\raggedright
Significado
\end{minipage} \\
\midrule\noalign{}
\endhead
\bottomrule\noalign{}
\endlastfoot
\textbf{Magic Number} & \texttt{P2} & Identifica que é um PGM em formato
de texto (ASCII). \\
\textbf{Comentário} & \texttt{\#\ CREATOR...} & Linha informativa
ignorada pelo processador. \\
\textbf{Dimensões} & \texttt{397\ 343} & \textbf{397} colunas (Largura)
e \textbf{343} linhas (Altura). \\
\textbf{Intensidade} & \texttt{255} & Define o valor do branco puro
(escala de 0 a 255). \\
\end{longtable}

\begin{center}\rule{0.5\linewidth}{0.5pt}\end{center}

\subsubsection{\texorpdfstring{\emoji{clipboard} Tarefa (especificação
para
VPL)}{ Tarefa (especificação para VPL)}}\label{tarefa-especificauxe7uxe3o-para-vpl-7}

\textbf{Entrada:}

Nenhuma entrada via teclado. O programa deve ler o arquivo
\texttt{"aula01fig03b.pgm"} presente no diretório de execução.

\textbf{Saída:}

Uma tupla contendo \texttt{(Altura,\ Largura,\ 1)}.

\begin{center}\rule{0.5\linewidth}{0.5pt}\end{center}

\subsubsection{\texorpdfstring{\emoji{pushpin}
Exemplos}{ Exemplos}}\label{exemplos-7}

\begin{longtable}[]{@{}
  >{\raggedright\arraybackslash}p{(\linewidth - 4\tabcolsep) * \real{0.3333}}
  >{\raggedright\arraybackslash}p{(\linewidth - 4\tabcolsep) * \real{0.3333}}
  >{\raggedright\arraybackslash}p{(\linewidth - 4\tabcolsep) * \real{0.3333}}@{}}
\toprule\noalign{}
\begin{minipage}[b]{\linewidth}\raggedright
Nome do Arquivo
\end{minipage} & \begin{minipage}[b]{\linewidth}\raggedright
Saída Esperada
\end{minipage} & \begin{minipage}[b]{\linewidth}\raggedright
Observação
\end{minipage} \\
\midrule\noalign{}
\endhead
\bottomrule\noalign{}
\endlastfoot
``aula01fig03b.pgm'' & (343, 397, 1) & Note a inversão da ordem: Altura
primeiro \\
\end{longtable}

\begin{figure}

\centering{

\begin{verbatim}
<IPython.core.display.HTML object>
\end{verbatim}

}

\caption{\label{fig-01-estrutura-pgm}Simulador: Estrutura do Arquivo
PGM}

\end{figure}%

\begin{tcolorbox}[enhanced jigsaw, titlerule=0mm, left=2mm, opacityback=0, toprule=.15mm, arc=.35mm, opacitybacktitle=0.6, colbacktitle=quarto-callout-note-color!10!white, colframe=quarto-callout-note-color-frame, rightrule=.15mm, breakable, coltitle=black, bottomtitle=1mm, toptitle=1mm, title=\textcolor{quarto-callout-note-color}{\faInfo}\hspace{0.5em}{Nota}, colback=white, bottomrule=.15mm, leftrule=.75mm]

O arquivo necessário para este EP será baixado automaticamente do
repositório por meio do seguinte código: ::

\begin{Shaded}
\begin{Highlighting}[]
\ImportTok{import}\NormalTok{ os, urllib.request}

\NormalTok{BASE\_URL }\OperatorTok{=} \StringTok{"https://raw.githubusercontent.com/fzampirolli/pdi{-}vc/master/all/cap01/imagens"}
\BuiltInTok{file} \OperatorTok{=} \StringTok{"aula01fig03b.pgm"}

\NormalTok{opener }\OperatorTok{=}\NormalTok{ urllib.request.build\_opener()}
\NormalTok{opener.addheaders }\OperatorTok{=}\NormalTok{ [(}\StringTok{\textquotesingle{}User{-}agent\textquotesingle{}}\NormalTok{, }\StringTok{\textquotesingle{}Mozilla/5.0\textquotesingle{}}\NormalTok{)]}
\NormalTok{urllib.request.install\_opener(opener)}

\ControlFlowTok{for}\NormalTok{ f }\KeywordTok{in}\NormalTok{ [}\BuiltInTok{file}\NormalTok{]:}
    \ControlFlowTok{if} \KeywordTok{not}\NormalTok{ os.path.exists(f):}
\NormalTok{        url }\OperatorTok{=} \SpecialStringTok{f"}\SpecialCharTok{\{}\NormalTok{BASE\_URL}\SpecialCharTok{\}}\SpecialStringTok{/}\SpecialCharTok{\{}\NormalTok{f}\SpecialCharTok{\}}\SpecialStringTok{"}
        \ControlFlowTok{try}\NormalTok{:}
\NormalTok{            urllib.request.urlretrieve(url, f)}
            \BuiltInTok{print}\NormalTok{(}\SpecialStringTok{f"✅ Arquivo baixado: }\SpecialCharTok{\{}\NormalTok{f}\SpecialCharTok{\}}\SpecialStringTok{"}\NormalTok{)}
        \ControlFlowTok{except}\NormalTok{ urllib.error.HTTPError }\ImportTok{as}\NormalTok{ e:}
            \BuiltInTok{print}\NormalTok{(}\SpecialStringTok{f"❌ Erro }\SpecialCharTok{\{}\NormalTok{e}\SpecialCharTok{.}\NormalTok{code}\SpecialCharTok{\}}\SpecialStringTok{: Não encontrado em }\SpecialCharTok{\{}\NormalTok{url}\SpecialCharTok{\}}\SpecialStringTok{"}\NormalTok{)}
\end{Highlighting}
\end{Shaded}

\begin{verbatim}
✅ Arquivo baixado: aula01fig03b.pgm
\end{verbatim}

\begin{Shaded}
\begin{Highlighting}[]
\OperatorTok{\%\%}\NormalTok{writefile EP01\_10.py}
\CommentTok{\# sua solução}
\end{Highlighting}
\end{Shaded}

\begin{verbatim}
Writing EP01_10.py
\end{verbatim}

\begin{Shaded}
\begin{Highlighting}[]
\NormalTok{TestSuite(}\StringTok{"EP01\_10.py"}\NormalTok{).run()}
\end{Highlighting}
\end{Shaded}

\begin{verbatim}
<IPython.core.display.HTML object>
\end{verbatim}

\begin{center}\rule{0.5\linewidth}{0.5pt}\end{center}

\subsection{\texorpdfstring{EP01\_11 $\nearrow$ Análise de Vizinhança:
Filtro de Máximo
1D}{EP01\_11  Análise de Vizinhança: Filtro de Máximo 1D}}\label{ep01_11-anuxe1lise-de-vizinhanuxe7a-filtro-de-muxe1ximo-1d}

Nesta atividade, você deve implementar um filtro morfológico simples de
máximo operando em um sinal unidimensional (vetor).

\begin{itemize}
\tightlist
\item
  Leia um inteiro \textbf{n}, representando o tamanho do vetor.
\item
  Leia os \textbf{n} elementos inteiros que compõem o vetor original
  \textbf{v1}.
\item
  Crie um novo vetor \textbf{v2}, onde cada posição \(i\) é o resultado
  da comparação entre o elemento atual e seus vizinhos imediatos:
\end{itemize}

\[v2[i] = \max(v1[i-1],\; v1[i],\; v1[i+1])\]

\emoji{pushpin} \textbf{Importante}:

\begin{itemize}
\tightlist
\item
  \textbf{Bordas:} Nas extremidades do vetor (índices \(0\) e \(n-1\)),
  a vizinhança é menor. No índice \(0\), compare apenas \(v1[0]\) e
  \(v1[1]\). No último índice, compare apenas \(v1[n-2]\) e \(v1[n-1]\).
\item
  \textbf{Saída:} Imprima o cabeçalho ``v2:'' seguido pelos valores do
  vetor resultante, um por linha.
\item
  Ver um simulador interativo para esta questão na
  \textbf{Figura~\ref{fig-01-filtro-maximo-1d}} (passe o mouse sobre os
  resultados para visualizar a janela de vizinhança utilizada no
  cálculo).
\end{itemize}

\begin{center}\rule{0.5\linewidth}{0.5pt}\end{center}

\subsubsection{\texorpdfstring{\emoji{brain} Por que analisar
vizinhos?}{ Por que analisar vizinhos?}}\label{por-que-analisar-vizinhos}

Em processamento de imagens, o valor de um pixel raramente é isolado;
ele depende do contexto ao seu redor. O \textbf{Filtro de Máximo} é a
base da operação de \textbf{Dilatação} em morfologia matemática,
servindo para:

\begin{longtable}[]{@{}
  >{\raggedright\arraybackslash}p{(\linewidth - 2\tabcolsep) * \real{0.5000}}
  >{\raggedright\arraybackslash}p{(\linewidth - 2\tabcolsep) * \real{0.5000}}@{}}
\toprule\noalign{}
\begin{minipage}[b]{\linewidth}\raggedright
Função
\end{minipage} & \begin{minipage}[b]{\linewidth}\raggedright
Efeito Visual
\end{minipage} \\
\midrule\noalign{}
\endhead
\bottomrule\noalign{}
\endlastfoot
\textbf{Realce} & Expande estruturas brilhantes e ``engorda'' objetos
claros. \\
\textbf{Remoção de Ruído} & Elimina pequenos pontos pretos (ruído de
``sal e pimenta'' escuro). \\
\textbf{Preenchimento} & Fecha pequenos buracos ou lacunas em formas
binárias. \\
\end{longtable}

\begin{center}\rule{0.5\linewidth}{0.5pt}\end{center}

\subsubsection{\texorpdfstring{\emoji{clipboard} Tarefa (especificação
para
VPL)}{ Tarefa (especificação para VPL)}}\label{tarefa-especificauxe7uxe3o-para-vpl-8}

\textbf{Entrada:}

Um inteiro \texttt{n}.

Nas linhas seguintes, os \texttt{n} elementos inteiros do vetor.

\textbf{Saída:}

A string \texttt{v2:} na primeira linha.

Nas linhas seguintes, cada elemento de \texttt{v2} (um por linha).

\begin{center}\rule{0.5\linewidth}{0.5pt}\end{center}

\subsubsection{\texorpdfstring{\emoji{pushpin}
Exemplos}{ Exemplos}}\label{exemplos-8}

\begin{longtable}[]{@{}
  >{\raggedright\arraybackslash}p{(\linewidth - 4\tabcolsep) * \real{0.3333}}
  >{\raggedright\arraybackslash}p{(\linewidth - 4\tabcolsep) * \real{0.3333}}
  >{\raggedright\arraybackslash}p{(\linewidth - 4\tabcolsep) * \real{0.3333}}@{}}
\toprule\noalign{}
\begin{minipage}[b]{\linewidth}\raggedright
Entrada
\end{minipage} & \begin{minipage}[b]{\linewidth}\raggedright
Saída
\end{minipage} & \begin{minipage}[b]{\linewidth}\raggedright
Observação
\end{minipage} \\
\midrule\noalign{}
\endhead
\bottomrule\noalign{}
\endlastfoot
5102053015 & v2:2020303030 & No índice 1: max(10, 20, 5) = 20 \\
\end{longtable}

\begin{figure}[H]

\centering{

\begin{verbatim}
<IPython.core.display.HTML object>
\end{verbatim}

}

\caption{\label{fig-01-filtro-maximo-1d}Simulador: Filtro de Máximo
Local 1D}

\end{figure}%

\begin{Shaded}
\begin{Highlighting}[]
\OperatorTok{\%\%}\NormalTok{writefile EP01\_11.py}
\CommentTok{\# sua solução}
\end{Highlighting}
\end{Shaded}

\begin{verbatim}
Writing EP01_11.py
\end{verbatim}

\begin{Shaded}
\begin{Highlighting}[]
\NormalTok{TestSuite(}\StringTok{"EP01\_11.py"}\NormalTok{).run()}
\end{Highlighting}
\end{Shaded}

\begin{verbatim}
<IPython.core.display.HTML object>
\end{verbatim}

\chapter{Da Captura ao Pixel - Amostragem, Quantização e
Conectividade}\label{da-captura-ao-pixel---amostragem-quantizauxe7uxe3o-e-conectividade}

\href{https://colab.research.google.com/github/fzampirolli/pdi-vc/blob/master/notebooks_alunos/cap02/cap02_aluno.ipynb}{
\includegraphics[width=0.16\linewidth,height=\textheight,keepaspectratio]{cap02/imagens/colab-badge.png}}
\href{https://github.com/fzampirolli/pdi-vc}{
\includegraphics[width=0.2\linewidth,height=\textheight,keepaspectratio]{cap02/imagens/github-badge.png}}

\section{2.1. Objetivos}\label{objetivos-1}

Ao final deste capítulo, você será capaz de:

\begin{itemize}
\tightlist
\item
  Explicar o modelo simples de formação da imagem baseado em iluminação
  e refletância.
\item
  Diferenciar os principais elementos da percepção visual humana e dos
  sensores digitais.
\item
  Compreender os processos de \textbf{amostragem} (discretização do
  espaço) e \textbf{quantização} (discretização da intensidade).
\item
  Descrever as relações topológicas entre pixels: vizinhança,
  adjacência, conectividade e distâncias.
\item
  Realizar transformações geométricas básicas (translação, rotação,
  escala) preservando a qualidade da imagem.
\item
  Visualizar na prática os efeitos da variação da resolução espacial e
  da profundidade de bits.
\end{itemize}

\begin{center}\rule{0.5\linewidth}{0.5pt}\end{center}

\section{2.2. O Olho e a Câmera -- Elementos da Percepção
Visual}\label{o-olho-e-a-cuxe2mera-elementos-da-percepuxe7uxe3o-visual}

A formação de uma imagem digital começa com a captura da luz refletida
pelos objetos. Tanto o olho humano quanto as câmeras digitais seguem
princípios semelhantes, mas com diferenças importantes.

\subsection{2.2.1. Aspectos da visão
humana}\label{aspectos-da-visuxe3o-humana}

O olho é um sistema óptico complexo. A luz passa pela córnea, pelo humor
aquoso, pelo cristalino (que ajusta o foco) e pelo humor vítreo, até
atingir a retina. As células fotorreceptoras (cones e bastonetes)
convertem a luz em impulsos nervosos. Os cones são responsáveis pela
visão de cores (três tipos sensíveis ao vermelho, verde e azul) e os
bastonetes pela visão em baixa luminosidade (apenas intensidade). O
ponto cego é a região de saída do nervo óptico, sem fotorreceptores.

\subsection{2.2.2. Sensores digitais}\label{sensores-digitais}

As câmeras digitais utilizam sensores CCD ou CMOS que substituem a
retina. Cada elemento do sensor (fotossítio) acumula carga elétrica
proporcional à luz incidente. Um filtro de Bayer (matriz de filtros
coloridos) permite capturar as três componentes de cor. O processo de
leitura analógico-digital converte a carga em um valor numérico
(quantização). Diferentemente do olho, o sensor digital tem resolução
espacial fixa (número de pixels) e profundidade de bits definida (ex.: 8
bits por canal).

\begin{quote}
\textbf{Curiosidade:} O olho humano possui cerca de 120 milhões de
bastonetes e 6 milhões de cones, mas a resolução efetiva da visão
central (fóvea) é equivalente a aproximadamente 120\,×\,120 pixels de
alta definição. O processamento cerebral realiza enorme
pós-processamento.
\end{quote}

\begin{center}\rule{0.5\linewidth}{0.5pt}\end{center}

\section{2.3. O Modelo Matemático da Formação da
Imagem}\label{o-modelo-matemuxe1tico-da-formauxe7uxe3o-da-imagem}

Uma imagem pode ser modelada como o produto de duas funções:

{[} f(x,y) = i(x,y) \cdot r(x,y) {]}

onde:

\begin{itemize}
\tightlist
\item
  (i(x,y)) é a \textbf{iluminação} incidente sobre a cena (energia
  luminosa por unidade de área). Depende da fonte de luz.
\item
  (r(x,y)) é a \textbf{refletância} do objeto (fração da luz refletida
  em cada ponto). Depende das propriedades do material e da superfície.
\end{itemize}

Na prática, as duas componentes variam em faixas distintas: (i(x,y))
varia lentamente no espaço, enquanto (r(x,y)) pode variar rapidamente
(texturas, bordas). O processamento de imagens frequentemente tenta
separar ou compensar essas componentes (ex.: correção de iluminação não
uniforme).

\begin{center}\rule{0.5\linewidth}{0.5pt}\end{center}

\section{2.4. Digitalização: Amostragem e
Quantização}\label{digitalizauxe7uxe3o-amostragem-e-quantizauxe7uxe3o}

Para transformar uma cena contínua em uma imagem digital, dois processos
são necessários:

\subsection{2.4.1. Amostragem -- Discretização do
espaço}\label{amostragem-discretizauxe7uxe3o-do-espauxe7o}

A amostragem consiste em medir o valor da função (f(x,y)) em pontos
igualmente espaçados, formando uma matriz de (M) linhas e (N) colunas.
Cada elemento dessa matriz é um \textbf{pixel}. A \textbf{resolução
espacial} é dada por (M \times N). Quanto maior a resolução, mais
detalhes espaciais são preservados, mas maior também o custo
computacional e de armazenamento.

\textbf{Exemplo prático:} Uma imagem de 640\,×\,480 pixels tem resolução
espacial de 307\,200 pixels. Reduzir a resolução para 320\,×\,240
(subamostragem) reduz os detalhes e pode causar o efeito de
\emph{aliasing} (serrilhado). O aumento da resolução (sobre‐amostragem)
só é possível por interpolação, que não adiciona informação real.

\subsection{2.4.2. Quantização -- Discretização da
intensidade}\label{quantizauxe7uxe3o-discretizauxe7uxe3o-da-intensidade}

A quantização associa a cada pixel um valor numérico discreto,
geralmente representado por um inteiro de (b) bits. A
\textbf{profundidade de bits} ((b)) define o número de níveis de
intensidade: (2\^{}b) níveis. Imagens em tons de cinza costumam usar 8
bits (256 níveis). Imagens coloridas usam três canais de 8 bits (24 bits
no total).

\textbf{Ilustração:} Se usarmos apenas 1 bit por pixel (preto e branco),
perdemos todos os tons intermediários. Com 2 bits (4 níveis) já se
percebe degradês grosseiros. Com 8 bits, o olho humano dificilmente
percebe a discretização (visão contínua).

\begin{quote}
\textbf{Erro de quantização} é a diferença entre o valor analógico real
e o valor discreto atribuído. Ele se manifesta como ruído de
quantização, visível em regiões com gradiente suave quando se usa poucos
bits (efeito de ``posterização'').
\end{quote}

\begin{center}\rule{0.5\linewidth}{0.5pt}\end{center}

\section{2.5. Relações entre Pixels -- Topologia da
Imagem}\label{relauxe7uxf5es-entre-pixels-topologia-da-imagem}

Os pixels não são elementos isolados; suas posições relativas definem
importantes conceitos para processamento.

\subsection{2.5.1. Vizinhança}\label{vizinhanuxe7a}

Dado um pixel de coordenadas ((x,y)), definem‑se dois tipos principais
de vizinhança (para imagens em grid retangular):

\begin{itemize}
\tightlist
\item
  \textbf{Vizinhança‑4} (vizinhança de von Neumann): inclui os pixels
  nas posições ((x-1,y)), ((x+1,y)), ((x,y-1)), ((x,y+1)).
\item
  \textbf{Vizinhança‑8} (vizinhança de Moore): inclui todos os oito
  pixels adjacentes (acrescenta os quatro diagonais).
\end{itemize}

A escolha da vizinhança influencia operações como detecção de bordas,
cálculo de gradientes e conectividade.

\subsection{2.5.2. Adjacência, conectividade e
caminhos}\label{adjacuxeancia-conectividade-e-caminhos}

Dois pixels são \textbf{adjacentes} se estão em contato segundo uma
vizinhança definida e satisfazem um certo critério de valor (ex.: mesmo
nível de intensidade). Uma \textbf{conectividade} define uma relação de
equivalência entre pixels que formam uma região conexa. Um
\textbf{caminho} é uma sequência de pixels adjacentes. A conectividade‑4
e conectividade‑8 podem produzir resultados diferentes na segmentação e
no cálculo de componentes conexas (etiquetagem).

\subsection{2.5.3. Distâncias entre
pixels}\label{distuxe2ncias-entre-pixels}

Várias métricas podem ser usadas para medir distância:

Essas distâncias são usadas em algoritmos de interpolação, transformada
de distância, crescimento de regiões, etc.

\begin{center}\rule{0.5\linewidth}{0.5pt}\end{center}

\section{2.6. Armazenamento de Imagens}\label{armazenamento-de-imagens}

Uma imagem digital pode ser armazenada em diversos formatos, cada um com
características próprias:

\begin{longtable}[]{@{}
  >{\raggedright\arraybackslash}p{(\linewidth - 4\tabcolsep) * \real{0.2368}}
  >{\raggedright\arraybackslash}p{(\linewidth - 4\tabcolsep) * \real{0.4211}}
  >{\raggedright\arraybackslash}p{(\linewidth - 4\tabcolsep) * \real{0.3421}}@{}}
\toprule\noalign{}
\begin{minipage}[b]{\linewidth}\raggedright
Formato
\end{minipage} & \begin{minipage}[b]{\linewidth}\raggedright
Características
\end{minipage} & \begin{minipage}[b]{\linewidth}\raggedright
Uso típico
\end{minipage} \\
\midrule\noalign{}
\endhead
\bottomrule\noalign{}
\endlastfoot
\textbf{BMP} & Não comprimido (ou compressão simples). & Windows,
aplicações legado. \\
\textbf{PNG} & Compressão sem perdas (lossless). & Web, imagens com
transparência. \\
\textbf{JPEG} & Compressão com perdas (lossy), boa para fotografias. &
Fotos, imagens médicas (uso moderado). \\
\textbf{TIFF} & Suporta múltiplas camadas e compressão variada. &
Editoração, arquivamento. \\
\textbf{RAW} & Dados brutos do sensor, sem processamento. & Fotografia
profissional. \\
\end{longtable}

Os metadados de uma imagem incluem: largura, altura, profundidade de
bits, tipo de codificação de cor (grayscale, RGB, CMYK), informações de
calibração (ex.: resolução DPI), e metadados EXIF (data, modelo da
câmera, configurações). Ao ler uma imagem com a função
\texttt{mm.read()}, a biblioteca \texttt{morph} preserva automaticamente
os metadados relevantes.

\begin{center}\rule{0.5\linewidth}{0.5pt}\end{center}

\section{2.7. Transformações Geométricas
Básicas}\label{transformauxe7uxf5es-geomuxe9tricas-buxe1sicas}

As transformações geométricas alteram a posição dos pixels, mantendo os
valores de intensidade. São fundamentais para alinhamento, correção de
distorções e aumento de dados.

\subsection{2.7.1. Translação}\label{translauxe7uxe3o}

Desloca a imagem em (t\_x) e (t\_y) pixels:

Na prática, cria‑se uma nova imagem maior e copia‑se os pixels
deslocados. Pixels que saem da área original são descartados e áreas
vazias são preenchidas (geralmente com 0 -- preto).

\subsection{2.7.2. Rotação}\label{rotauxe7uxe3o}

Gira a imagem em torno da origem (ou de seu centro) por um ângulo
(\theta). A transformação direta pode produzir buracos; na prática,
usa‑se a transformação inversa (mapeamento de cada pixel da imagem
destino para a origem) e interpolação (vizinho mais próximo, bilinear,
bicúbica). A rotação geralmente requer redimensionamento da imagem para
que nenhuma área seja cortada.

\subsection{2.7.3. Escala
(redimensionamento)}\label{escala-redimensionamento}

Altera o tamanho da imagem pelo fator (s\_x) (horizontal) e (s\_y)
(vertical). O redimensionamento exige interpolação:

\begin{itemize}
\tightlist
\item
  \textbf{Vizinho mais próximo:} rápido, mas produz efeito de blocos.
\item
  \textbf{Bilinear:} suaviza melhor, é um bom compromisso.
\item
  \textbf{Bicúbica:} mais suave, usada em softwares profissionais.
\end{itemize}

A biblioteca \texttt{morph} oferece funções como
\texttt{mm.resize(img,\ scale,\ interpolation)} e
\texttt{mm.rotate(img,\ angle,\ center,\ interp)}.

\begin{center}\rule{0.5\linewidth}{0.5pt}\end{center}

\section{2.8. Laboratório Prático: Visualizando Efeitos da Amostragem e
Quantização}\label{laboratuxf3rio-pruxe1tico-visualizando-efeitos-da-amostragem-e-quantizauxe7uxe3o}

Neste experimento, vamos explorar como a redução da resolução espacial e
da profundidade de bits afetam a qualidade visual.

\subsection{2.8.1. Código para subamostragem (redução de
resolução)}\label{cuxf3digo-para-subamostragem-reduuxe7uxe3o-de-resoluuxe7uxe3o}

\begin{Shaded}
\begin{Highlighting}[]
\ImportTok{import}\NormalTok{ numpy }\ImportTok{as}\NormalTok{ np}
\ImportTok{import}\NormalTok{ matplotlib.pyplot }\ImportTok{as}\NormalTok{ plt}
\ImportTok{import}\NormalTok{ morph}

\CommentTok{\# Carregar uma imagem (ex.: Lenna)}
\NormalTok{img }\OperatorTok{=}\NormalTok{ mm.read(}\StringTok{\textquotesingle{}https://upload.wikimedia.org/wikipedia/en/7/7d/Lenna\_\%28test\_image\%29.png\textquotesingle{}}\NormalTok{)}
\NormalTok{img\_gray }\OperatorTok{=}\NormalTok{ mm.rgb2gray(img)}

\CommentTok{\# Função para subamostrar por fator (ex.: 2, 4, 8)}
\KeywordTok{def}\NormalTok{ subsample(image, factor):}
\NormalTok{    h, w }\OperatorTok{=}\NormalTok{ image.shape}
\NormalTok{    new\_h, new\_w }\OperatorTok{=}\NormalTok{ h }\OperatorTok{//}\NormalTok{ factor, w }\OperatorTok{//}\NormalTok{ factor}
    \CommentTok{\# Criar imagem reduzida: pegar um pixel a cada \textquotesingle{}factor\textquotesingle{}}
\NormalTok{    reduced }\OperatorTok{=}\NormalTok{ image[::factor, ::factor]}
    \CommentTok{\# Redimensionar de volta ao tamanho original para comparação (interpolação vizinho mais próximo)}
    \ControlFlowTok{return}\NormalTok{ mm.resize(reduced, (h, w), method}\OperatorTok{=}\StringTok{\textquotesingle{}nearest\textquotesingle{}}\NormalTok{)}

\NormalTok{plt.figure(figsize}\OperatorTok{=}\NormalTok{(}\DecValTok{15}\NormalTok{, }\DecValTok{5}\NormalTok{))}
\NormalTok{factors }\OperatorTok{=}\NormalTok{ [}\DecValTok{1}\NormalTok{, }\DecValTok{2}\NormalTok{, }\DecValTok{4}\NormalTok{, }\DecValTok{8}\NormalTok{]}
\ControlFlowTok{for}\NormalTok{ i, f }\KeywordTok{in} \BuiltInTok{enumerate}\NormalTok{(factors):}
\NormalTok{    sub }\OperatorTok{=}\NormalTok{ subsample(img\_gray, f)}
\NormalTok{    plt.subplot(}\DecValTok{1}\NormalTok{, }\DecValTok{4}\NormalTok{, i}\OperatorTok{+}\DecValTok{1}\NormalTok{)}
\NormalTok{    mm.imshow(sub)}
\NormalTok{    plt.title(}\SpecialStringTok{f\textquotesingle{}Fator = }\SpecialCharTok{\{}\NormalTok{f}\SpecialCharTok{\}}\SpecialStringTok{\textquotesingle{}}\NormalTok{)}
\NormalTok{plt.show()}
\end{Highlighting}
\end{Shaded}

\subsection{2.8.2. Código para redução da profundidade de bits
(quantização)}\label{cuxf3digo-para-reduuxe7uxe3o-da-profundidade-de-bits-quantizauxe7uxe3o}

\begin{Shaded}
\begin{Highlighting}[]
\KeywordTok{def}\NormalTok{ quantize(image, bits):}
    \CommentTok{"""Reduz a imagem para 2\^{}bits níveis de cinza."""}
\NormalTok{    levels }\OperatorTok{=} \DecValTok{2} \OperatorTok{**}\NormalTok{ bits}
    \CommentTok{\# Normalizar, quantizar e denormalizar}
    \ControlFlowTok{return}\NormalTok{ np.floor(image }\OperatorTok{/} \DecValTok{256} \OperatorTok{*}\NormalTok{ levels) }\OperatorTok{/}\NormalTok{ levels }\OperatorTok{*} \DecValTok{255}

\NormalTok{plt.figure(figsize}\OperatorTok{=}\NormalTok{(}\DecValTok{15}\NormalTok{, }\DecValTok{5}\NormalTok{))}
\NormalTok{bits }\OperatorTok{=}\NormalTok{ [}\DecValTok{8}\NormalTok{, }\DecValTok{4}\NormalTok{, }\DecValTok{2}\NormalTok{, }\DecValTok{1}\NormalTok{]}
\ControlFlowTok{for}\NormalTok{ i, b }\KeywordTok{in} \BuiltInTok{enumerate}\NormalTok{(bits):}
\NormalTok{    quant }\OperatorTok{=}\NormalTok{ quantize(img\_gray, b)}
\NormalTok{    plt.subplot(}\DecValTok{1}\NormalTok{, }\DecValTok{4}\NormalTok{, i}\OperatorTok{+}\DecValTok{1}\NormalTok{)}
\NormalTok{    mm.imshow(quant)}
\NormalTok{    plt.title(}\SpecialStringTok{f\textquotesingle{}}\SpecialCharTok{\{}\NormalTok{b}\SpecialCharTok{\}}\SpecialStringTok{ bits (}\SpecialCharTok{\{}\DecValTok{2}\OperatorTok{**}\NormalTok{b}\SpecialCharTok{\}}\SpecialStringTok{ níveis)\textquotesingle{}}\NormalTok{)}
\NormalTok{plt.show()}
\end{Highlighting}
\end{Shaded}

\subsection{2.8.3. Discussão dos
resultados}\label{discussuxe3o-dos-resultados}

\begin{itemize}
\tightlist
\item
  \textbf{Subamostragem (fator 8):} observa‑se o efeito de
  \emph{aliasing} -- bordas serrilhadas e perda de detalhes finos. Em
  fator 4 já se nota blocos, mas a imagem ainda é reconhecível.
\item
  \textbf{Quantização com 1 bit:} a imagem vira preto e branco, perde‑se
  totalmente a graduação de tons. Com 2 bits (4 níveis) ainda se percebe
  um efeito de ``posters'', mas a silhueta geral se mantém. Com 4 bits
  (16 níveis) a qualidade é razoável; o olho humano dificilmente
  distingue 8 bits de 16 bits em imagens comuns.
\end{itemize}

Esses experimentos mostram o compromisso entre qualidade e custo de
armazenamento, essencial para sistemas de processamento e compressão.

\begin{center}\rule{0.5\linewidth}{0.5pt}\end{center}

\section{2.9. Exercícios Propostos}\label{exercuxedcios-propostos}

\begin{enumerate}
\def\labelenumi{\arabic{enumi}.}
\tightlist
\item
  \textbf{Amostragem} -- Carregue uma imagem, aplique subamostragem com
  fatores 2, 3 e 5. Compare visualmente e calcule o fator de redução no
  número de pixels.
\item
  \textbf{Quantização} -- Implemente uma função que reduza a imagem a 3
  bits (8 níveis). Visualize e descreva o ruído introduzido.
\item
  \textbf{Conectividade} -- Escreva um código que identifique
  componentes conexas em uma imagem binária usando vizinhança‑4 e
  vizinhança‑8. Explique as diferenças.
\item
  \textbf{Transformação geométrica} -- Aplique rotação de 30° e 45° em
  uma imagem, utilizando interpolação bilinear. Compare com a rotação
  sem interpolação (vizinho mais próximo) e comente sobre a suavidade.
\end{enumerate}

\begin{center}\rule{0.5\linewidth}{0.5pt}\end{center}

\section{2.10. Referências}\label{referuxeancias}

\begin{itemize}
\tightlist
\item
  GONZALEZ, R. C.; WOODS, R. E. \emph{Processamento Digital de Imagens}.
  3. ed.~Pearson, 2010.
\item
  SZELISKI, R. \emph{Computer Vision: Algorithms and Applications}.
  Springer, 2010.
\item
  Otsu, N. (1979). \emph{A threshold selection method from gray‑level
  histograms}. IEEE Trans. Syst. Man Cybern., 9(1), 62--66. (citado no
  Capítulo 1).
\end{itemize}

\begin{center}\rule{0.5\linewidth}{0.5pt}\end{center}

\textbf{Nota:} As operações de convolução, filtragem e morfologia
matemática (erosão, dilatação) serão tratadas no \textbf{Capítulo 3 --
Operações Espaciais em Imagens}. O presente capítulo concentrou‑se na
formação, digitalização e topologia básica da imagem.

\section{Objetivos}\label{objetivos-2}

Ao final deste capítulo, o aluno será capaz de:

\begin{itemize}
\tightlist
\item
  Compreender a representação matemática de imagens \(f(x,y)\)
\item
  Ler, exibir e salvar imagens
\item
  Trabalhar com imagens binárias, tons de cinza e coloridas (RGB)
\item
  Acessar e modificar pixels individualmente
\end{itemize}

\section{O que é uma Imagem
Digital?}\label{o-que-uxe9-uma-imagem-digital-1}

Uma imagem digital é uma função discreta:

\begin{equation}\phantomsection\label{eq-1-imagem}{f: \mathbb{Z}^2 \rightarrow \mathbb{Z}^k
}\end{equation}

onde \((x,y)\) são coordenadas espaciais e \(f(x,y)\) é a intensidade do
pixel. Para imagens em \textbf{tons de cinza} \(k=1\) com valores em
\([0, 255]\); para imagens \textbf{coloridas} (RGB) \(k=3\), ver
Figura~\ref{fig-1-representacao}.

\begin{figure}[H]

\centering{

\includegraphics[width=0.7\linewidth,height=\textheight,keepaspectratio]{cap02/imagens/lena.jpg}

}

\caption{\label{fig-1-representacao}Exemplo de Transformação Não Linear
do Espaço de Entrada para o Espaço de Características.}

\end{figure}%

\section{Configuração do
Ambiente}\label{configurauxe7uxe3o-do-ambiente-1}

\begin{Shaded}
\begin{Highlighting}[]
\CommentTok{\# Instalação (Google Colab — executar uma vez)}
\CommentTok{\# !pip install opencv{-}python{-}headless scikit{-}image matplotlib numpy}

\ImportTok{import}\NormalTok{ sys}
\ImportTok{import}\NormalTok{ importlib}
\ImportTok{import}\NormalTok{ numpy }\ImportTok{as}\NormalTok{ np}
\ImportTok{import}\NormalTok{ cv2}
\ImportTok{import}\NormalTok{ matplotlib.pyplot }\ImportTok{as}\NormalTok{ plt}

\NormalTok{sys.path.insert(}\DecValTok{0}\NormalTok{, }\StringTok{\textquotesingle{}../../morph\textquotesingle{}}\NormalTok{)}
\ImportTok{import}\NormalTok{ morph}
\NormalTok{importlib.}\BuiltInTok{reload}\NormalTok{(morph)}
\ImportTok{from}\NormalTok{ morph }\ImportTok{import}\NormalTok{ mm}

\BuiltInTok{print}\NormalTok{(}\SpecialStringTok{f\textquotesingle{}OpenCV }\SpecialCharTok{\{}\NormalTok{cv2}\SpecialCharTok{.}\NormalTok{\_\_version\_\_}\SpecialCharTok{\}}\SpecialStringTok{ | NumPy }\SpecialCharTok{\{}\NormalTok{np}\SpecialCharTok{.}\NormalTok{\_\_version\_\_}\SpecialCharTok{\}}\SpecialStringTok{\textquotesingle{}}\NormalTok{)}
\BuiltInTok{print}\NormalTok{(}\StringTok{"Carregando imagem..."}\NormalTok{)}
\BuiltInTok{print}\NormalTok{(}\BuiltInTok{help}\NormalTok{(mm.show))}
\end{Highlighting}
\end{Shaded}

\begin{verbatim}
OpenCV 4.12.0 | NumPy 2.2.6
Carregando imagem...
Help on function show in module morph:

show(*args, title=None)
    This function will draw images f
    input: <*args> set of images f_i, where i>0 is binary image
          <title> optional string title for the plot
    output: image drawing
    Example:
    f1, f2 = np.zeros((100, 100,3)),  np.zeros((100, 100))
    f2[50:60, 50:60] = 1
    mm.show(f1, f2, title='Exemplo')

None
\end{verbatim}

\section{1.3 Lendo e Exibindo Imagens}\label{lendo-e-exibindo-imagens-1}

A operação mais básica em PDI é a leitura de uma imagem. A função
\texttt{mm.read()} da \texttt{morph.py} aceita caminhos locais e URLs.

\begin{Shaded}
\begin{Highlighting}[]
\ImportTok{import}\NormalTok{ urllib.request, os}

\CommentTok{\# Baixar imagem de exemplo}
\NormalTok{url }\OperatorTok{=} \StringTok{\textquotesingle{}https://raw.githubusercontent.com/opencv/opencv/master/samples/data/lena.jpg\textquotesingle{}}
\NormalTok{os.makedirs(}\StringTok{\textquotesingle{}imagens\textquotesingle{}}\NormalTok{, exist\_ok}\OperatorTok{=}\VariableTok{True}\NormalTok{)}
\ControlFlowTok{if} \KeywordTok{not}\NormalTok{ os.path.exists(}\StringTok{\textquotesingle{}imagens/lena.jpg\textquotesingle{}}\NormalTok{):}
\NormalTok{    urllib.request.urlretrieve(url, }\StringTok{\textquotesingle{}imagens/lena.jpg\textquotesingle{}}\NormalTok{)}

\CommentTok{\# Ler imagem com morph.py}
\NormalTok{img }\OperatorTok{=}\NormalTok{ mm.read(}\StringTok{\textquotesingle{}imagens/lena.jpg\textquotesingle{}}\NormalTok{)}

\BuiltInTok{print}\NormalTok{(}\SpecialStringTok{f\textquotesingle{}Shape  : }\SpecialCharTok{\{}\NormalTok{img}\SpecialCharTok{.}\NormalTok{shape}\SpecialCharTok{\}}\SpecialStringTok{\textquotesingle{}}\NormalTok{)    }\CommentTok{\# (altura, largura, canais)}
\BuiltInTok{print}\NormalTok{(}\SpecialStringTok{f\textquotesingle{}Dtype  : }\SpecialCharTok{\{}\NormalTok{img}\SpecialCharTok{.}\NormalTok{dtype}\SpecialCharTok{\}}\SpecialStringTok{\textquotesingle{}}\NormalTok{)    }\CommentTok{\# uint8}
\BuiltInTok{print}\NormalTok{(}\SpecialStringTok{f\textquotesingle{}Min/Max: }\SpecialCharTok{\{}\NormalTok{img}\SpecialCharTok{.}\BuiltInTok{min}\NormalTok{()}\SpecialCharTok{\}}\SpecialStringTok{ / }\SpecialCharTok{\{}\NormalTok{img}\SpecialCharTok{.}\BuiltInTok{max}\NormalTok{()}\SpecialCharTok{\}}\SpecialStringTok{\textquotesingle{}}\NormalTok{)}

\NormalTok{img\_rgb }\OperatorTok{=}\NormalTok{ cv2.cvtColor(img, cv2.COLOR\_BGR2RGB)}
\NormalTok{mm.show(img\_rgb, title}\OperatorTok{=}\StringTok{\textquotesingle{}Imagem original\textquotesingle{}}\NormalTok{)}
\end{Highlighting}
\end{Shaded}

\begin{verbatim}
Shape  : (512, 512, 3)
Dtype  : uint8
Min/Max: 0 / 255
\end{verbatim}

\begin{figure}[H]

\centering{

\pandocbounded{\includegraphics[keepaspectratio]{cap02/cap02.py.pt_files/figure-pdf/fig-1-lena-output-2.png}}

}

\caption{\label{fig-1-lena}MVS Linear --- Margem Máxima entre Íris
Setosa e Versicolor (atributos de pétala).}

\end{figure}%

\section{1.4 Tipos de Imagem}\label{tipos-de-imagem}

\subsection{1.4.1 Conversão para Tons de
Cinza}\label{conversuxe3o-para-tons-de-cinza-1}

A conversão de RGB para tons de cinza aplica pesos perceptuais definidos
pela ITU-R BT.601:

\begin{equation}\phantomsection\label{eq-1-cinza}{g = 0.299\,R + 0.587\,G + 0.114\,B}\end{equation}

\begin{Shaded}
\begin{Highlighting}[]
\CommentTok{\# Converter para tons de cinza}
\NormalTok{img\_gray }\OperatorTok{=}\NormalTok{ mm.gray(img)}
\BuiltInTok{print}\NormalTok{(}\SpecialStringTok{f\textquotesingle{}Shape cinza: }\SpecialCharTok{\{}\NormalTok{img\_gray}\SpecialCharTok{.}\NormalTok{shape}\SpecialCharTok{\}}\SpecialStringTok{\textquotesingle{}}\NormalTok{)}

\NormalTok{mm.show([img, img\_gray], title}\OperatorTok{=}\NormalTok{[}\StringTok{\textquotesingle{}Original (RGB)\textquotesingle{}}\NormalTok{, }\StringTok{\textquotesingle{}Tons de Cinza\textquotesingle{}}\NormalTok{])}
\end{Highlighting}
\end{Shaded}

\subsection{1.4.2 Limiarização
(Binarização)}\label{limiarizauxe7uxe3o-binarizauxe7uxe3o}

A limiarização converte uma imagem em tons de cinza para binária:

\begin{equation}\phantomsection\label{eq-1-limiar}{g(x,y) = \begin{cases} 255 & \text{se } f(x,y) \geq T \\ 0 & \text{caso contrário} \end{cases}}\end{equation}

O método de \textbf{Otsu} (Otsu, 1979) determina \(T\) automaticamente
maximizando a variância inter-classes do histograma.

\begin{Shaded}
\begin{Highlighting}[]
\CommentTok{\# Limiarização com T fixo}
\NormalTok{T }\OperatorTok{=} \DecValTok{128}
\NormalTok{img\_bin }\OperatorTok{=}\NormalTok{ mm.threshold(img\_gray, T)}

\CommentTok{\# Limiarização Otsu (T automático)}
\NormalTok{\_, img\_otsu }\OperatorTok{=}\NormalTok{ cv2.threshold(img\_gray, }\DecValTok{0}\NormalTok{, }\DecValTok{255}\NormalTok{,}
\NormalTok{                             cv2.THRESH\_BINARY }\OperatorTok{+}\NormalTok{ cv2.THRESH\_OTSU)}

\BuiltInTok{print}\NormalTok{(}\SpecialStringTok{f\textquotesingle{}Valores únicos (T=}\SpecialCharTok{\{}\NormalTok{T}\SpecialCharTok{\}}\SpecialStringTok{): }\SpecialCharTok{\{}\BuiltInTok{sorted}\NormalTok{(}\BuiltInTok{set}\NormalTok{(img\_bin.flatten()))}\SpecialCharTok{\}}\SpecialStringTok{\textquotesingle{}}\NormalTok{)}

\NormalTok{mm.show([img\_gray, img\_bin, img\_otsu],}
\NormalTok{        title}\OperatorTok{=}\NormalTok{[}\SpecialStringTok{f\textquotesingle{}Cinza\textquotesingle{}}\NormalTok{, }\SpecialStringTok{f\textquotesingle{}Binária (T=}\SpecialCharTok{\{}\NormalTok{T}\SpecialCharTok{\}}\SpecialStringTok{)\textquotesingle{}}\NormalTok{, }\StringTok{\textquotesingle{}Otsu\textquotesingle{}}\NormalTok{])}
\end{Highlighting}
\end{Shaded}

\section{1.5 Acesso a Pixels}\label{acesso-a-pixels-1}

Em Python/NumPy, o acesso a pixels usa indexação de arrays
\texttt{img{[}linha,\ coluna{]}}.

\begin{Shaded}
\begin{Highlighting}[]
\NormalTok{r, c }\OperatorTok{=} \DecValTok{100}\NormalTok{, }\DecValTok{100}

\CommentTok{\# Acessar pixel}
\NormalTok{pixel\_cinza }\OperatorTok{=}\NormalTok{ img\_gray[r, c]}
\NormalTok{pixel\_rgb   }\OperatorTok{=}\NormalTok{ img[r, c]       }\CommentTok{\# array [R, G, B]}

\BuiltInTok{print}\NormalTok{(}\SpecialStringTok{f\textquotesingle{}Pixel cinza (}\SpecialCharTok{\{}\NormalTok{r}\SpecialCharTok{\}}\SpecialStringTok{,}\SpecialCharTok{\{}\NormalTok{c}\SpecialCharTok{\}}\SpecialStringTok{): }\SpecialCharTok{\{}\NormalTok{pixel\_cinza}\SpecialCharTok{\}}\SpecialStringTok{\textquotesingle{}}\NormalTok{)}
\BuiltInTok{print}\NormalTok{(}\SpecialStringTok{f\textquotesingle{}Pixel RGB   (}\SpecialCharTok{\{}\NormalTok{r}\SpecialCharTok{\}}\SpecialStringTok{,}\SpecialCharTok{\{}\NormalTok{c}\SpecialCharTok{\}}\SpecialStringTok{): R=}\SpecialCharTok{\{}\NormalTok{pixel\_rgb[}\DecValTok{0}\NormalTok{]}\SpecialCharTok{\}}\SpecialStringTok{  G=}\SpecialCharTok{\{}\NormalTok{pixel\_rgb[}\DecValTok{1}\NormalTok{]}\SpecialCharTok{\}}\SpecialStringTok{  B=}\SpecialCharTok{\{}\NormalTok{pixel\_rgb[}\DecValTok{2}\NormalTok{]}\SpecialCharTok{\}}\SpecialStringTok{\textquotesingle{}}\NormalTok{)}

\CommentTok{\# Imagem sintética 5×5 com pixel central branco}
\NormalTok{syn }\OperatorTok{=}\NormalTok{ np.zeros((}\DecValTok{5}\NormalTok{, }\DecValTok{5}\NormalTok{), dtype}\OperatorTok{=}\NormalTok{np.uint8)}
\NormalTok{syn[}\DecValTok{2}\NormalTok{, }\DecValTok{2}\NormalTok{] }\OperatorTok{=} \DecValTok{255}
\BuiltInTok{print}\NormalTok{(}\StringTok{\textquotesingle{}}\CharTok{\textbackslash{}n}\StringTok{Matriz 5×5:\textquotesingle{}}\NormalTok{)}
\BuiltInTok{print}\NormalTok{(syn)}
\end{Highlighting}
\end{Shaded}

\section{1.6 Resumo}\label{resumo-1}

Neste capítulo foram apresentados os fundamentos de representação de
imagens digitais:

\begin{itemize}
\tightlist
\item
  Imagem digital = função \(f(x,y)\) mapeando coordenadas para
  intensidades
\item
  Tipos: binária (\(k=1\), \(\{0,255\}\)), tons de cinza (\(k=1\),
  \([0,255]\)), RGB (\(k=3\))
\item
  \texttt{morph.py}: biblioteca didática que simplifica operações de PDI
  em Python
\item
  Limiarização: operação fundamental de binarização
\end{itemize}

O Capítulo 2 estuda \textbf{histogramas e equalização}.

\subsection{Questões de Revisão}\label{questuxf5es-de-revisuxe3o}

\begin{enumerate}
\def\labelenumi{\arabic{enumi}.}
\tightlist
\item
  Qual a diferença entre imagem em tons de cinza e imagem binária?
\item
  Como o método de Otsu determina o limiar \(T\) automaticamente?
\item
  O que representa \texttt{img.shape} para uma imagem NumPy RGB?
\end{enumerate}

\chapter{Operações Espaciais: Transformações Geométricas, Convolução e
Morfologia}\label{operauxe7uxf5es-espaciais-transformauxe7uxf5es-geomuxe9tricas-convoluuxe7uxe3o-e-morfologia}

\href{https://colab.research.google.com/github/fzampirolli/pdi-vc/blob/master/notebooks_alunos/cap03/cap03_aluno.ipynb}{
\includegraphics[width=0.16\linewidth,height=\textheight,keepaspectratio]{cap03/imagens/colab-badge.png}}
\href{https://github.com/fzampirolli/pdi-vc}{
\includegraphics[width=0.2\linewidth,height=\textheight,keepaspectratio]{cap03/imagens/github-badge.png}}

Objetivos: Aplicar transformações geométricas (rotação, escala,
translação); realizar operações de convolução e filtragem; introduzir
operações morfológicas (erosão, dilatação, abertura, fechamento).

Conteúdo:

Transformações geométricas (interpolação)

Convolução e filtros (média, gaussiano, bordas)

Morfologia matemática (erosão, dilatação, operações derivadas)

\section{Objetivos}\label{objetivos-3}

Ao final deste capítulo, o aluno será capaz de:

\begin{itemize}
\tightlist
\item
  Compreender a representação matemática de imagens \(f(x,y)\)
\item
  Ler, exibir e salvar imagens
\item
  Trabalhar com imagens binárias, tons de cinza e coloridas (RGB)
\item
  Acessar e modificar pixels individualmente
\end{itemize}

\section{O que é uma Imagem
Digital?}\label{o-que-uxe9-uma-imagem-digital-2}

Uma imagem digital é uma função discreta:

\begin{equation}\phantomsection\label{eq-1-imagem}{f: \mathbb{Z}^2 \rightarrow \mathbb{Z}^k
}\end{equation}

onde \((x,y)\) são coordenadas espaciais e \(f(x,y)\) é a intensidade do
pixel. Para imagens em \textbf{tons de cinza} \(k=1\) com valores em
\([0, 255]\); para imagens \textbf{coloridas} (RGB) \(k=3\), ver
Figura~\ref{fig-1-representacao}.

\begin{figure}[H]

\centering{

\includegraphics[width=0.7\linewidth,height=\textheight,keepaspectratio]{cap03/imagens/lena.jpg}

}

\caption{\label{fig-1-representacao}Exemplo de Transformação Não Linear
do Espaço de Entrada para o Espaço de Características.}

\end{figure}%

\section{Configuração do
Ambiente}\label{configurauxe7uxe3o-do-ambiente-2}

\begin{Shaded}
\begin{Highlighting}[]
\CommentTok{\# Instalação (Google Colab — executar uma vez)}
\CommentTok{\# !pip install opencv{-}python{-}headless scikit{-}image matplotlib numpy}

\ImportTok{import}\NormalTok{ sys}
\ImportTok{import}\NormalTok{ importlib}
\ImportTok{import}\NormalTok{ numpy }\ImportTok{as}\NormalTok{ np}
\ImportTok{import}\NormalTok{ cv2}
\ImportTok{import}\NormalTok{ matplotlib.pyplot }\ImportTok{as}\NormalTok{ plt}

\NormalTok{sys.path.insert(}\DecValTok{0}\NormalTok{, }\StringTok{\textquotesingle{}../../morph\textquotesingle{}}\NormalTok{)}
\ImportTok{import}\NormalTok{ morph}
\NormalTok{importlib.}\BuiltInTok{reload}\NormalTok{(morph)}
\ImportTok{from}\NormalTok{ morph }\ImportTok{import}\NormalTok{ mm}

\BuiltInTok{print}\NormalTok{(}\SpecialStringTok{f\textquotesingle{}OpenCV }\SpecialCharTok{\{}\NormalTok{cv2}\SpecialCharTok{.}\NormalTok{\_\_version\_\_}\SpecialCharTok{\}}\SpecialStringTok{ | NumPy }\SpecialCharTok{\{}\NormalTok{np}\SpecialCharTok{.}\NormalTok{\_\_version\_\_}\SpecialCharTok{\}}\SpecialStringTok{\textquotesingle{}}\NormalTok{)}
\BuiltInTok{print}\NormalTok{(}\StringTok{"Carregando imagem..."}\NormalTok{)}
\BuiltInTok{print}\NormalTok{(}\BuiltInTok{help}\NormalTok{(mm.show))}
\end{Highlighting}
\end{Shaded}

\begin{verbatim}
OpenCV 4.12.0 | NumPy 2.2.6
Carregando imagem...
Help on function show in module morph:

show(*args, title=None)
    This function will draw images f
    input: <*args> set of images f_i, where i>0 is binary image
          <title> optional string title for the plot
    output: image drawing
    Example:
    f1, f2 = np.zeros((100, 100,3)),  np.zeros((100, 100))
    f2[50:60, 50:60] = 1
    mm.show(f1, f2, title='Exemplo')

None
\end{verbatim}

\section{1.3 Lendo e Exibindo Imagens}\label{lendo-e-exibindo-imagens-2}

A operação mais básica em PDI é a leitura de uma imagem. A função
\texttt{mm.read()} da \texttt{morph.py} aceita caminhos locais e URLs.

\begin{Shaded}
\begin{Highlighting}[]
\ImportTok{import}\NormalTok{ urllib.request, os}

\CommentTok{\# Baixar imagem de exemplo}
\NormalTok{url }\OperatorTok{=} \StringTok{\textquotesingle{}https://raw.githubusercontent.com/opencv/opencv/master/samples/data/lena.jpg\textquotesingle{}}
\NormalTok{os.makedirs(}\StringTok{\textquotesingle{}imagens\textquotesingle{}}\NormalTok{, exist\_ok}\OperatorTok{=}\VariableTok{True}\NormalTok{)}
\ControlFlowTok{if} \KeywordTok{not}\NormalTok{ os.path.exists(}\StringTok{\textquotesingle{}imagens/lena.jpg\textquotesingle{}}\NormalTok{):}
\NormalTok{    urllib.request.urlretrieve(url, }\StringTok{\textquotesingle{}imagens/lena.jpg\textquotesingle{}}\NormalTok{)}

\CommentTok{\# Ler imagem com morph.py}
\NormalTok{img }\OperatorTok{=}\NormalTok{ mm.read(}\StringTok{\textquotesingle{}imagens/lena.jpg\textquotesingle{}}\NormalTok{)}

\BuiltInTok{print}\NormalTok{(}\SpecialStringTok{f\textquotesingle{}Shape  : }\SpecialCharTok{\{}\NormalTok{img}\SpecialCharTok{.}\NormalTok{shape}\SpecialCharTok{\}}\SpecialStringTok{\textquotesingle{}}\NormalTok{)    }\CommentTok{\# (altura, largura, canais)}
\BuiltInTok{print}\NormalTok{(}\SpecialStringTok{f\textquotesingle{}Dtype  : }\SpecialCharTok{\{}\NormalTok{img}\SpecialCharTok{.}\NormalTok{dtype}\SpecialCharTok{\}}\SpecialStringTok{\textquotesingle{}}\NormalTok{)    }\CommentTok{\# uint8}
\BuiltInTok{print}\NormalTok{(}\SpecialStringTok{f\textquotesingle{}Min/Max: }\SpecialCharTok{\{}\NormalTok{img}\SpecialCharTok{.}\BuiltInTok{min}\NormalTok{()}\SpecialCharTok{\}}\SpecialStringTok{ / }\SpecialCharTok{\{}\NormalTok{img}\SpecialCharTok{.}\BuiltInTok{max}\NormalTok{()}\SpecialCharTok{\}}\SpecialStringTok{\textquotesingle{}}\NormalTok{)}

\NormalTok{img\_rgb }\OperatorTok{=}\NormalTok{ cv2.cvtColor(img, cv2.COLOR\_BGR2RGB)}
\NormalTok{mm.show(img\_rgb, title}\OperatorTok{=}\StringTok{\textquotesingle{}Imagem original\textquotesingle{}}\NormalTok{)}
\end{Highlighting}
\end{Shaded}

\begin{verbatim}
Shape  : (512, 512, 3)
Dtype  : uint8
Min/Max: 0 / 255
\end{verbatim}

\begin{figure}[H]

\centering{

\pandocbounded{\includegraphics[keepaspectratio]{cap03/cap03.py.pt_files/figure-pdf/fig-1-lena-output-2.png}}

}

\caption{\label{fig-1-lena}MVS Linear --- Margem Máxima entre Íris
Setosa e Versicolor (atributos de pétala).}

\end{figure}%

\section{1.4 Tipos de Imagem}\label{tipos-de-imagem-1}

\subsection{1.4.1 Conversão para Tons de
Cinza}\label{conversuxe3o-para-tons-de-cinza-2}

A conversão de RGB para tons de cinza aplica pesos perceptuais definidos
pela ITU-R BT.601:

\begin{equation}\phantomsection\label{eq-1-cinza}{g = 0.299\,R + 0.587\,G + 0.114\,B}\end{equation}

\begin{Shaded}
\begin{Highlighting}[]
\CommentTok{\# Converter para tons de cinza}
\NormalTok{img\_gray }\OperatorTok{=}\NormalTok{ mm.gray(img)}
\BuiltInTok{print}\NormalTok{(}\SpecialStringTok{f\textquotesingle{}Shape cinza: }\SpecialCharTok{\{}\NormalTok{img\_gray}\SpecialCharTok{.}\NormalTok{shape}\SpecialCharTok{\}}\SpecialStringTok{\textquotesingle{}}\NormalTok{)}

\NormalTok{mm.show([img, img\_gray], title}\OperatorTok{=}\NormalTok{[}\StringTok{\textquotesingle{}Original (RGB)\textquotesingle{}}\NormalTok{, }\StringTok{\textquotesingle{}Tons de Cinza\textquotesingle{}}\NormalTok{])}
\end{Highlighting}
\end{Shaded}

\subsection{1.4.2 Limiarização
(Binarização)}\label{limiarizauxe7uxe3o-binarizauxe7uxe3o-1}

A limiarização converte uma imagem em tons de cinza para binária:

\begin{equation}\phantomsection\label{eq-1-limiar}{g(x,y) = \begin{cases} 255 & \text{se } f(x,y) \geq T \\ 0 & \text{caso contrário} \end{cases}}\end{equation}

O método de \textbf{Otsu} (Otsu, 1979) determina \(T\) automaticamente
maximizando a variância inter-classes do histograma.

\begin{Shaded}
\begin{Highlighting}[]
\CommentTok{\# Limiarização com T fixo}
\NormalTok{T }\OperatorTok{=} \DecValTok{128}
\NormalTok{img\_bin }\OperatorTok{=}\NormalTok{ mm.threshold(img\_gray, T)}

\CommentTok{\# Limiarização Otsu (T automático)}
\NormalTok{\_, img\_otsu }\OperatorTok{=}\NormalTok{ cv2.threshold(img\_gray, }\DecValTok{0}\NormalTok{, }\DecValTok{255}\NormalTok{,}
\NormalTok{                             cv2.THRESH\_BINARY }\OperatorTok{+}\NormalTok{ cv2.THRESH\_OTSU)}

\BuiltInTok{print}\NormalTok{(}\SpecialStringTok{f\textquotesingle{}Valores únicos (T=}\SpecialCharTok{\{}\NormalTok{T}\SpecialCharTok{\}}\SpecialStringTok{): }\SpecialCharTok{\{}\BuiltInTok{sorted}\NormalTok{(}\BuiltInTok{set}\NormalTok{(img\_bin.flatten()))}\SpecialCharTok{\}}\SpecialStringTok{\textquotesingle{}}\NormalTok{)}

\NormalTok{mm.show([img\_gray, img\_bin, img\_otsu],}
\NormalTok{        title}\OperatorTok{=}\NormalTok{[}\SpecialStringTok{f\textquotesingle{}Cinza\textquotesingle{}}\NormalTok{, }\SpecialStringTok{f\textquotesingle{}Binária (T=}\SpecialCharTok{\{}\NormalTok{T}\SpecialCharTok{\}}\SpecialStringTok{)\textquotesingle{}}\NormalTok{, }\StringTok{\textquotesingle{}Otsu\textquotesingle{}}\NormalTok{])}
\end{Highlighting}
\end{Shaded}

\section{1.5 Acesso a Pixels}\label{acesso-a-pixels-2}

Em Python/NumPy, o acesso a pixels usa indexação de arrays
\texttt{img{[}linha,\ coluna{]}}.

\begin{Shaded}
\begin{Highlighting}[]
\NormalTok{r, c }\OperatorTok{=} \DecValTok{100}\NormalTok{, }\DecValTok{100}

\CommentTok{\# Acessar pixel}
\NormalTok{pixel\_cinza }\OperatorTok{=}\NormalTok{ img\_gray[r, c]}
\NormalTok{pixel\_rgb   }\OperatorTok{=}\NormalTok{ img[r, c]       }\CommentTok{\# array [R, G, B]}

\BuiltInTok{print}\NormalTok{(}\SpecialStringTok{f\textquotesingle{}Pixel cinza (}\SpecialCharTok{\{}\NormalTok{r}\SpecialCharTok{\}}\SpecialStringTok{,}\SpecialCharTok{\{}\NormalTok{c}\SpecialCharTok{\}}\SpecialStringTok{): }\SpecialCharTok{\{}\NormalTok{pixel\_cinza}\SpecialCharTok{\}}\SpecialStringTok{\textquotesingle{}}\NormalTok{)}
\BuiltInTok{print}\NormalTok{(}\SpecialStringTok{f\textquotesingle{}Pixel RGB   (}\SpecialCharTok{\{}\NormalTok{r}\SpecialCharTok{\}}\SpecialStringTok{,}\SpecialCharTok{\{}\NormalTok{c}\SpecialCharTok{\}}\SpecialStringTok{): R=}\SpecialCharTok{\{}\NormalTok{pixel\_rgb[}\DecValTok{0}\NormalTok{]}\SpecialCharTok{\}}\SpecialStringTok{  G=}\SpecialCharTok{\{}\NormalTok{pixel\_rgb[}\DecValTok{1}\NormalTok{]}\SpecialCharTok{\}}\SpecialStringTok{  B=}\SpecialCharTok{\{}\NormalTok{pixel\_rgb[}\DecValTok{2}\NormalTok{]}\SpecialCharTok{\}}\SpecialStringTok{\textquotesingle{}}\NormalTok{)}

\CommentTok{\# Imagem sintética 5×5 com pixel central branco}
\NormalTok{syn }\OperatorTok{=}\NormalTok{ np.zeros((}\DecValTok{5}\NormalTok{, }\DecValTok{5}\NormalTok{), dtype}\OperatorTok{=}\NormalTok{np.uint8)}
\NormalTok{syn[}\DecValTok{2}\NormalTok{, }\DecValTok{2}\NormalTok{] }\OperatorTok{=} \DecValTok{255}
\BuiltInTok{print}\NormalTok{(}\StringTok{\textquotesingle{}}\CharTok{\textbackslash{}n}\StringTok{Matriz 5×5:\textquotesingle{}}\NormalTok{)}
\BuiltInTok{print}\NormalTok{(syn)}
\end{Highlighting}
\end{Shaded}

\section{1.6 Resumo}\label{resumo-2}

Neste capítulo foram apresentados os fundamentos de representação de
imagens digitais:

\begin{itemize}
\tightlist
\item
  Imagem digital = função \(f(x,y)\) mapeando coordenadas para
  intensidades
\item
  Tipos: binária (\(k=1\), \(\{0,255\}\)), tons de cinza (\(k=1\),
  \([0,255]\)), RGB (\(k=3\))
\item
  \texttt{morph.py}: biblioteca didática que simplifica operações de PDI
  em Python
\item
  Limiarização: operação fundamental de binarização
\end{itemize}

O Capítulo 2 estuda \textbf{histogramas e equalização}.

\subsection{Questões de Revisão}\label{questuxf5es-de-revisuxe3o-1}

\begin{enumerate}
\def\labelenumi{\arabic{enumi}.}
\tightlist
\item
  Qual a diferença entre imagem em tons de cinza e imagem binária?
\item
  Como o método de Otsu determina o limiar \(T\) automaticamente?
\item
  O que representa \texttt{img.shape} para uma imagem NumPy RGB?
\end{enumerate}

\cleardoublepage
\phantomsection
\addcontentsline{toc}{part}{Appendices}
\appendix

\chapter*{Referências}\label{referuxeancias-1}
\addcontentsline{toc}{chapter}{Referências}

\markboth{Referências}{Referências}

\phantomsection\label{refs}
\begin{CSLReferences}{0}{1}
\bibitem[\citeproctext]{ref-bradski2016}
BRADSKI, Gary; KAEHLER, Adrian. \textbf{Learning OpenCV 4: Computer
Vision with the OpenCV Library}. 2. ed. Sebastopol, CA: O'Reilly Media,
2016.

\bibitem[\citeproctext]{ref-gonzalez2018digital}
GONZALEZ, R. C.; WOODS, R. E. \textbf{Digital Image Processing}. 4th.
ed. New York: Pearson, 2018.

\bibitem[\citeproctext]{ref-itu_bt601}
ITU-R. \textbf{Recommendation BT.601: Studio encoding parameters of
digital television for standard 4:3 and wide screen 16:9 aspect ratios}.
\url{https://www.itu.int/rec/R-REC-BT.601/}, 2011.

\bibitem[\citeproctext]{ref-knuth_literate_1984}
KNUTH, Donald E. \href{https://doi.org/10.1093/comjnl/27.2.97}{Literate
Programming}. \textbf{The Computer Journal}, v. 27, n. 2, p. 97--111,
1984.

\bibitem[\citeproctext]{ref-knuth_art_1997}
KNUTH, Donald E. \textbf{The Art of Computer Programming}. 3rd. ed.
Reading, MA: Addison-Wesley, 1997.

\bibitem[\citeproctext]{ref-lewis2020retrieval}
LEWIS, Patrick \emph{et al.} Retrieval-augmented generation for
knowledge-intensive nlp tasks. \textbf{Advances in neural information
processing systems}, v. 33, p. 9459--9474, 2020.

\bibitem[\citeproctext]{ref-otsu1979}
OTSU, Nobuyuki. \href{https://doi.org/10.1109/TSMC.1979.4310076}{A
threshold selection method from gray-level histograms}. \textbf{IEEE
Transactions on Systems, Man, and Cybernetics}, v. 9, n. 1, p. 62--66,
1979.

\bibitem[\citeproctext]{ref-russ2023}
RUSS, John C.; RUSS, John C. (Jr. ). \textbf{The Image Processing
Handbook}. 8. ed. \emph{{[}S.l.{]}}: CRC Press, 2023.

\bibitem[\citeproctext]{ref-szeliski2022}
SZELISKI, Richard.
\textbf{\href{https://doi.org/10.1007/978-3-030-87137-2}{Computer
Vision: Algorithms and Applications}}. 2. ed. Cham, Switzerland:
Springer, 2022.

\bibitem[\citeproctext]{ref-zampirolli2024practical}
ZAMPIROLLI, Francisco de Assis \emph{et al.} A Practical Digital Image
Processing Course with morph.py. \emph{Em}: Porto Alegre, RS, Brasil:
Sociedade Brasileira de Computação (SBC), 2024. Disponível em:
\textless{}\url{https://sol.sbc.org.br/index.php/educomp/article/view/28199}\textgreater{}

\bibitem[\citeproctext]{ref-zampirolli2025morph}
ZAMPIROLLI, Francisco de Assis \emph{et al.}
\href{https://doi.org/10.5753/rbie.2025.5009}{Teaching Hands-On Digital
Image Processing with morph.py: Methods and Comprehensive Results}.
\textbf{Revista Brasileira de Informática na Educação}, v. 33, p.
990--1026, 2025.

\end{CSLReferences}

\end{tcolorbox}


\backmatter


\end{document}
